% ============================================================
% TFG — Anatomía Emocional de un Modelo Transformer:
% Compresión Selectiva e Interpretabilidad Mecánica de BERT
% Autor: Guido Biosca Lasa
% Versión MH: archivos renombrados, includes actualizados
% ============================================================
\documentclass[11pt, a4paper]{article}

% --- Tipografía ---
\usepackage[utf8]{inputenc}
\usepackage[T1]{fontenc}
\usepackage{newpxtext}          % Palatino para texto
\usepackage{newpxmath}          % Palatino para matemáticas
\usepackage[scaled=0.85]{beramono} % Bera Mono para monoespaciado
\usepackage{microtype}          % Mejora kerning y justificación

% --- Idioma ---
\usepackage[catalan, english, main=spanish]{babel}

% --- Geometría ---
\usepackage{geometry}
\geometry{
    margin=2.5cm,
    headheight=14pt
}

% --- Espaciado ---
\usepackage{setspace}
\setstretch{1.25}
\usepackage{parskip}
\setlength{\parskip}{4pt}

% --- Tablas y figuras ---
\usepackage{graphicx}
\usepackage{booktabs}
\usepackage{array}
\usepackage{longtable}
\usepackage{float}
\usepackage{caption}
\captionsetup{
    font=small,
    labelfont=bf,
    skip=8pt
}

% --- Listas ---
\usepackage{enumitem}

% --- Matemáticas y símbolos ---
\usepackage{amsmath}
\usepackage{amssymb}
\usepackage{eurosym}

% --- Color ---
\usepackage[table]{xcolor}
\definecolor{linkblue}{RGB}{10,60,120}

% --- Hiperenlaces ---
\usepackage{hyperref}
\hypersetup{
    colorlinks=true,
    linkcolor=linkblue,
    citecolor=linkblue,
    urlcolor=linkblue,
    pdfauthor={Guido Biosca Lasa},
    pdftitle={Anatomía Emocional de un Modelo Transformer},
}

% --- Encabezados y pies ---
\usepackage{fancyhdr}
\pagestyle{fancy}
\fancyhf{}
\fancyhead[L]{\small\itshape Anatomía Emocional de un Modelo Transformer}
\fancyhead[R]{\small\itshape G.\ Biosca Lasa}
\fancyfoot[C]{\thepage}
\renewcommand{\headrulewidth}{0.4pt}

\fancypagestyle{plain}{
    \fancyhf{}
    \fancyfoot[C]{\thepage}
    \renewcommand{\headrulewidth}{0pt}
}

% --- Formato de secciones (serif bold, sin sans-serif) ---
\usepackage{titlesec}

\titleformat{\section}{\Large\bfseries}{\thesection.}{0.5em}{}
\titleformat{\subsection}{\large\bfseries}{\thesubsection.}{0.5em}{}
\titleformat{\subsubsection}{\normalsize\bfseries}{\thesubsubsection.}{0.5em}{}

\titlespacing{\section}{0pt}{12pt}{6pt}
\titlespacing{\subsection}{0pt}{10pt}{4pt}
\titlespacing{\subsubsection}{0pt}{8pt}{3pt}

% --- Gantt (anexo de gestión) ---
\usepackage{pgfgantt}
\usepackage{pdflscape}
\definecolor{ganttG0}{RGB}{173,216,230}
\definecolor{ganttG1}{RGB}{144,238,144}
\definecolor{ganttG2}{RGB}{255,200,100}
\definecolor{ganttGE}{RGB}{255,140,100}
\definecolor{ganttGO}{RGB}{180,180,180}
\definecolor{ganttG9}{RGB}{200,170,210}

% --- Algoritmos ---
\usepackage{algorithm}
\usepackage{algpseudocode}
\renewcommand{\listalgorithmname}{Lista de algoritmos}

% --- Glosario / acrónimos ligero (tabla manual en anexo) ---

% --- TikZ extra para diagramas de cap2/cap3 ---
\usepackage{tikz}
\usetikzlibrary{arrows.meta, positioning, shapes.geometric, decorations.pathreplacing,
                fit, backgrounds, calc, shadings}

% --- Bibliografía ---
\usepackage[backend=biber, style=numeric, sorting=none]{biblatex}
\addbibresource{referencias.bib}

% ============================================================
\begin{document}

% --- Portada ---
\begin{titlepage}
    \centering
    \vspace*{1.5cm}

    % \includegraphics[height=2cm]{images/logo-fib.png}\\[0.5cm]

    {\Large\textbf{Universitat Politècnica de Catalunya}}\\[0.2cm]
    {\large Facultat d'Informàtica de Barcelona}\\[0.15cm]
    {\normalsize Grau en Enginyeria Informàtica --- Especialitat: Computació}

    \vspace{2cm}

    \rule{0.8\textwidth}{0.5pt}
    \vspace{0.6cm}

    {\LARGE\bfseries Anatomía Emocional de un\\[0.3em]
    Modelo Transformer}\\[0.6em]
    {\Large Compresión Selectiva e Interpretabilidad\\[0.2em]
    Mecánica de BERT}

    \vspace{0.6cm}
    \rule{0.8\textwidth}{0.5pt}

    \vspace{1cm}

    {\large\textbf{Trabajo Final de Grado}}

    \vfill

    \begin{tabular}{rl}
        \textbf{Autor:}     & Guido Biosca Lasa \\[0.3cm]
        \textbf{Director:}  & Lluís Padró Cirera \\[0.3cm]
        \textbf{Tutora GEP:} & Carolina Consolación Segura \\
    \end{tabular}

    \vspace{2cm}
    {\large Mayo 2026}
    \vspace{1cm}
\end{titlepage}

% --- Resúmenes (catalán, español, inglés) ---
\newpage
\thispagestyle{plain}

\begin{otherlanguage}{catalan}
\section*{Resum}
\addcontentsline{toc}{section}{Resum}
\noindent
On viu, dins d'un Transformer fine-tuneado, la capacitat de classificar
emocions? Aquest treball aborda la pregunta obrint BERT-base amb un
bisturí matemàtic ---la compressió per descomposició en valors singulars
(SVD)---: comprimir cadascun dels 72 components lineals del model i
observar què es trenca mesura la seva contribució funcional amb l'eina
més directa possible, la seva absència. Sobre GoEmotions (classificació
multietiqueta de 23 emocions) s'apliquen cinc tècniques
d'interpretabilitat mecànica ---\emph{probing} lineal, \emph{logit lens},
\emph{activation patching}, ablació de les 144 capçaleres d'atenció i
anàlisi de selectivitat de les 36\,864 neurones \emph{feed-forward}---
coordinades amb un estudi sistemàtic de sensibilitat a la compressió.

Els components no contribueixen per igual: comprimir Query a rang 64
reté el 98{,}6\% del rendiment, mentre que comprimir FFN Intermediate al
mateix rang el destrueix completament, una asimetria de factor 72$\times$
normalitzada per paràmetres eliminats. Les cinc tècniques convergeixen
en una arquitectura funcional de tres nivells ---senyal lèxica en capes
primerenques, còmput de transició en intermèdies i alineació geomètrica
amb la base del classificador en tardanes--- on la \emph{xarxa
feed-forward de la capa 11 actua com a coll d'ampolla suficient}:
restaurar-la sola en un model prèviament col·lapsat per compressió
recupera el 100\% del rendiment, un resultat que matisa la narrativa
habitual sobre el paper central de l'atenció en favor d'un repartiment
funcional on l'atenció encamina i la FFN computa.

Com a aplicació derivada del mapa funcional, un algorisme greedy
informat per dades empíriques de sensibilitat domina la frontera de
Pareto sobre 21 estratègies avaluades (8 dels 9 punts Pareto-òptims).
Un resultat negatiu el precedeix i el justifica metodològicament: l'intent
previ de traduir interpretabilitat qualitativa a heurístiques explícites
convergeix exactament sobre la mateixa frontera que les baselines cegues,
evidenciant que el coneixement mecanicista qualitatiu és necessari però
no suficient sense quantificació empírica. El treball es tanca amb cinc
prediccions falsables sobre la generalització del framework a altres
models i tasques.

\textbf{Paraules clau:} interpretabilitat mecànica de Transformers, BERT,
compressió SVD, classificació d'emocions, GoEmotions, \emph{probing},
\emph{activation patching}, frontera de Pareto.
\end{otherlanguage}

\vspace{1cm}

\section*{Resumen}
\addcontentsline{toc}{section}{Resumen}
\noindent
¿Dónde vive, dentro de un Transformer fine-tuneado, la capacidad de
clasificar emociones? Este trabajo aborda la pregunta abriendo BERT-base
con un bisturí matemático ---la compresión por descomposición en valores
singulares (SVD)---: comprimir cada uno de los 72 componentes lineales del
modelo y observar qué se rompe mide su contribución funcional con la
herramienta más directa posible, su ausencia. Sobre GoEmotions
(clasificación multi-etiqueta de 23 emociones) se aplican cinco técnicas de
interpretabilidad mecánica ---\emph{probing} lineal, \emph{logit lens},
\emph{activation patching}, ablación de las 144 cabezas de atención y
análisis de selectividad de las 36\,864 neuronas \emph{feed-forward}---
coordinadas con un estudio sistemático de sensibilidad a la compresión.

Los componentes no contribuyen por igual: comprimir Query a rango 64
retiene el 98{,}6\% del rendimiento, mientras que comprimir FFN
Intermediate al mismo rango lo destruye por completo, una asimetría de
factor 72$\times$ normalizada por parámetros eliminados. Las cinco
técnicas convergen en una arquitectura funcional de tres niveles ---señal
léxica en capas tempranas, cómputo de transición en intermedias y
alineación geométrica con la base del clasificador en tardías--- donde la
\emph{red feed-forward de la capa 11 actúa como cuello de botella
suficiente}: restaurarla sola en un modelo previamente colapsado por
compresión recupera el 100\% del rendimiento, un resultado que matiza la
narrativa habitual sobre el papel central de la atención en favor de un
reparto funcional donde la atención enruta y la FFN computa.

Como aplicación derivada del mapa funcional, un algoritmo greedy
informado por datos empíricos de sensibilidad domina la frontera de
Pareto sobre 21 estrategias evaluadas (8 de los 9 puntos
Pareto-óptimos). Un resultado negativo lo precede y lo justifica
metodológicamente: el intento previo de traducir interpretabilidad
cualitativa a heurísticas explícitas converge exactamente sobre la misma
frontera que las baselines ciegas, evidenciando que el conocimiento
mecanicista cualitativo es necesario pero no suficiente sin cuantificación
empírica. El trabajo cierra con cinco predicciones falsables sobre la
generalización del framework a otros modelos y tareas.

\textbf{Palabras clave:} interpretabilidad mecánica de Transformers, BERT,
compresión SVD, clasificación de emociones, GoEmotions, \emph{probing},
\emph{activation patching}, frontera de Pareto.

\vspace{1cm}

\begin{otherlanguage}{english}
\section*{Abstract}
\addcontentsline{toc}{section}{Abstract}
\noindent
Where, inside a fine-tuned Transformer, does the ability to classify
emotions live? This work tackles the question by opening BERT-base with a
mathematical scalpel ---selective singular value decomposition (SVD)
compression---: compressing each of the 72 linear components and observing
what breaks measures its functional contribution as directly as possible,
through its absence. On GoEmotions (multi-label classification of 23
emotions), five mechanistic interpretability techniques ---linear probing,
logit lens, activation patching, ablation of the 144 attention heads and
selectivity analysis of the 36{,}864 \emph{feed-forward} neurons--- are
combined with a systematic compression sensitivity study.

Components do not contribute equally: compressing Query to rank 64
retains 98.6\% of the performance, whereas compressing FFN Intermediate
at the same rank destroys the model entirely, a 72$\times$ asymmetry once
normalised by removed parameters. The five techniques converge on a
three-tier functional architecture ---lexical signal in early layers,
transition computation in middle layers, geometric alignment with the
classifier basis in late layers--- where the \emph{layer-11 feed-forward
network acts as a sufficient bottleneck}: restoring it alone in a model
previously collapsed by compression recovers 100\% of the original
performance, a result that nuances the conventional narrative about the
central role of attention in favour of a functional split where attention
routes and FFN computes.

As a practical application of the functional map obtained, a greedy
algorithm informed by empirical sensitivity data dominates the Pareto
frontier across 21 evaluated strategies (8 out of 9 Pareto-optimal
points). A negative result precedes and methodologically justifies it:
attempts to translate qualitative interpretability into explicit
heuristics converge exactly onto the same frontier as blind baselines,
showing that qualitative mechanistic knowledge is necessary but not
sufficient without empirical quantification. The work closes with five
falsifiable predictions on the generalisation of the framework to other
models and tasks.

\textbf{Keywords:} Transformer mechanistic interpretability, BERT, SVD
compression, emotion classification, GoEmotions, probing, activation
patching, Pareto frontier.
\end{otherlanguage}

% --- Índice y listas auxiliares ---
\newpage
\tableofcontents
\newpage
\listoffigures
\newpage
\listoftables
\newpage

% --- Glosario de acrónimos ---
\section*{Glosario de acrónimos}
\addcontentsline{toc}{section}{Glosario de acrónimos}
\begin{tabular}{@{}ll@{}}
    \textbf{BCE}    & Binary Cross-Entropy (entropía cruzada binaria) \\
    \textbf{BERT}   & Bidirectional Encoder Representations from Transformers \\
    \textbf{CLS}    & Classification token (token especial de clasificación) \\
    \textbf{CPI}    & Cost Performance Index (índice de rendimiento del coste) \\
    \textbf{CSV}    & Comma-Separated Values \\
    \textbf{DTU}    & Technical University of Denmark \\
    \textbf{F1}     & Métrica armónica de precisión y recall (macro/micro) \\
    \textbf{FFN}    & Feed-Forward Network \\
    \textbf{FIB}    & Facultat d'Informàtica de Barcelona \\
    \textbf{GeLU}   & Gaussian Error Linear Unit \\
    \textbf{GEP}    & Gestió de Projectes (asignatura de gestión del proyecto) \\
    \textbf{GPU}    & Graphics Processing Unit \\
    \textbf{KL}     & Kullback--Leibler (divergencia) \\
    \textbf{LoRA}   & Low-Rank Adaptation \\
    \textbf{MHA}    & Multi-Head Attention \\
    \textbf{MLM}    & Masked Language Modeling \\
    \textbf{NB}     & Notebook (Jupyter) \\
    \textbf{NLP}    & Natural Language Processing \\
    \textbf{NSP}    & Next Sentence Prediction \\
    \textbf{Q, K, V} & Query, Key, Value (proyecciones de atención) \\
    \textbf{SVD}    & Singular Value Decomposition (descomposición en valores singulares) \\
    \textbf{TFG}    & Trabajo Fin de Grado \\
    \textbf{t-SNE}  & t-distributed Stochastic Neighbor Embedding \\
    \textbf{UPC}    & Universitat Politècnica de Catalunya \\
\end{tabular}
\newpage

% --- Capítulos ---
% =============================================================
% CAPÍTULO 1 — Introducción
% Versión: matrícula de honor
% =============================================================

\section{Introducción}
\label{sec:introduccion}

\subsection{Contexto y motivación}
\label{sec:contexto-motivacion}

Los modelos de lenguaje basados en la arquitectura Transformer
\cite{vaswani2017attention} han transformado el procesamiento del lenguaje natural
en la última década. BERT \cite{devlin2019bert}, con sus 110~millones de parámetros,
demostró en 2018 que una sola red pre-entrenada podía alcanzar resultados de
referencia en tareas tan diversas como clasificación de texto, extracción de
información o respuesta a preguntas. Hoy, 110~millones de parámetros es la talla
\emph{pequeña}: GPT-3 alcanzó los 175\,000~millones \cite{brown2020language}, y los
modelos abiertos recientes operan en el rango de 8\,000 a 405\,000~millones. El coste
de desplegarlos en dispositivos móviles, sistemas embebidos o infraestructuras con
restricciones de latencia y energía sigue siendo un problema abierto, y proyecciones
recientes estiman que la inferencia ---no el entrenamiento--- concentra el grueso del
gasto energético a largo plazo de un modelo desplegado \cite{strubell2019energy,
schwartz2020green}.

Frente a este escenario, una pregunta directa orienta este trabajo: ¿se puede hacer
un modelo más pequeño sin destruir lo que ha aprendido? La compresión de modelos
lleva años intentando responderla con cuatro herramientas principales ---poda,
cuantización, destilación y factorización de bajo rango. Aquí se utiliza la última
en su variante más clásica: la descomposición en valores singulares (SVD), aplicada
de forma selectiva a las 72 matrices lineales del encoder de BERT-base. La SVD
ofrece tres ventajas para los objetivos de este trabajo: no requiere
re-entrenamiento, admite granularidad arbitraria (rango específico por matriz), y
proporciona garantías matemáticas de optimalidad respecto a la norma de Frobenius
mediante el teorema de Eckart-Young \cite{eckart1936approximation}.

Pero el proyecto no se limita a la compresión. Lo que comenzó como un análisis de
cuánto puede comprimirse cada componente derivó en algo más interesante: al observar
qué se rompe y qué sobrevive a la eliminación selectiva de partes del modelo, se
obtiene un mapa empírico de cómo está organizada internamente la información. La
compresión funciona, en este sentido, como un \emph{bisturí experimental}: cada
configuración comprimida es una intervención controlada que revela la importancia
funcional del componente afectado. Esta perspectiva conecta directamente la
compresión con la \emph{interpretabilidad mecánica}, una línea de investigación
reciente que busca abrir el modelo y caracterizar los circuitos internos
responsables de comportamientos específicos \cite{elhage2021mathematical,
meng2022locating, conmy2023towards}.

El trabajo combina así dos perspectivas habitualmente disjuntas en la literatura.
Por un lado, la compresión SVD selectiva por componente y profundidad. Por otro,
cinco técnicas de interpretabilidad mecánica ---probing lineal, \emph{logit lens},
\emph{activation patching}, ablación de cabezas de atención y análisis de selectividad
neuronal en las redes feed-forward--- que examinan la organización interna del
modelo desde ángulos diferentes pero convergentes. Y cierra el bucle traduciendo
esa comprensión en una estrategia de compresión \emph{informada por datos} que
domina la frontera de Pareto frente a las aproximaciones ciegas.

La elección de la clasificación de emociones como dominio de aplicación no es
arbitraria. GoEmotions \cite{demszky2020goemotions} ofrece 23~categorías
emocionales con propiedades semánticas heterogéneas: algunas emociones tienen
vocabulario distintivo (\emph{gratitude} $\leftrightarrow$ ``thanks''), mientras que
otras requieren comprensión pragmática profunda (\emph{disappointment},
\emph{realization}). Esta diversidad convierte la tarea en un test exigente para la
compresión, porque permite observar cómo la degradación del modelo afecta de forma
\emph{diferencial} a emociones con distintos perfiles lingüísticos. Un dataset
binario (sentimiento positivo/negativo) no haría visible esta vulnerabilidad
selectiva, que constituye uno de los hallazgos centrales del proyecto.

La motivación tiene tres vertientes. Desde el punto de vista \emph{práctico}, la
compresión SVD se aplica después del entrenamiento, no requiere hardware
especializado y es directamente escalable a cualquier modelo con capas lineales.
Desde el punto de vista \emph{científico}, la compresión selectiva proporciona
acceso a la organización interna del modelo de un modo que las técnicas
correlacionales (como el probing) no pueden replicar: cuando comprimir un componente
destruye una capacidad, esa relación es funcional, no solo asociativa. Y desde el
punto de vista de la \emph{sostenibilidad computacional}, reducir el tamaño de los
modelos para inferencia contribuye directamente a la filosofía \textit{Green AI}
\cite{schwartz2020green}.


\subsection{Formulación del problema}
\label{sec:formulacion-problema}

El problema central de este trabajo se articula en torno a tres preguntas
encadenadas. La primera describe; la segunda explica; la tercera prescribe.

\begin{quote}
    \textbf{P1 (descriptiva).} \textit{¿Cómo afecta la compresión SVD selectiva,
    aplicada por tipo de componente y por profundidad de capa, al rendimiento y a
    la calidad de las representaciones internas de un modelo Transformer
    \emph{fine-tuneado} para clasificación de emociones?}
\end{quote}

\begin{quote}
    \textbf{P2 (explicativa).} \textit{¿Qué revelan las técnicas de interpretabilidad
    mecánica sobre la organización interna de la información emocional en el
    modelo, y qué propiedades estructurales explican la sensibilidad diferencial
    observada en P1?}
\end{quote}

\begin{quote}
    \textbf{P3 (prescriptiva).} \textit{¿Puede el conocimiento mecanicista
    obtenido en P2 traducirse en estrategias de compresión que dominen a las
    aproximaciones ciegas establecidas en la literatura?}
\end{quote}

Los trabajos previos de compresión SVD en Transformers \cite{noach2020compressing,
hsu2022language} se han centrado mayoritariamente en compresión uniforme: el mismo
grado de reducción para todas las capas y componentes. Esta uniformidad ignora dos
asimetrías documentadas. Primero, las distintas capas de BERT codifican tipos de
información diferentes ---las tempranas capturan rasgos léxicos y sintácticos, las
tardías capturan rasgos semánticos y específicos de la tarea \cite{tenney2019bert,
jawahar2019what}. Segundo, los componentes de atención y las redes feed-forward
desempeñan roles funcionales distintos \cite{geva2021transformer, voita2019analyzing}.
La hipótesis subyacente del proyecto es que tratarlos igual es subóptimo, y que
medir cuánto subóptimo es ---y por qué--- requiere combinar compresión sistemática
con interpretabilidad mecánica.


\subsection{Objetivos}
\label{sec:objetivos}

El objetivo general del proyecto es:

\begin{quote}
    Producir un mapa funcional de cómo BERT-base \emph{fine-tuneado}
    organiza internamente la información emocional, utilizando la
    compresión SVD selectiva como instrumento experimental
    ---\emph{bisturí} que mide la contribución funcional de cada
    componente por la magnitud de lo que se rompe al eliminarlo--- y
    triangulando los hallazgos con cinco técnicas de interpretabilidad
    mecánica. Como aplicación derivada del mapa, se diseña una
    estrategia de compresión informada por datos empíricos de
    sensibilidad y se compara con las aproximaciones ciegas
    establecidas en la literatura.
\end{quote}

Este objetivo se descompone en cinco objetivos específicos que estructuran la
contribución y delimitan el alcance:

\begin{enumerate}[nosep, label=\textbf{O\arabic*.}]
    \item \textbf{Caracterizar la sensibilidad a la compresión SVD} por tipo de
    componente (Q, K, V, Output, FFN Intermediate, FFN Output), por profundidad
    de capa y por emoción, estableciendo un mapa cuantitativo del trade-off
    compresión--rendimiento bajo compresión uniforme y adaptativa.

    \item \textbf{Localizar y mapear la organización interna de la información
    emocional} mediante cinco técnicas de interpretabilidad mecánica
    complementarias: probing lineal para identificar la profundidad de
    cristalización de cada emoción, \emph{logit lens} para distinguir información
    presente de información alineada con la base del clasificador, activation
    patching para establecer dependencias funcionales, ablación de cabezas de
    atención para mapear la especialización atencional y análisis de selectividad
    neuronal para caracterizar la codificación en las redes feed-forward.

    \item \textbf{Establecer conexiones cuantitativas entre estructura interna y
    vulnerabilidad a la compresión}, identificando qué propiedades del modelo
    (profundidad de cristalización, intensidad de codificación neuronal,
    distribución espectral) predicen la caída de rendimiento bajo SVD.

    \item \textbf{Diseñar y evaluar una estrategia de compresión informada} que
    traduzca los hallazgos de interpretabilidad en un algoritmo de asignación de
    rangos, y demostrar su superioridad sobre las aproximaciones ciegas en la
    frontera de Pareto.

    \item \textbf{Evaluar la recuperación por fine-tuning} del modelo
    comprimido, investigando si la compresión SVD seguida de un breve
    re-entrenamiento puede actuar como regularización implícita que beneficie
    a las emociones infrarrepresentadas.
\end{enumerate}


\subsection{Contribuciones}
\label{sec:contribuciones-intro}

Las contribuciones principales de este trabajo son las siguientes:

\begin{enumerate}[label=\textbf{C\arabic*.}]

    \item \textbf{Mapa de sensibilidad a la compresión} desagregado por tipo de
    componente, profundidad de capa y emoción (O1). El mapa documenta una
    asimetría radical entre componentes: las proyecciones Query y Key toleran
    compresión a rango~128 con retención del 99\% del F1 macro, mientras que la
    FFN Intermediate al mismo rango colapsa al 7\%. La diferencia de
    sensibilidad alcanza un factor 14$\times$ en retención absoluta de F1 a
    rango~128, que se amplifica hasta un factor 72$\times$ a rango~64
    cuando se normaliza por el porcentaje de parámetros eliminados (Query
    retiene 98{,}6\%, FFN Intermediate colapsa al 0\%). Hasta donde sabemos,
    esta asimetría no había sido documentada cuantitativamente en la
    literatura sobre SVD en Transformers.

    \item \textbf{Triangulación de interpretabilidad mecánica} mediante cinco
    técnicas independientes (probing, \emph{logit lens}, \emph{activation patching}, ablación
    de cabezas y análisis neuronal) (O2). Los resultados convergen en una
    arquitectura funcional de tres niveles ---señal léxica en capas tempranas,
    cómputo de transición en capas medias, alineación geométrica con la base
    del clasificador en capas tardías--- que explica de forma unificada
    fenómenos previamente disjuntos: el patrón en U del \emph{logit lens}, la
    suficiencia causal de la capa~11 en \emph{activation patching}, el gradiente de
    criticidad atencional documentado por ablación, y la concentración del 84\%
    de las neuronas significativas en las capas 8--11.

    \item \textbf{Resolución de la paradoja información-dependencia.} El
    probing demuestra que las capas tardías apenas aportan información nueva
    (ganancia $<0{,}004$ en L10--L11), pero el \emph{activation patching} las
    identifica como las únicas suficientes para restaurar el modelo desde un
    colapso total. Mediante el \emph{logit lens} documentamos que esta aparente
    contradicción refleja una distinción funcional precisa: las capas tardías
    no \emph{crean} señal emocional ---la \emph{rotan} hacia la base sobre la
    que opera la cabeza de clasificación. La distinción entre información
    \emph{presente} (probing), información \emph{leíble} (\emph{logit lens}) y
    componentes \emph{suficientes} (patching) constituye una clarificación
    metodológica que va más allá del caso estudiado.

    \item \textbf{Conexión cuantitativa entre interpretabilidad y compresión}
    (O3). La profundidad de cristalización por probing predice la
    vulnerabilidad a la compresión SVD a rango 128
    (Spearman~$\rho \approx 0{,}48$), y la norma del vector de selectividad
    neuronal es el mejor predictor individual de la caída de F1 bajo
    compresión ($\rho = 0{,}64$, $p = 0{,}001$, $n=23$). La
    geometría que la SVD optimiza es \emph{ortogonal} a la organización
    funcional del modelo (correlación entre fracción de señal eliminada y
    caída de F1: $\rho \approx 0$): la SVD no daña selectivamente las neuronas
    emocionales, sino que las daña uniformemente, y la vulnerabilidad
    diferencial proviene de la intensidad de codificación, no de la geometría
    SVD.

    \item \textbf{Algoritmo greedy de compresión informada por datos} (O4),
    alimentado con datos empíricos de sensibilidad obtenidos en C1, que domina
    la frontera de Pareto con 8 de los 9 puntos Pareto-óptimos del estudio. A
    ratio de parámetros del 80\%, la estrategia greedy retiene el 87\% del F1
    macro, mientras que la compresión uniforme al mismo ratio retiene solo el
    43\%. El algoritmo descubre de forma independiente la jerarquía de
    importancia identificada por interpretabilidad ---comprime primero Q y K,
    nunca toca la FFN Intermediate de las capas tardías---, lo que constituye
    una validación cruzada entre dos líneas metodológicamente independientes.

    \item \textbf{Resultado negativo informativo: la heurística cualitativa
    no extrae valor adicional.} Tres estrategias informadas basadas en reglas
    heurísticas derivadas de los hallazgos de interpretabilidad
    \emph{convergen exactamente} sobre la frontera de las baselines ciegas
    (\texttt{informed\_aggressive} reproduce \texttt{uniform\_r256},
    \texttt{informed\_moderate} reproduce \texttt{uniform\_r512}), o requieren
    más parámetros que el modelo original sin alcanzar las retenciones más
    altas (\texttt{informed\_light} a ratio 1{,}285). Este resultado delimita
    el valor aplicado de la interpretabilidad mecánica para la compresión: el
    conocimiento cualitativo identifica \emph{qué} medir, pero solo los datos
    cuantitativos de sensibilidad permiten determinar \emph{cuánto}
    comprimir. La narrativa completa (heurística no innova $\rightarrow$
    pivot data-driven $\rightarrow$ dominancia Pareto) es, en sí misma, una
    observación metodológica transferible.

    \item \textbf{Evidencia preliminar de regularización por compresión}
    (O5). Un modelo comprimido al 86\% de parámetros, tras tres épocas de
    \emph{fine-tuning}, alcanza un F1 macro de 0{,}591 frente a 0{,}577 del baseline
    sin comprimir. Las ganancias se concentran en emociones
    infrarrepresentadas: \emph{embarrassment} pasa de F1 = 0{,}267 a F1 =
    0{,}509 ($+90\%$ relativo). Reportamos este hallazgo como observación
    consistente con la hipótesis de que la SVD, al eliminar direcciones de
    baja energía, actúa como regularización estructural que reduce el
    sobreajuste a patrones espurios. El alcance experimental (un modelo, una
    tarea, sin grupo de control con epochs adicionales) impide afirmar
    causalidad: en el Capítulo~\ref{sec:conclusiones} se discuten los
    experimentos de control que verificarían la hipótesis.

\end{enumerate}


\subsection{Marco del proyecto en la FIB}
\label{sec:marco-fib}

Este proyecto se enmarca en el Grado en Ingeniería Informática de la Facultat
d'Informàtica de Barcelona (FIB-UPC), dentro de la especialidad de Computación.
El trabajo integra conocimientos de múltiples asignaturas del grado. La
descomposición SVD y el análisis espectral se fundamentan en álgebra lineal y
cálculo matricial (M1, M2), mientras que la estabilidad numérica de las
aproximaciones de bajo rango conecta directamente con los contenidos de
Computación Numérica (CN). El diseño y la evaluación del algoritmo greedy de
compresión aplican técnicas de diseño algorítmico y análisis de complejidad (A,
EDA). Los fundamentos de aprendizaje automático (IA) sustentan tanto el
\emph{fine-tuning} de BERT como los clasificadores de probing, y la ejecución
eficiente en GPU (PAR) ha sido necesaria para el entrenamiento y la evaluación
sobre el dataset completo. Las metodologías de desarrollo de software (PROP)
se reflejan en la organización modular del código experimental. Un intercambio
en la Technical University of Denmark (DTU, Q1 2024--25) proporcionó formación
complementaria de nivel máster en \textit{Machine Learning}, \textit{Deep
Learning} y \textit{Graph Neural Networks}, incluyendo contacto con técnicas de
interpretabilidad mecánica que han sido instrumentales en el diseño del
Capítulo~\ref{sec:resultados-interpretabilidad}.

El proyecto tiene carácter eminentemente experimental y de investigación: no
se desarrolla un producto de software, sino que se lleva a cabo un estudio
sistemático sobre cómo la compresión SVD afecta a los distintos componentes de
un Transformer y qué revela esto sobre la organización interna de las
representaciones emocionales. La metodología es comparable a la de la
investigación empírica en NLP: hipótesis explícitas, experimentos controlados,
métricas estándar (F1 macro y micro, retención, ratio de parámetros), análisis
estadístico (correlaciones de Spearman con reporte de p-valores) y reporte
honesto de resultados ---incluidos los nulos.


\subsection{Alcance y limitaciones explícitas}
\label{sec:alcance-limitaciones}

Por transparencia metodológica, los resultados del trabajo deben interpretarse
dentro de las siguientes restricciones, que se discuten en mayor detalle en el
Capítulo~\ref{sec:conclusiones}:

\begin{itemize}[nosep]
    \item \textbf{Modelo único}: los experimentos se realizan sobre BERT-base
    (110\,M parámetros). La generalización a BERT-large, RoBERTa, o
    arquitecturas decoder-only (GPT-2, LLaMA) no se verifica empíricamente
    en este trabajo.

    \item \textbf{Tarea única}: clasificación multi-etiqueta de emociones
    sobre GoEmotions. La extensión a NER, inferencia textual o tareas
    generativas queda como trabajo futuro.

    \item \textbf{Compresión post-hoc}: la SVD se aplica una vez fine-tuneado
    el modelo, sin compresión durante el entrenamiento ni cuantización
    posterior.

    \item \textbf{Potencia estadística limitada}: las correlaciones se
    calculan sobre $n = 23$ emociones, lo que permite detectar efectos
    grandes con confianza pero no efectos moderados.
\end{itemize}

Estas limitaciones se asumen por diseño: el alcance del trabajo prioriza la
profundidad analítica sobre un único modelo y tarea frente a estudios más
superficiales sobre múltiples configuraciones.


\subsection{Estructura de la memoria}
\label{sec:estructura-memoria}

La memoria se organiza en siete capítulos. El Capítulo~\ref{sec:estado-arte}
revisa el estado del arte en compresión de Transformers e interpretabilidad
mecánica, sitúa los fundamentos matemáticos de la SVD y posiciona el trabajo
respecto a la literatura previa. El Capítulo~\ref{sec:metodologia-tfg} describe
la metodología experimental: dataset, modelo, técnicas de compresión y
herramientas de interpretabilidad.

Los tres capítulos siguientes contienen los resultados, en orden de granularidad
creciente. El Capítulo~\ref{sec:resultados-compresion} presenta los resultados
de compresión SVD: análisis espectral, transición de fase bajo compresión
uniforme, sensibilidad por componente y profundidad, y compresión adaptativa
energética. El Capítulo~\ref{sec:resultados-interpretabilidad} expone los
resultados de las cinco técnicas de interpretabilidad mecánica y argumenta su
convergencia en una arquitectura funcional unificada. El
Capítulo~\ref{sec:sintesis} cierra el bucle: utiliza los hallazgos de los
capítulos anteriores para diseñar y evaluar la estrategia greedy de compresión
informada, documenta el resultado negativo de la heurística previa y analiza la
recuperación mediante \emph{fine-tuning}.

El Capítulo~\ref{sec:conclusiones} sintetiza las contribuciones, discute las
limitaciones y propone líneas de investigación futura. La gestión del proyecto
(planificación temporal, presupuesto, sostenibilidad) se incluye como anexo
final.
% =============================================================
% CAPÍTULO 2 — Estado del arte y marco teórico
% =============================================================

\section{Estado del arte y marco teórico}
\label{sec:estado-arte}

Este capítulo proporciona los fundamentos teóricos y el contexto de investigación
necesarios para comprender el resto de la memoria. La exposición sigue un orden
constructivo: primero se traza la evolución del NLP hasta la arquitectura Transformer
(\S\ref{sec:sota-historia}); después se presenta el Transformer en detalle
(\S\ref{sec:sota-transformer}); a continuación se describe BERT
(\S\ref{sec:sota-bert}); se desarrollan los fundamentos matemáticos de la SVD
(\S\ref{sec:sota-svd-teoria}); se sitúa el problema de la clasificación de emociones
(\S\ref{sec:sota-emociones}); se revisan las técnicas de compresión
(\S\ref{sec:sota-compresion}); se presentan las herramientas de interpretabilidad
mecánica (\S\ref{sec:sota-interpretabilidad}); y se cierra con el posicionamiento
del trabajo (\S\ref{sec:sota-posicionamiento}).

% =============================================================
\subsection{Del procesamiento secuencial al Transformer}
\label{sec:sota-historia}
% =============================================================

El procesamiento del lenguaje natural fue durante décadas un problema fundamentalmente
secuencial. Las redes neuronales recurrentes (RNN) \cite{elman1990finding} procesaban
los tokens uno a uno, manteniendo un vector de estado oculto que acumulaba información
de los tokens anteriores. En la práctica, esta información se degradaba
exponencialmente con la distancia: las RNN sufrían el problema del
\emph{desvanecimiento del gradiente} \cite{bengio1994learning}, que impedía aprender
dependencias a larga distancia. Las LSTM \cite{hochreiter1997long} mitigaron
parcialmente este problema introduciendo mecanismos de puertas que controlan qué
información se almacena, se olvida o se transmite, y dominaron el NLP durante la
década de 2014--2017.

Sin embargo, las arquitecturas recurrentes tienen una limitación estructural: al
procesar los tokens secuencialmente, la computación no es paralelizable (el cálculo
de $h_t$ depende de $h_{t-1}$), lo que limita la velocidad de entrenamiento e impide
aprovechar eficientemente las GPUs.

Un primer avance fue el mecanismo de atención de Bahdanau et al.\
\cite{bahdanau2015neural}, que permitió al decoder de un modelo de traducción atender
selectivamente a distintas posiciones del encoder, eliminando el cuello de botella de
comprimir toda la frase de origen en un único vector. Pero la recurrencia seguía
presente.

La contribución de Vaswani et al.\ \cite{vaswani2017attention} fue radical: el
Transformer elimina completamente las operaciones recurrentes y se basa
exclusivamente en mecanismos de atención, logrando tres ventajas:

\begin{enumerate}[nosep]
    \item \textbf{Paralelismo total}: todas las posiciones se procesan simultáneamente.
    \item \textbf{Distancia constante}: la atención conecta cualquier par de posiciones
    en $O(1)$, frente a los $O(n)$ pasos de una RNN.
    \item \textbf{Escalabilidad}: permitió entrenar modelos significativamente más
    grandes, abriendo la puerta a los modelos de lenguaje pre-entrenados.
\end{enumerate}

El impacto fue inmediato en traducción automática, pero la revolución real vino
después: la arquitectura Transformer, pre-entrenada sobre grandes corpus, producía
representaciones lingüísticas de propósito general transferibles a prácticamente
cualquier tarea de NLP. Esto dio lugar a la era de los \textit{foundation models}
(BERT, GPT, T5).

% =============================================================
\subsection{La arquitectura Transformer}
\label{sec:sota-transformer}
% =============================================================

La arquitectura Transformer original \cite{vaswani2017attention} consta de dos
bloques apilados (encoder y decoder) conectados mediante atención cruzada
(Figura~\ref{fig:transformer-original}). El encoder procesa la secuencia de
entrada y produce representaciones contextualizadas; el decoder genera la secuencia
de salida de forma autoregresiva. BERT utiliza exclusivamente la torre del encoder,
descartando el decoder por completo, ya que su objetivo no es generar secuencias
sino producir representaciones bidireccionales reutilizables.

\begin{figure}[H]
    \centering
    \includegraphics[width=0.45\textwidth]{figures/transformer_original.png}
    \caption{Arquitectura Transformer original (Vaswani et al., 2017). El encoder
    (izquierda) y el decoder (derecha) se componen de bloques apilados $N$ veces.
    BERT utiliza exclusivamente la torre del encoder. Figura adaptada
    de~\cite{vaswani2017attention}.}
    \label{fig:transformer-original}
\end{figure}

La Figura~\ref{fig:bert-encoder} detalla la estructura interna de un bloque
encoder, que se apila 12 veces en BERT-base. El resto de esta sección describe
sus componentes.

\begin{figure}[H]
    \centering
    \shorthandoff{><}
    \begin{tikzpicture}[
        block/.style={draw, rounded corners=3pt, minimum height=0.85cm,
                      font=\small, minimum width=5.6cm, anchor=center},
        attn/.style={block, fill=blue!8},
        ffn/.style={block, fill=orange!8},
        norm/.style={block, fill=gray!6},
        emb/.style={block, fill=green!6},
        add/.style={draw, circle, minimum size=0.45cm, fill=yellow!12,
                    font=\scriptsize, inner sep=0pt},
        arr/.style={->, >=stealth, thick},
        res/.style={->, >=stealth, thick, dashed, gray!55},
        brace/.style={decorate, decoration={brace, amplitude=6pt}, thick, gray!70},
        lab/.style={font=\scriptsize, gray!60!black},
    ]
        % === Nodes ===
        \node[emb]  (emb)  at (0, 0.0) {Token + Position + Segment Embeddings};
        \node[norm] (ln1)  at (0, 1.4) {Layer Norm};
        \node[attn] (attn) at (0, 2.8) {Multi-Head Attention ($H{=}12$, $d_k{=}64$)};
        \node[add]  (add1) at (0, 3.9) {$+$};
        \node[norm] (ln2)  at (0, 5.0) {Layer Norm};
        \node[ffn]  (ffn)  at (0, 6.4) {Feed-Forward: $768 \rightarrow 3072 \rightarrow 768$};
        \node[add]  (add2) at (0, 7.5) {$+$};
        \node[font=\small] (out) at (0, 8.6) {Salida (al siguiente bloque o al clasificador)};

        % === Main vertical flow ===
        \draw[arr] (emb)  -- (ln1);
        \draw[arr] (ln1)  -- (attn);
        \draw[arr] (attn) -- (add1);
        \draw[arr] (add1) -- (ln2);
        \draw[arr] (ln2)  -- (ffn);
        \draw[arr] (ffn)  -- (add2);
        \draw[arr] (add2) -- (out);

        % === Residual 1: emb bypasses attn, joins at add1 ===
        \draw[res] (emb.east) -- (3.8, 0.0) -- (3.8, 3.9) -- (add1.east);

        % === Residual 2: add1 output bypasses ffn, joins at add2 ===
        \draw[res] (add1.east) -- (3.8, 3.9) -- (4.6, 3.9) -- (4.6, 7.5) -- (add2.east);

        % === Weight matrix annotations (left side) ===
        \node[lab] at (-4.2, 2.8) {$W_Q, W_K, W_V, W_O$};
        \node[lab] at (-4.2, 6.4) {$W_{\text{int}}, W_{\text{out}}$};

        % === Brace ===
        \draw[brace] (-5.6, 1.1) -- (-5.6, 8.0)
            node[midway, left=8pt, font=\small, gray!70!black, rotate=90]
            {$\times\, 12$ capas};
    \end{tikzpicture}
    \shorthandon{><}
    \caption{Arquitectura de un bloque encoder de BERT-base. Cada bloque contiene
    un sub-bloque de atención multi-cabeza (con 12 cabezas de dimensión 64) y un
    sub-bloque feed-forward (con expansión $768 \rightarrow 3072 \rightarrow 768$),
    ambos con conexión residual (líneas discontinuas) y layer normalization. Las 6
    matrices de pesos anotadas ($W_Q, W_K, W_V, W_O, W_{\text{int}}, W_{\text{out}}$)
    son el objetivo de la compresión SVD. BERT-base apila 12 de estos bloques.}
    \label{fig:bert-encoder}
\end{figure}

% -------------------------------------------------------------
\subsubsection{Atención multi-cabeza}
\label{sec:sota-mha}
% -------------------------------------------------------------

El mecanismo central del Transformer es la \emph{atención escalada por producto
punto}. Dada una secuencia de $n$ tokens representada como $X \in \mathbb{R}^{n
\times d}$, cada cabeza de atención $h$ (de un total de $H$) computa tres
proyecciones lineales:
\begin{equation}
    Q_h = X W_Q^h, \quad K_h = X W_K^h, \quad V_h = X W_V^h
    \label{eq:qkv}
\end{equation}
donde $W_Q^h, W_K^h, W_V^h \in \mathbb{R}^{d \times d_k}$ con $d_k = d/H$. La
\emph{Query} codifica qué busca cada token, la \emph{Key} codifica qué ofrece, y
el \emph{Value} contiene la información que se transmitirá.

La atención se computa en tres pasos. Primero, los scores de compatibilidad:
\begin{equation}
    S = \frac{Q_h K_h^T}{\sqrt{d_k}} \in \mathbb{R}^{n \times n}
    \label{eq:scores}
\end{equation}
El factor $1/\sqrt{d_k}$ evita que los productos punto crezcan con la dimensión,
estabilizando los gradientes a través de la softmax. Segundo, la normalización:
\begin{equation}
    A = \text{softmax}(S), \quad A_{ij} = \frac{e^{S_{ij}}}{\sum_{k=1}^{n} e^{S_{ik}}}
    \label{eq:softmax-attn}
\end{equation}
Cada fila $A_i$ es una distribución de probabilidad que determina cuánto atiende
el token $i$ a cada posición. Tercero, la salida como promedio ponderado de los
valores:
\begin{equation}
    \text{head}_h = A \cdot V_h \in \mathbb{R}^{n \times d_k}
    \label{eq:head-output}
\end{equation}

Las salidas de las $H$ cabezas se concatenan y se proyectan:
\begin{equation}
    \text{MultiHead}(X) = [\text{head}_1; \ldots; \text{head}_H]\, W_O
    \label{eq:multihead}
\end{equation}
donde $W_O \in \mathbb{R}^{d \times d}$. Cada cabeza opera en un subespacio de
dimensión $d_k$ y puede especializarse en un tipo diferente de relación entre
tokens. Investigaciones posteriores \cite{voita2019analyzing, clark2019what}
confirmaron que esta especialización emerge espontáneamente: algunas cabezas
atienden a posiciones sintácticas, otras a tokens semánticamente relacionados, y
otras a patrones posicionales fijos.

\begin{figure}[H]
    \centering
    \shorthandoff{><}
    \begin{tikzpicture}[
        block/.style={draw, rounded corners=2pt, minimum height=0.7cm, minimum width=1.4cm,
                      fill=blue!8, font=\small, anchor=center},
        op/.style={draw, circle, minimum size=0.55cm, fill=orange!12, font=\scriptsize,
                   inner sep=0pt},
        arr/.style={->, >=stealth, thick},
    ]
        % === Row 1: Input X (centered below projections) ===
        \node (X) at (1.8, 0) {$X$};

        % === Row 2: Linear projections ===
        \node[block] (WQ) at (0,   1.5) {$W_Q^h$};
        \node[block] (WK) at (1.8, 1.5) {$W_K^h$};
        \node[block] (WV) at (3.6, 1.5) {$W_V^h$};

        % === Row 3: Q, K, V labels ===
        \node[font=\small] (Q) at (0,   2.6) {$Q_h$};
        \node[font=\small] (K) at (1.8, 2.6) {$K_h$};
        \node[font=\small] (V) at (3.6, 2.6) {$V_h$};

        % === Row 4: MatMul Q·K^T ===
        \node[op] (mm1) at (0.9, 3.7) {$\times$};
        \node[font=\scriptsize, gray, right=0.7cm of mm1] {$QK^T\!/\!\sqrt{d_k}$};

        % === Row 5: Softmax ===
        \node[block, fill=green!8] (sm) at (0.9, 4.8) {softmax};

        % === Row 6: MatMul with V ===
        \node[op] (mm2) at (0.9, 5.9) {$\times$};

        % === Row 7: Output ===
        \node[font=\small] (out) at (0.9, 6.8) {$\text{head}_h$};

        % === Arrows: X to projections ===
        \draw[arr] (X) -- (WQ);
        \draw[arr] (X) -- (WK);
        \draw[arr] (X) -- (WV);

        % === Arrows: projections to labels ===
        \draw[arr] (WQ) -- (Q);
        \draw[arr] (WK) -- (K);
        \draw[arr] (WV) -- (V);

        % === Arrows: Q and K into mm1 ===
        \draw[arr] (Q) -- (0, 3.7) -- (mm1);
        \draw[arr] (K) -- (1.8, 3.7) -- (mm1);

        % === Arrow: mm1 to softmax to mm2 to output ===
        \draw[arr] (mm1) -- (sm);
        \draw[arr] (sm) -- (mm2);
        \draw[arr] (mm2) -- (out);

        % === Arrow: V bypasses to mm2 ===
        \draw[arr] (V) -- (3.6, 5.9) -- (mm2);
    \end{tikzpicture}
    \shorthandon{><}
    \caption{Mecanismo de atención para una cabeza $h$. La entrada $X$ se proyecta
    en Query, Key y Value. Los scores $QK^T/\sqrt{d_k}$ se normalizan con softmax
    y ponderan los valores para generar la salida.}
    \label{fig:attention-mechanism}
\end{figure}

% -------------------------------------------------------------
\subsubsection{El bloque feed-forward}
\label{sec:sota-ffn}
% -------------------------------------------------------------

Cada capa incluye un bloque feed-forward (FFN) que se aplica posición por posición:
\begin{equation}
    \text{FFN}(x) = \text{GELU}(x\, W_{\text{int}} + b_1)\, W_{\text{out}} + b_2
    \label{eq:ffn}
\end{equation}
donde $W_{\text{int}} \in \mathbb{R}^{d \times d_{ff}}$ y $W_{\text{out}} \in
\mathbb{R}^{d_{ff} \times d}$, con $d_{ff} = 4d = 3072$ en BERT-base. La función
de activación GELU (\textit{Gaussian Error Linear Unit}) \cite{hendrycks2016gaussian}
se define como:
\begin{equation}
    \text{GELU}(x) = x \cdot \Phi(x) \approx 0{,}5\, x \left(1 + \tanh\!\left[\sqrt{2/\pi}\,(x + 0{,}044715\, x^3)\right]\right)
    \label{eq:gelu}
\end{equation}
donde $\Phi(x)$ es la función de distribución acumulada de la normal estándar. A
diferencia de ReLU, GELU es suave y permite gradientes no-nulos para valores negativos
pequeños. Esta no-linealidad entre $W_{\text{int}}$ y $W_{\text{out}}$ es lo que
convierte a la FFN en un bloque computacional genuino: sin ella, las dos
multiplicaciones matriciales colapsarían en una sola transformación lineal.

La dimensión
intermedia expandida crea un espacio de representación cuatro veces mayor que el
espacio residual, donde el modelo puede computar transformaciones complejas antes de
proyectar de vuelta. Geva et al.\ \cite{geva2021transformer} propusieron que las
neuronas de esta capa funcionan como \emph{memorias clave-valor}: cada neurona
almacena un patrón de activación (clave, en $W_{\text{int}}$) y un vector de
actualización del espacio residual (valor, en $W_{\text{out}}$). Esta perspectiva
es directamente relevante para el análisis de especialización neuronal emocional
del Capítulo~\ref{sec:resultados-interpretabilidad}.

% -------------------------------------------------------------
\subsubsection{Conexiones residuales y flujo residual}
\label{sec:sota-residual}
% -------------------------------------------------------------

El output de cada sub-bloque se combina con su input mediante una \emph{conexión
residual} \cite{he2016deep} seguida de \emph{layer normalization} \cite{ba2016layer}:
\begin{align}
    h'  &= \text{LayerNorm}(x + \text{MultiHead}(x)) \label{eq:residual-attn} \\
    h'' &= \text{LayerNorm}(h' + \text{FFN}(h')) \label{eq:residual-ffn}
\end{align}

Las conexiones residuales crean un \emph{flujo residual} (\textit{residual stream})
que atraviesa todas las capas. Elhage et al.\ \cite{elhage2021mathematical}
formalizaron esta perspectiva: conceptualmente, cada capa lee del flujo, computa una
actualización, y la escribe de vuelta sumándola. La representación final de un token
es la suma acumulada de todas las contribuciones:
\begin{equation}
    x_L = x_0 + \sum_{l=1}^{L} \Delta_l^{\text{attn}} + \sum_{l=1}^{L} \Delta_l^{\text{ffn}}
    \label{eq:residual-stream}
\end{equation}

Esta arquitectura tiene una implicación crucial para la compresión: cuando se comprime
una capa intermedia, la señal corrupta se propaga por el flujo residual a todas las
capas posteriores, generando un \emph{efecto cascada} que amplifica el daño.

% =============================================================
\subsection{BERT}
\label{sec:sota-bert}
% =============================================================

BERT (\textit{Bidirectional Encoder Representations from Transformers})
\cite{devlin2019bert} demostró el poder del paradigma \textit{pre-train + fine-tune}:
un Transformer encoder pre-entrenado sobre grandes corpus produce representaciones
transferibles a tareas específicas con un \emph{fine-tuning} mínimo.

BERT-base consta de $L = 12$ capas de encoder, con $d = 768$, $H = 12$ cabezas
($d_k = 64$), y $d_{ff} = 3072$, sumando 110 millones de parámetros. De estos,
84{,}9 millones corresponden a las 72 capas lineales del encoder (6 matrices por
capa $\times$ 12 capas), que son el objetivo de la compresión SVD.

% -------------------------------------------------------------
\subsubsection{Pre-entrenamiento bidireccional}
\label{sec:sota-bert-pretraining}
% -------------------------------------------------------------

BERT se pre-entrena sobre BookCorpus + Wikipedia (3.300 millones de palabras)
mediante \textit{Masked Language Modeling} (MLM): se enmascara aleatoriamente
el 15\% de los tokens y el modelo debe predecirlos a partir del contexto
bidireccional. A diferencia de los modelos autoregresivos como GPT
\cite{radford2018improving}, que solo ven el contexto izquierdo, BERT accede
simultáneamente al contexto izquierdo y derecho, lo que produce representaciones
más ricas. Para la clasificación de emociones, donde el tono emocional puede
depender de partículas negativas al final de la frase (``no estoy enfadado
en absoluto''), esta bidireccionalidad es especialmente valiosa.

BERT también incluye un segundo objetivo de pre-entrenamiento, \textit{Next Sentence
Prediction} (NSP), que entrena al modelo a predecir si dos frases son consecutivas
en el texto original. Aunque trabajos posteriores como RoBERTa \cite{liu2019roberta}
demostraron que NSP no es necesario e incluso puede ser contraproducente, forma
parte del diseño original de BERT-base, que es el modelo utilizado en este trabajo.

\begin{figure}[H]
    \centering
    \shorthandoff{><}
    \begin{tikzpicture}[
        tok/.style={draw, rounded corners=2pt, minimum height=0.65cm, minimum width=1.3cm,
                    font=\small},
        inp/.style={tok, fill=blue!6},
        mask/.style={tok, fill=red!12},
        pred/.style={tok, fill=green!10, dashed},
        arr/.style={->, >=stealth, thick, gray},
    ]
        \node[inp] (t1) at (0,0) {\texttt{[CLS]}};
        \node[inp] (t2) at (1.6,0) {I};
        \node[inp] (t3) at (3.2,0) {really};
        \node[mask] (t4) at (4.8,0) {\texttt{[MASK]}};
        \node[inp] (t5) at (6.4,0) {this};
        \node[inp] (t6) at (8.0,0) {movie};

        \node[draw, thick, rounded corners, fill=gray!4, minimum width=9.8cm,
              minimum height=1.0cm, font=\small] (bert) at (4.0, 1.4) {BERT Encoder (12 capas)};

        \node[pred] (p4) at (4.8, 2.8) {love};

        \foreach \i in {t1,t2,t3,t4,t5,t6} {
            \draw[arr] (\i.north) -- (\i.north |- bert.south);
        }
        \draw[arr, thick, green!50!black] (bert.north -| t4) -- (p4.south);

        \node[font=\scriptsize, red!60!black, below=0.08cm of t4] {enmascarado};
        \node[font=\scriptsize, green!50!black, above=0.08cm of p4] {predecir original};
    \end{tikzpicture}
    \shorthandon{><}
    \caption{Masked Language Modeling. El modelo predice los tokens enmascarados
    usando el contexto bidireccional completo.}
    \label{fig:mlm}
\end{figure}

% -------------------------------------------------------------
\subsubsection{El token \texttt{[CLS]} y el \emph{fine-tuning}}
\label{sec:sota-cls}
% -------------------------------------------------------------

BERT añade un token especial \texttt{[CLS]} al inicio de cada secuencia. Al no
corresponder a ninguna palabra real, \texttt{[CLS]} no tiene significado léxico
propio y funciona como recolector de información global: a lo largo de las 12 capas,
acumula información de toda la secuencia a través del mecanismo de atención. Su
representación en la última capa se usa como input al clasificador de la tarea.

El \emph{fine-tuning} consiste en añadir una capa lineal sobre \texttt{[CLS]} y entrenar
el modelo completo sobre la tarea objetivo. Esto transfiere el conocimiento
lingüístico del pre-entrenamiento a la tarea específica con un coste
significativamente menor que entrenar desde cero.

La propiedad de \texttt{[CLS]} como agregador lo convierte en el punto de lectura
natural para los \textit{probing classifiers} de este trabajo: su activación en
cada capa intermedia captura el estado de procesamiento del modelo en ese punto de
la jerarquía.

% =============================================================
\subsection{Fundamentos de la descomposición en valores singulares}
\label{sec:sota-svd-teoria}
% =============================================================

La SVD es la herramienta matemática central de este proyecto. Esta sección
desarrolla sus fundamentos teóricos y la intuición geométrica que conecta la
estructura espectral con la función computacional.

% -------------------------------------------------------------
\subsubsection{Definición y propiedades}
\label{sec:sota-svd-definicion}
% -------------------------------------------------------------

Toda matriz $W \in \mathbb{R}^{m \times n}$ admite una descomposición:
\begin{equation}
    W = U \Sigma V^T
    \label{eq:svd-completa}
\end{equation}
donde $U \in \mathbb{R}^{m \times m}$ y $V \in \mathbb{R}^{n \times n}$ son
matrices ortogonales, y $\Sigma \in \mathbb{R}^{m \times n}$ tiene los
\emph{valores singulares} $\sigma_1 \geq \sigma_2 \geq \cdots \geq \sigma_r > 0$
en la diagonal ($r = \text{rank}(W)$).

Las columnas de $U$ (\emph{vectores singulares izquierdos}) forman una base del
espacio de salida; las de $V$ (\emph{vectores singulares derechos}), del espacio de
entrada. Los valores singulares cuantifican la ganancia de la transformación en
cada dirección. Se relacionan con los autovalores de $W^T W$ mediante:
\begin{equation}
    \sigma_i = \sqrt{\lambda_i(W^T W)}
    \label{eq:svd-eigenvalues}
\end{equation}

% -------------------------------------------------------------
\subsubsection{Interpretación geométrica}
\label{sec:sota-svd-geometria}
% -------------------------------------------------------------

La SVD descompone cualquier transformación lineal en tres operaciones elementales:
$V^T$ rota el espacio de entrada alineando los ejes con las direcciones naturales
de la transformación; $\Sigma$ escala cada eje por su valor singular; y $U$ rota
el resultado al espacio de salida. Visualmente, $W$ convierte la esfera unitaria
en un elipsoide cuyos semiejes tienen longitudes $\sigma_1, \sigma_2, \ldots$

\begin{figure}[H]
    \centering
    \shorthandoff{><}
    \begin{tikzpicture}[scale=0.85]
        % Unit circle
        \begin{scope}[shift={(-4,0)}]
            \draw[thick, blue!50] (0,0) circle (1);
            \draw[->, thick, red!60] (0,0) -- (1,0) node[right, font=\scriptsize] {$v_1$};
            \draw[->, thick, green!50!black] (0,0) -- (0,1) node[above, font=\scriptsize] {$v_2$};
            \node[below=0.9cm, font=\small] {Entrada};
            \node[above=1.3cm, font=\small] {$V^T$: rotar};
        \end{scope}

        \draw[->, ultra thick, gray!50] (-2.3,0) -- (-1.3,0);

        \begin{scope}[shift={(1,0)}]
            \draw[thick, blue!50] (0,0) ellipse (1.8 and 0.6);
            \draw[->, thick, red!60] (0,0) -- (1.8,0) node[right, font=\scriptsize] {$\sigma_1$};
            \draw[->, thick, green!50!black] (0,0) -- (0,0.6) node[above, font=\scriptsize] {$\sigma_2$};
            \node[below=0.9cm, font=\small] {Escalado};
            \node[above=1.3cm, font=\small] {$\Sigma$: escalar};
        \end{scope}

        \draw[->, ultra thick, gray!50] (3.5,0) -- (4.5,0);

        \begin{scope}[shift={(6.5,0)}, rotate=30]
            \draw[thick, blue!50] (0,0) ellipse (1.8 and 0.6);
            \draw[->, thick, red!60] (0,0) -- (1.8,0) node[right, font=\scriptsize] {$u_1$};
            \draw[->, thick, green!50!black] (0,0) -- (0,0.6) node[above left, font=\scriptsize] {$u_2$};
        \end{scope}
        \node[below=0.9cm, font=\small] at (5.5,0) {Salida};
        \node[above=1.3cm, font=\small] at (5.5,0) {$U$: rotar};
    \end{tikzpicture}
    \shorthandon{><}
    \caption{Interpretación geométrica de la SVD. La transformación $W$ convierte la
    esfera unitaria en un elipsoide mediante rotación ($V^T$), escalado ($\Sigma$)
    y rotación ($U$). Los semiejes del elipsoide son los valores singulares.}
    \label{fig:svd-geometric}
\end{figure}

En el contexto de BERT, las direcciones con valores singulares grandes corresponden
a las transformaciones que el modelo utiliza con más intensidad. Las de valores
pequeños pueden representar ruido, redundancia, o señales funcionales codificadas
con baja energía pero alta relevancia.

% -------------------------------------------------------------
\subsubsection{Aproximación de bajo rango y el teorema de Eckart-Young}
\label{sec:sota-eckart-young}
% -------------------------------------------------------------

La \emph{aproximación truncada de rango $k$} conserva los $k$ primeros valores
singulares:
\begin{equation}
    W_k = U_k \Sigma_k V_k^T = \sum_{i=1}^{k} \sigma_i\, u_i v_i^T
    \label{eq:svd-truncada-teoria}
\end{equation}

El \emph{teorema de Eckart-Young-Mirsky} \cite{eckart1936approximation} establece
que esta aproximación es óptima bajo la norma de Frobenius:
\begin{equation}
    W_k = \arg\min_{\text{rank}(\hat{W}) \leq k} \| W - \hat{W} \|_F
    \label{eq:eckart-young}
\end{equation}
con error:
\begin{equation}
    \| W - W_k \|_F^2 = \sum_{i=k+1}^{r} \sigma_i^2
    \label{eq:error-svd}
\end{equation}

Intuitivamente, la optimalidad se debe a que la norma de Frobenius es invariante
bajo rotaciones ortogonales ($\|W\|_F^2 = \sum_i \sigma_i^2$), de modo que retener
los $k$ valores singulares más grandes maximiza la energía capturada.

Sin embargo, la optimalidad es respecto a una métrica algebraica, no funcional. Dos
matrices con el mismo error de Frobenius pueden producir efectos radicalmente
distintos en el rendimiento del modelo, dependiendo de si los valores singulares
eliminados codifican redundancia o información crítica. Esta discrepancia es lo que
motiva el uso de interpretabilidad para guiar la compresión.

% -------------------------------------------------------------
\subsubsection{Energía espectral y compresión efectiva}
\label{sec:sota-energia-espectral}
% -------------------------------------------------------------

La fracción de energía retenida mide la calidad de la aproximación:
\begin{equation}
    E(k) = \frac{\sum_{i=1}^{k} \sigma_i^2}{\sum_{i=1}^{r} \sigma_i^2}
    \label{eq:energia}
\end{equation}

Matrices con decaimiento espectral rápido admiten compresión agresiva; matrices
con espectro plano resisten cualquier truncamiento.

\begin{figure}[H]
    \centering
    \shorthandoff{><}
    \begin{tikzpicture}[scale=0.85]
        \draw[->, >=stealth, thick] (0,0) -- (10.5,0) node[right, font=\small] {$i$};
        \draw[->, >=stealth, thick] (0,0) -- (0,5) node[above, font=\small] {$\sigma_i^2$};

        % Retained
        \foreach \x/\h in {0.4/4.5, 1.0/3.8, 1.6/3.0, 2.2/2.3, 2.8/1.7, 3.4/1.2} {
            \fill[green!25] (\x,0) rectangle (\x+0.45, \h);
            \draw[green!40!black] (\x,0) rectangle (\x+0.45, \h);
        }
        % Discarded
        \foreach \x/\h in {4.0/0.85, 4.6/0.6, 5.2/0.45, 5.8/0.32, 6.4/0.22, 7.0/0.15, 7.6/0.1, 8.2/0.07, 8.8/0.04, 9.4/0.02} {
            \fill[red!15] (\x,0) rectangle (\x+0.45, \h);
            \draw[red!30] (\x,0) rectangle (\x+0.45, \h);
        }

        \draw[dashed, thick, gray!60] (3.95, -0.3) -- (3.95, 5.0);
        \node[font=\small, gray] at (3.75, 5.3) {rango $k$};
        \node[font=\small, green!40!black] at (2.0, -0.7) {retenido};
        \node[font=\small, red!40!black] at (7.0, -0.7) {descartado};

        \draw[decorate, decoration={brace, amplitude=5pt, mirror}, thick, green!40!black]
            (0.3, -1.1) -- (3.95, -1.1) node[midway, below=6pt, font=\scriptsize] {$E(k) \approx 95\%$};
    \end{tikzpicture}
    \shorthandon{><}
    \caption{Espectro de valores singulares y truncamiento. Los primeros $k$ valores
    singulares capturan la mayor parte de la energía. La pregunta central de este
    trabajo es si la energía descartada contiene redundancia o señal funcional.}
    \label{fig:svd-spectrum}
\end{figure}

En la práctica, la aproximación reemplaza una capa $W$ de dimensión $m \times n$
por dos capas: $V_k^T$ ($k \times n$) y $U_k \Sigma_k$ ($m \times k$), reduciendo
los parámetros de $mn$ a $k(m+n)$. La compresión solo es efectiva cuando
$k < mn/(m+n)$, lo que da rangos de break-even de 384 para las matrices de atención
($768 \times 768$) y 614 para las FFN ($768 \times 3072$).

% =============================================================
\subsection{Clasificación de emociones en NLP}
\label{sec:sota-emociones}
% =============================================================

% -------------------------------------------------------------
\subsubsection{Del análisis de sentimientos a la detección de emociones}
\label{sec:sota-sentiment}
% -------------------------------------------------------------

El análisis de sentimientos, clasificar textos como positivos, negativos o
neutros \cite{pang2008opinion}, fue una de las primeras aplicaciones exitosas
del NLP, pero su granularidad de tres categorías es insuficiente para muchas
aplicaciones. La \emph{detección de emociones} aborda esta limitación usando
taxonomías más ricas, como las seis emociones básicas de Ekman
\cite{ekman1992argument} o el modelo de ocho emociones de Plutchik
\cite{plutchik1980general}. El dataset GoEmotions \cite{demszky2020goemotions},
con 28 categorías emocionales (27 emociones + \textit{neutral}) sobre 58.009
comentarios de Reddit, representa el estándar actual para esta tarea.

% -------------------------------------------------------------
\subsubsection{Clasificación multi-etiqueta}
\label{sec:sota-multilabel}
% -------------------------------------------------------------

Un mismo texto puede expresar múltiples emociones simultáneamente, lo que convierte
el problema en \emph{clasificación multi-etiqueta}. A diferencia de la clasificación
multi-clase, donde softmax distribuye una probabilidad total de 1 entre las clases
(forzando que una domine), la clasificación multi-etiqueta aplica sigmoid a cada
clase independientemente, permitiendo múltiples clases activas. La función de
pérdida correspondiente es la \emph{binary cross-entropy}:
\begin{equation}
    \mathcal{L} = -\frac{1}{C} \sum_{c=1}^{C} \left[ y_c \log(\hat{p}_c) + (1 - y_c) \log(1 - \hat{p}_c) \right]
    \label{eq:bce}
\end{equation}
donde $C = 23$, $y_c \in \{0, 1\}$, y $\hat{p}_c = \sigma(z_c)$.

\begin{figure}[H]
    \centering
    \shorthandoff{><}
    \begin{tikzpicture}[
        box/.style={draw, rounded corners=2pt, minimum height=0.5cm,
                    minimum width=1.3cm, font=\small, inner sep=3pt},
        arr/.style={->, >=stealth, thick, gray},
    ]
        % Multi-class
        \node[font=\small, gray!70!black] at (-3.5, 2.8) {Multi-clase};
        \node[box, fill=blue!6] (mc-logit) at (-3.5, 2.1) {logits $z$};
        \node[box, fill=orange!8] (mc-soft) at (-3.5, 1.2) {softmax};
        \node[box, fill=green!8] (mc-p1) at (-5.2, 0.1) {\small $p_1{=}0{,}7$};
        \node[box, fill=gray!5] (mc-p2) at (-3.5, 0.1) {\small $p_2{=}0{,}2$};
        \node[box, fill=gray!5] (mc-p3) at (-1.8, 0.1) {\small $p_3{=}0{,}1$};
        \node[font=\scriptsize, gray] at (-3.5, -0.5) {$\sum p_i = 1$};

        \draw[arr] (mc-logit) -- (mc-soft);
        \draw[arr] (mc-soft) -- (mc-p1);
        \draw[arr] (mc-soft) -- (mc-p2);
        \draw[arr] (mc-soft) -- (mc-p3);

        % Multi-label
        \node[font=\small, gray!70!black] at (3.5, 2.8) {Multi-etiqueta};
        \node[box, fill=blue!6] (ml-logit) at (3.5, 2.1) {logits $z$};
        \node[box, fill=orange!8] (ml-s1) at (1.8, 1.2) {$\sigma(z_1)$};
        \node[box, fill=orange!8] (ml-s2) at (3.5, 1.2) {$\sigma(z_2)$};
        \node[box, fill=orange!8] (ml-s3) at (5.2, 1.2) {$\sigma(z_3)$};
        \node[box, fill=green!8] (ml-p1) at (1.8, 0.1) {\small $p_1{=}0{,}9$};
        \node[box, fill=green!8] (ml-p2) at (3.5, 0.1) {\small $p_2{=}0{,}8$};
        \node[box, fill=gray!5] (ml-p3) at (5.2, 0.1) {\small $p_3{=}0{,}1$};
        \node[font=\scriptsize, gray] at (3.5, -0.5) {cada $p_i$ independiente};

        \draw[arr] (ml-logit) -- (ml-s1);
        \draw[arr] (ml-logit) -- (ml-s2);
        \draw[arr] (ml-logit) -- (ml-s3);
        \draw[arr] (ml-s1) -- (ml-p1);
        \draw[arr] (ml-s2) -- (ml-p2);
        \draw[arr] (ml-s3) -- (ml-p3);

        \draw[dashed, gray!40] (0, -0.7) -- (0, 3.1);
    \end{tikzpicture}
    \shorthandon{><}
    \caption{Multi-clase vs.\ multi-etiqueta. Softmax fuerza que una clase domine;
    sigmoid permite múltiples clases activas simultáneamente.}
    \label{fig:multiclass-multilabel}
\end{figure}

% -------------------------------------------------------------
\subsubsection{Desafíos de la clasificación emocional}
\label{sec:sota-desafios}
% -------------------------------------------------------------

La clasificación de emociones presenta desafíos específicos relevantes para
interpretar los resultados de este proyecto:

\paragraph{Ambigüedad semántica.}
Emociones como \textit{annoyance} y \textit{anger} ocupan regiones cercanas del
espacio semántico y se distinguen solo por intensidad o contexto pragmático. Esta
ambigüedad explica por qué algunas emociones necesitan más capas de procesamiento:
el modelo debe integrar información contextual que las primeras capas no capturan.

\paragraph{Desequilibrio de clases.}
En GoEmotions, \textit{admiration} tiene 4.130 ejemplos mientras que
\textit{embarrassment} tiene 303, un factor de 13{,}6$\times$. Este
desequilibrio afecta tanto al rendimiento como a la robustez frente a la compresión.

\paragraph{Señal léxica vs.\ contextual.}
Algunas emociones tienen marcadores léxicos claros (``gracias'' $\rightarrow$
\textit{gratitude}), mientras que otras dependen enteramente del contexto. Esta
distinción, que el probing captura como cristalización temprana vs.\ tardía, tiene
consecuencias directas para la compresión: las emociones léxicas son robustas
(su señal se captura en capas tempranas, redundantes y compresibles), mientras que
las contextuales dependen de las capas tardías (frágiles bajo SVD).

% =============================================================
\subsection{Compresión de modelos Transformer}
\label{sec:sota-compresion}
% =============================================================

La compresión de modelos busca reducir el tamaño y el coste computacional sin
destruir lo aprendido \cite{ganesh2021compressing}. En el contexto de los
Transformers, esto es especialmente relevante para el despliegue en entornos con
restricciones de memoria, latencia o energía \cite{strubell2019energy}.

\subsubsection{Poda}

La poda elimina pesos o estructuras completas cuya contribución es marginal
\cite{han2015deep}. Wang et al.\ \cite{wang2019structured} mostraron que es posible
eliminar hasta el 40\% de las cabezas de BERT sin pérdida significativa. Michel
et al.\ \cite{michel2019sixteen} confirmaron que en muchos casos una sola cabeza
por capa es suficiente. El inconveniente es que típicamente requiere re-entrenamiento.

\subsubsection{Cuantización}

La cuantización reduce la precisión numérica de los pesos (de float32 a int8 o
menos) \cite{jacob2018quantization}. Dettmers et al.\ \cite{dettmers2023qlora}
combinaron cuantización a 4 bits con adaptadores de bajo rango (QLoRA). La
limitación es la dependencia de hardware con soporte para operaciones de baja
precisión.

\subsubsection{Destilación}

En la destilación, un modelo pequeño aprende a imitar a uno grande. DistilBERT
\cite{sanh2019distilbert} redujo BERT a 6 capas conservando el 97\% del
rendimiento. El coste es un ciclo completo de entrenamiento.

\subsubsection{Factorización de bajo rango}

La SVD (\S\ref{sec:sota-svd-teoria}) es el método más directo. Noach y Goldberg
\cite{noach2020compressing} la aplicaron a BERT con compresión uniforme. Hsu et al.\
\cite{hsu2022language} propusieron una variante ponderada por la información de
Fisher que prioriza las direcciones más relevantes para la tarea. LoRA
\cite{hu2022lora} aborda el problema inverso: en lugar de comprimir un modelo
existente, añade adaptadores de bajo rango para \emph{fine-tuning} eficiente. La
literatura ha analizado también las propiedades estructurales de bajo rango en
los Transformers desde dos ángulos complementarios: Bhojanapalli et
al.~\cite{bhojanapalli2020low} demostraron que las cabezas de atención presentan
un \emph{cuello de botella de bajo rango} cuando la dimensión por cabeza es
inferior a la longitud de secuencia, resultado que conecta directamente con la
asimetría espectral entre componentes Q/K y FFN documentada en este trabajo. Por su parte,
Wang et al.~\cite{wang2020linformer} (Linformer) explotaron explícitamente la
naturaleza intrínsecamente de bajo rango de la matriz de atención
$\text{softmax}(QK^T)$ para reducir su complejidad de $O(n^2)$ a $O(n)$. Estos
trabajos sugieren que la atención es estructuralmente comprimible, lo que es
consistente con la robustez de Q/K bajo SVD que documentaremos en el
Capítulo~\ref{sec:resultados-compresion}.

La línea de SVD para LLMs ha continuado activa en 2023--2024.
ASVD~\cite{yuan2023asvd} introduce una variante consciente de las
activaciones que pondera los valores singulares por la magnitud de las
activaciones de entrada típicas, mejorando la fidelidad respecto a la
SVD ingenua sobre LLaMA-7B y similares. SVD-LLM~\cite{wang2024svdllm}
incorpora una corrección \emph{truncation-aware} que minimiza el error
en el espacio de salida y reporta resultados competitivos hasta el 60\%
de compresión. ShortGPT~\cite{men2024shortgpt} aborda una pregunta
ortogonal pero relacionada ---qué \emph{capas enteras} son
prescindibles--- y muestra que en LLaMA-2-13B se pueden eliminar hasta
9~de las 40~capas con caída moderada. Otra línea relacionada,
SparseGPT~\cite{frantar2023sparsegpt}, demuestra que la poda no
estructurada \emph{one-shot} sin re-entrenamiento es viable a escala
de 175\,000~M de parámetros. La Tabla~\ref{tab:comparativa-tecnicas}
resume las características de las cinco familias de técnicas; los
trabajos recientes sobre SVD para LLMs operan dentro de la familia
``SVD'' pero con variantes algorítmicas más sofisticadas que la
truncación clásica empleada en este TFG.

\begin{table}[H]
    \centering
    \caption{Comparativa de técnicas de compresión para Transformers.}
    \label{tab:comparativa-tecnicas}
    \small
    \begin{tabular}{@{}l c c c c c@{}}
        \toprule
        & \textbf{Poda} & \textbf{Cuantiz.} & \textbf{Destilación} & \textbf{SVD} & \textbf{LoRA} \\
        \midrule
        Re-entrenamiento      & Sí/Parcial & No/Parcial & Sí (completo) & No\textsuperscript{*}   & Sí (parcial) \\
        Granularidad          & Media       & Baja       & Baja          & Alta & Alta \\
        Garantías teóricas    & No          & Parcial    & No            & Sí   & Parcial \\
        Aplicable post-hoc    & Parcial     & Sí         & No            & Sí   & No \\
        Hardware especial     & No          & Sí         & No            & No   & No \\
        \bottomrule
    \end{tabular}
    \\[0.15cm]{\scriptsize \textsuperscript{*}No requiere re-entrenamiento para
    su aplicación inicial; el efecto de un \emph{fine-tuning} posterior se evalúa
    empíricamente en el Capítulo~\ref{sec:sintesis}.}
\end{table}

\subsubsection{Justificación del enfoque SVD}

La SVD encaja con los objetivos de este proyecto por cuatro razones: no requiere
re-entrenamiento para su aplicación, lo que permite evaluar decenas de
configuraciones a coste computacional bajo y desacoplar el efecto de la
compresión del efecto del re-entrenamiento; ofrece granularidad máxima
(rango específico por capa); el teorema de Eckart-Young
(\S\ref{sec:sota-eckart-young}) proporciona garantías de optimalidad respecto a
la norma de Frobenius; y el espectro de valores singulares aporta
interpretabilidad directa sobre la estructura de cada matriz de pesos. La
posibilidad de combinar la SVD con un \emph{fine-tuning} posterior se evalúa
empíricamente en el Capítulo~\ref{sec:sintesis}.

Un aspecto común a todas las técnicas anteriores es que son \emph{ciegas} respecto
a la función del modelo: la poda elimina pesos de baja magnitud, la cuantización
reduce precisión, la SVD trunca valores singulares pequeños, pero ninguna utiliza
información sobre qué computan los componentes. La posibilidad de guiar la
compresión con conocimiento funcional (proteger los componentes que el modelo
realmente necesita y comprimir los redundantes) está poco explorada en la
literatura. Este trabajo investiga precisamente esta dirección, utilizando las
técnicas de interpretabilidad mecánica descritas a continuación para informar la
estrategia de compresión.

% =============================================================
\subsection{Interpretabilidad mecánica}
\label{sec:sota-interpretabilidad}
% =============================================================

La interpretabilidad mecánica busca entender \emph{cómo} procesan información las
redes neuronales, identificando los circuitos internos responsables de
comportamientos específicos. A diferencia de métodos post-hoc como LIME
\cite{ribeiro2016should} o SHAP \cite{lundberg2017unified}, que tratan al modelo
como caja negra, la interpretabilidad mecánica abre el modelo y examina sus
componentes.

% -------------------------------------------------------------
\subsubsection{Probing de representaciones}
\label{sec:sota-probing}
% -------------------------------------------------------------

Los \textit{probing classifiers} \cite{conneau2018cram, hewitt2019structural} son
clasificadores simples (regresiones logísticas) entrenados sobre las representaciones
intermedias de un modelo para detectar qué información está codificada en cada capa.

La restricción a clasificadores lineales es deliberada: un probe no-lineal complejo
podría extraer información de cualquier representación, pero un probe lineal solo
detecta información que ya está separada linealmente en el espacio de activaciones.
Si un probe lineal clasifica una propiedad con alto F1 a partir de la capa $l$, esa
propiedad está codificada de forma \emph{linealmente accesible} en dicha capa.

Tenney et al.\ \cite{tenney2019bert} documentaron una jerarquía en BERT: capas
tempranas codifican información léxica, intermedias capturan relaciones semánticas,
y tardías se especializan en la tarea. Hewitt y Liang \cite{hewitt2019designing}
establecieron que los probes deben mantenerse simples ($C = 1{,}0$, sin tuning)
para medir la representación, no la capacidad del clasificador.

\begin{figure}[H]
    \centering
    \shorthandoff{><}
    \begin{tikzpicture}[
        layer/.style={draw, rounded corners=2pt, minimum height=0.7cm,
                      minimum width=1.1cm, font=\scriptsize, anchor=center},
        probe/.style={draw, rounded corners=2pt, minimum height=0.55cm,
                      minimum width=1.1cm, font=\scriptsize, fill=orange!10,
                      anchor=center},
        arr/.style={->, >=stealth, thick},
        probarr/.style={->, >=stealth, thin, gray!60},
        f1lab/.style={font=\tiny, gray!60!black},
    ]
        % === BERT layers (horizontal chain) ===
        \node[layer, fill=green!8]  (emb) at (0,    0) {Emb};
        \node[layer, fill=blue!6]   (l0)  at (1.4,  0) {L0};
        \node[layer, fill=blue!6]   (l1)  at (2.8,  0) {L1};
        \node[font=\small, gray]          at (4.0,  0) {$\cdots$};
        \node[layer, fill=blue!6]   (l8)  at (5.2,  0) {L8};
        \node[font=\small, gray]          at (6.4,  0) {$\cdots$};
        \node[layer, fill=blue!8]   (l11) at (7.6,  0) {L11};
        \node[layer, fill=teal!10]  (cls) at (9.4,  0) {Clasif.};

        % === Main flow arrows ===
        \draw[arr] (emb) -- (l0);
        \draw[arr] (l0)  -- (l1);
        \draw[arr] (l1)  -- (3.4, 0);
        \draw[arr] (4.5, 0) -- (l8);
        \draw[arr] (l8)  -- (5.9, 0);
        \draw[arr] (6.9, 0) -- (l11);
        \draw[arr] (l11) -- (cls);

        % === [CLS] extractions + probes ===
        % Emb
        \draw[probarr] (emb.south) -- ++(0, -0.4);
        \node[probe] (p0) at (0, -1.1) {Probe};
        \node[f1lab] at (0, -1.7) {F1 = 0};

        % L0
        \draw[probarr] (l0.south) -- ++(0, -0.4);
        \node[probe] (p1) at (1.4, -1.1) {Probe};
        \node[f1lab] at (1.4, -1.7) {F1 = 0{,}35};

        % L1
        \draw[probarr] (l1.south) -- ++(0, -0.4);
        \node[probe] (p2) at (2.8, -1.1) {Probe};
        \node[f1lab] at (2.8, -1.7) {F1 = 0{,}42};

        % L8
        \draw[probarr] (l8.south) -- ++(0, -0.4);
        \node[probe] (p3) at (5.2, -1.1) {Probe};
        \node[f1lab] at (5.2, -1.7) {F1 = 0{,}53};

        % L11
        \draw[probarr] (l11.south) -- ++(0, -0.4);
        \node[probe] (p4) at (7.6, -1.1) {Probe};
        \node[f1lab] at (7.6, -1.7) {F1 = 0{,}58};

        % === Labels ===
        \node[font=\scriptsize, gray!50!black, above=0.15cm of emb.north west, anchor=south west]
            {BERT encoder};

        % === Brace under probes ===
        \draw[decorate, decoration={brace, amplitude=4pt, mirror}, thick, gray!50]
            (-0.7, -2.0) -- (8.3, -2.0)
            node[midway, below=5pt, font=\scriptsize, gray!60!black]
            {¿En qué capa \emph{cristaliza} cada emoción?};
    \end{tikzpicture}
    \shorthandon{><}
    \caption{Probing de representaciones. Se entrena un clasificador lineal
    (regresión logística) sobre el token \texttt{[CLS]} de cada capa intermedia.
    El F1 del probe indica cuánta información emocional es linealmente accesible
    en esa capa. Los valores mostrados son ilustrativos.}
    \label{fig:probing-schema}
\end{figure}

En este trabajo se entrenan probes sobre \texttt{[CLS]} para cada capa y emoción.
El concepto de \emph{cristalización}, la primera capa donde una emoción alcanza
el 80\% de su F1 máximo, conecta profundidad de procesamiento con complejidad
semántica.

% -------------------------------------------------------------
\subsubsection{\emph{Activation patching}}
\label{sec:sota-patching}
% -------------------------------------------------------------

El probing es correlacional: detecta qué información \emph{está presente}, pero no
si la capa la \emph{computa} o solo la \emph{transmite} por el flujo residual. Las
\emph{intervenciones causales} resuelven esto: si manipulamos un componente y
observamos un cambio en el comportamiento, ese componente es causalmente responsable.

El \emph{análisis de mediación causal} fue introducido en NLP por Vig et
al.~\cite{vig2020causal} para estudiar sesgos de género en modelos de
lenguaje, formalizando la distinción entre efecto directo (la
información fluye hacia la salida sin pasar por el componente
intervenido) y efecto indirecto. Sobre esta base, Meng et al.\
\cite{meng2022locating} introdujeron el \emph{causal tracing}:
(1)~ejecutar el modelo normalmente (\emph{clean run}), (2)~corromper
las representaciones (\emph{corrupted run}), y (3)~restaurar
selectivamente un componente (\emph{patching}) para ver si el
comportamiento se recupera. Descubrieron que los hechos factuales se
localizan en las capas FFN, no en la atención. Conmy et al.\
\cite{conmy2023towards} formalizaron el framework para descubrir
circuitos computacionales completos.

\begin{figure}[H]
    \centering
    \shorthandoff{><}
    \begin{tikzpicture}[
        box/.style={draw, rounded corners=3pt, minimum height=0.85cm, minimum width=2.6cm,
                    font=\small, align=center},
        arr/.style={->, >=stealth, thick},
    ]
        \node[box, fill=green!8] (clean) at (0, 0) {Clean run\\F1 = 0{,}577};
        \node[box, fill=red!8] (corrupt) at (4.5, 0) {Corrupted\\(SVD r64)\\F1 = 0{,}000};
        \node[box, fill=blue!8] (patch) at (9, 0) {Patched\\restaurar capa $l$\\F1 = ?};

        \draw[arr] (clean) -- node[above, font=\scriptsize] {comprimir} (corrupt);
        \draw[arr] (corrupt) -- node[above, font=\scriptsize] {restaurar $l$} (patch);

        \node[font=\scriptsize, gray, below=0.25cm of patch] {Si F1 se recupera $\Rightarrow$ $l$ es causal};
    \end{tikzpicture}
    \shorthandon{><}
    \caption{\emph{Activation patching} en este trabajo. Se parte de un modelo colapsado
    por compresión (F1 = 0) y se restauran los pesos originales de un componente.
    La recuperación del F1 indica responsabilidad causal.}
    \label{fig:activation-patching}
\end{figure}

En este trabajo, la compresión SVD a rango 64 sirve como herramienta de
corrupción, lo que permite estudiar simultáneamente la localización emocional y
la interacción entre compresión y estructura funcional. El alcance
epistémico de este uso requiere una precisión: la corrupción mediante SVD
truncada no es
\emph{neutral} en sentido estadístico, a diferencia del \emph{causal tracing}
clásico de Meng et al.\ (que emplea ruido gaussiano sobre los embeddings), sino
una distorsión estructurada. Esto significa que las conclusiones de patching
extraídas en este trabajo son de naturaleza funcional (qué componentes
son suficientes para reactivar el modelo desde el colapso) más que estrictamente
causal en el sentido de Pearl. La distinción se discute con detalle en
\S\ref{sec:activation-patching}.

% -------------------------------------------------------------
\subsubsection{\emph{Logit lens} y proyecciones intermedias}
\label{sec:sota-logitlens}
% -------------------------------------------------------------

El probing detecta información presente; el \emph{activation patching} identifica
componentes suficientes. Una técnica complementaria, originalmente propuesta
por nostalgebraist~\cite{nostalgebraist2020logit} bajo el nombre de
\textit{logit lens}, aborda una pregunta intermedia: ¿qué información de las
capas intermedias es \emph{leíble} directamente por la cabeza de
clasificación entrenada del modelo? La técnica consiste en aplicar la
proyección final (la cabeza de clasificación de la última capa) a las
representaciones de capas intermedias, sin re-entrenar nada, y observar las
predicciones que produciría el modelo si tuviera que decidir desde cada capa.

La técnica permite distinguir dos fenómenos que el probing no separa:
información geométricamente alineada con la base del clasificador frente a
información presente pero \emph{rotada} hacia bases internas del modelo no
proyectables sobre el clasificador. Tiene una limitación de calibración
conocida: la cabeza fue entrenada exclusivamente sobre activaciones de la
última capa, lo que descalibra las magnitudes absolutas en capas intermedias.
limitación que Belrose et al.~\cite{belrose2023eliciting} mitigaron con
\textit{tuned lens}: una transformación lineal por capa que corrige el
desajuste de distribución. En este trabajo se aplica el \emph{logit lens} estándar
con interpretación cualitativa de los patrones (\S\ref{sec:logit-lens}); la
extensión a \emph{tuned lens} se discute como trabajo futuro
(\S\ref{sec:trabajo-futuro}).

% -------------------------------------------------------------
\subsubsection{Ablación de cabezas de atención}
\label{sec:sota-ablacion}
% -------------------------------------------------------------

Voita et al.\ \cite{voita2019analyzing} identificaron cabezas especializadas en
relaciones posicionales, sintácticas y de tokens infrecuentes. Clark et al.\
\cite{clark2019what} documentaron patrones de atención a gran escala. Michel et al.\
\cite{michel2019sixteen} mostraron que la mayoría de cabezas son redundantes, pero
la importancia varía entre capas y tareas.

Este trabajo extiende la ablación al dominio emocional, anulando cada una de las
144 cabezas de BERT y midiendo el impacto por emoción, construyendo una taxonomía
bidimensional basada en importancia y selectividad emocional.

% -------------------------------------------------------------
\subsubsection{Análisis de neuronas FFN}
\label{sec:sota-neuronas}
% -------------------------------------------------------------

Geva et al.\ \cite{geva2021transformer} propusieron que las capas FFN funcionan
como memorias clave-valor. Dai et al.\ \cite{dai2022knowledge} identificaron
\textit{knowledge neurons} que se activan selectivamente para hechos factuales
específicos: suprimirlas o amplificarlas modifica el conocimiento del modelo.

Este trabajo extiende el concepto al dominio emocional, analizando las 36.864
neuronas FFN de BERT mediante scores de selectividad (Cohen's $d$) e investigando
si la compresión SVD daña selectivamente las neuronas especializadas.

% -------------------------------------------------------------
\subsubsection{Superposition y sus implicaciones para la compresión}
\label{sec:sota-superposition}
% -------------------------------------------------------------

Elhage et al.\ \cite{elhage2022toy} formalizaron el concepto de \emph{superposition}:
cuando el modelo necesita representar más conceptos que dimensiones disponibles
($d = 768$), distribuye la información en direcciones casi-ortogonales, aceptando
pequeñas interferencias a cambio de mayor capacidad representacional. En alta
dimensión, es posible colocar del orden de $d^2$ vectores con ángulos mutuos
cercanos a 90 grados.

Esto tiene consecuencias directas para la compresión SVD: si el modelo codifica
información funcional en direcciones de baja energía, la SVD podría destruir señales
que la norma de Frobenius no identifica como importantes. Este trabajo investiga
empíricamente esta tensión. Los resultados
(Capítulo~\ref{sec:resultados-interpretabilidad}) revelan que la SVD no discrimina
selectivamente entre señales emocionales, pero las emociones con codificación
neuronal más intensa son más vulnerables a la perturbación general de los pesos.

% =============================================================
\subsection{Posicionamiento de este trabajo}
\label{sec:sota-posicionamiento}
% =============================================================

La revisión anterior identifica tres huecos que este trabajo cubre.

\paragraph{En compresión.} Los trabajos de compresión SVD en Transformers
\cite{noach2020compressing, hsu2022language} analizan la compresión de forma
agregada, sin distinguir entre componentes (Q, K, V, FFN) ni profundidades.
Ninguno complementa los resultados con evidencia causal.

\paragraph{En interpretabilidad.} Las técnicas de probing
\cite{tenney2019bert}, activation patching \cite{meng2022locating}, ablación
\cite{voita2019analyzing} y logit lens \cite{nostalgebraist2020logit} se han
aplicado mayoritariamente a tareas sintácticas o factuales (sujeto-verbo,
recuperación de hechos, generación). La aplicación sistemática y coordinada al
dominio emocional, y específicamente la triangulación entre cinco técnicas
sobre el mismo modelo y dataset, está poco explorada.

\paragraph{En la intersección.} No existe, hasta donde sabemos, un trabajo que use
la interpretabilidad mecánica para guiar la compresión. La conexión cuantitativa entre
cristalización y vulnerabilidad a la SVD es una contribución original.

La Tabla~\ref{tab:posicionamiento} resume el posicionamiento.

\begin{table}[H]
    \centering
    \caption{Posicionamiento respecto al estado del arte.}
    \label{tab:posicionamiento}
    \small
    \begin{tabular}{@{}l c c c c@{}}
        \toprule
        & Compr. & Interp. & Desagregado & Causal \\
        \midrule
        Noach \& Goldberg (2020) & SVD & --- & No & No \\
        Hsu et al.\ (2022) & SVD pond. & --- & Parcial & No \\
        Tenney et al.\ (2019) & --- & Probing & Sí (capa) & No \\
        Meng et al.\ (2022) & --- & Patching & Sí (FFN) & Sí \\
        Voita et al.\ (2019) & --- & Ablación & Sí (cabezas) & Parcial \\
        Clark et al.\ (2019) & --- & Atención & Sí (cabezas) & No \\
        Dai et al.\ (2022) & --- & Neuronas & Sí (FFN) & Sí \\
        \midrule
        \textit{Este trabajo} & \textit{SVD select.} & \textit{5 técnicas} & \textit{Sí (comp.+capa+em.)} & \textit{Sí} \\
        \bottomrule
    \end{tabular}
\end{table}
% =============================================================
% CAPÍTULO 3 — Metodología
% Versión limpia: sin negritas en cuerpo de texto, con figura
% de pipeline y expansiones integradas.
%
% Requisitos en el preámbulo (añadir si no están):
%   \usepackage{tikz}
%   \usetikzlibrary{arrows.meta, positioning, fit, backgrounds, calc}
% =============================================================

\section{Metodología}
\label{sec:metodologia-tfg}

Este capítulo describe el diseño experimental del proyecto: el dataset y su
preprocesamiento (\S\ref{sec:met-dataset}), el modelo base y su \emph{fine-tuning}
(\S\ref{sec:met-modelo}), el módulo de compresión SVD (\S\ref{sec:met-svd}),
las cuatro técnicas de interpretabilidad mecánica empleadas
(\S\ref{sec:met-interpretabilidad}) y las métricas de evaluación
(\S\ref{sec:met-metricas}). La Figura~\ref{fig:pipeline} muestra el flujo
experimental completo.

% ── Figura de pipeline ──────────────────────────────────────

\begin{figure}[H]
    \centering
    \begin{tikzpicture}[
        box/.style={
            rectangle, rounded corners=3pt, draw=black!35,
            minimum height=0.9cm, minimum width=2.6cm,
            align=center, font=\small,
            inner xsep=8pt, inner ysep=5pt
        },
        smallbox/.style={
            rectangle, rounded corners=2pt, draw=black!25,
            minimum height=0.7cm, align=center, font=\scriptsize,
            inner sep=4pt, fill=gray!6
        },
        arr/.style={-{Stealth[length=4pt]}, semithick, black!45},
        dasharr/.style={-{Stealth[length=4pt]}, thin, black!30, dashed},
        lbl/.style={font=\scriptsize\itshape, text=black!40},
    ]

    % Row 1 (y=0): GoEmotions → BERT → Baseline
    \node[box, fill=gray!8]  (data)     at (0,   0) {GoEmotions\\[-1pt]{\scriptsize 23 emociones}};
    \node[box, fill=blue!6]  (model)    at (3.6, 0) {BERT-base\\[-1pt]{\scriptsize Fine-tuning}};
    \node[box, fill=teal!8]  (baseline) at (7.2, 0) {Baseline\\[-1pt]{\scriptsize F1\textsubscript{macro} = 0{,}577}};

    \draw[arr] (data) -- (model);
    \draw[arr] (model) -- (baseline);

    % Row 2 (y=-1.6): Compresión SVD
    \node[box, fill=orange!8] (svd) at (3.6, -1.6) {Compresión SVD\\[-1pt]{\scriptsize Uniforme / Adaptativa}};
    \draw[arr] (baseline.south) -- ++(0,-0.35) -| (svd.north);

    % Row 3 (y=-3.5): Interpretabilidad mecánica — manual rectangle
    \draw[rounded corners=4pt, draw=black!15, fill=blue!2]
        (-1.6, -2.85) rectangle (9.0, -4.55);
    \node[font=\small, text=black!60] at (3.6, -3.15) {Interpretabilidad mecánica};

    \node[smallbox] (probing)  at (0,   -3.85) {Probing\\[-1pt]{\tiny NB04 · capa}};
    \node[smallbox] (patching) at (2.6, -3.85) {Act.\ patching\\[-1pt]{\tiny NB05 · componente}};
    \node[smallbox] (ablation) at (5.2, -3.85) {Head ablation\\[-1pt]{\tiny NB06 · cabeza}};
    \node[smallbox] (neurons)  at (7.8, -3.85) {Neuronas FFN\\[-1pt]{\tiny NB07 · neurona}};

    \draw[arr] (svd.south) -- ++(0, -0.7);

    % Row 4 (y=-5.6): Compresión informada
    \node[box, fill=teal!10] (informed) at (3.6, -5.6) {Compresión informada\\[-1pt]{\scriptsize NB09 · Greedy data-driven}};
    \draw[arr] (3.6, -4.55) -- (informed.north);

    % Row 5 (y=-7.1): Fine-tuning recovery
    \node[box, fill=green!8] (recovery) at (3.6, -7.1) {Fine-tuning recovery\\[-1pt]{\scriptsize F1 = 0{,}591 (+2{,}4\%)}};
    \draw[arr] (informed) -- (recovery);

    % Feedback loop (right side)
    \coordinate (fb1) at (9.6, -5.6);
    \coordinate (fb2) at (9.6, -1.6);
    \draw[dasharr, rounded corners=3pt]
        (informed.east) -- (fb1) -- (fb2) -- (svd.east);
    \node[lbl, rotate=90, anchor=south] at (9.6, -3.6) {retroalimentación};

    \end{tikzpicture}
    \caption{Pipeline experimental. El modelo baseline se somete a compresión SVD y a
    cuatro técnicas de interpretabilidad mecánica con granularidad creciente (de capa
    completa a neurona individual). Los hallazgos retroalimentan el diseño de
    estrategias de compresión informadas (línea discontinua). El ciclo se cierra
    con \emph{fine-tuning} del modelo comprimido.}
    \label{fig:pipeline}
\end{figure}

% -------------------------------------------------------------
\subsection{Dataset: GoEmotions}
\label{sec:met-dataset}
% -------------------------------------------------------------

El dataset GoEmotions \cite{demszky2020goemotions} contiene 58.009 comentarios de
Reddit anotados por tres evaluadores con una taxonomía de 28 categorías emocionales
(27 emociones + \textit{neutral}). A diferencia de la clasificación de sentimientos
estándar (positivo/negativo), GoEmotions es un dataset de clasificación
multi-etiqueta: cada texto puede expresar múltiples emociones simultáneamente (por
ejemplo, \textit{admiración} y \textit{alegría}), lo que impone al modelo la necesidad
de resolver co-ocurrencias con matices semánticos cercanos.

\subsubsection{Filtrado de emociones}

Del conjunto original de 28 categorías, se eliminaron cinco que presentaban problemas
de rendimiento o representatividad:

\begin{itemize}[nosep]
    \item \textit{neutral}: excluida por no ser una emoción propiamente dicha y por
    dominar el dataset (14.000+ ejemplos), lo que distorsionaba las métricas
    macro-promediadas.
    \item \textit{grief}, \textit{nervousness}, \textit{pride}, \textit{relief}:
    excluidas por producir F1 = 0 en el baseline, debido a su bajísima frecuencia en
    el dataset (menos de 100 ejemplos cada una en el conjunto de entrenamiento). La
    presencia de emociones con F1 nulo degradaba el F1 macro sin aportar información
    útil sobre la capacidad del modelo.
\end{itemize}

El filtrado se implementó a nivel del pipeline de datos: se eliminaron las cinco
columnas correspondientes del vector de etiquetas (pasando de 28 a 23 dimensiones) y
se descartaron los ejemplos cuyo vector resultante quedaba vacío (aquellos anotados
exclusivamente con alguna de las emociones excluidas). El resultado es un dataset de
23 emociones donde todas las categorías producen un F1 funcional (mínimo: 0{,}267
para \textit{embarrassment}).

La Tabla~\ref{tab:emociones-23} lista las 23 emociones del dataset final, ordenadas
por F1 en el baseline.

\begin{table}[H]
    \centering
    \caption{Las 23 emociones del dataset final, con F1 en el baseline.}
    \label{tab:emociones-23}
    \small
    \begin{tabular}{@{}l r r l@{}}
        \toprule
        Emoción & F1 baseline & Ejs.\ train & Grupo semántico \\
        \midrule
        gratitude      & 0{,}927 & 2\,662 & Positiva social \\
        amusement      & 0{,}853 & 2\,328 & Positiva energética \\
        love           & 0{,}816 & 2\,086 & Positiva social \\
        admiration     & 0{,}757 & 4\,130 & Positiva social \\
        fear           & 0{,}710 & \phantom{0\,}596 & Negativa interna \\
        curiosity      & 0{,}707 & 2\,191 & Epistémica \\
        remorse        & 0{,}672 & \phantom{0\,}545 & Negativa interna \\
        joy            & 0{,}635 & 1\,452 & Positiva energética \\
        sadness        & 0{,}632 & 1\,326 & Negativa interna \\
        optimism       & 0{,}573 & 1\,581 & Positiva orientada \\
        surprise       & 0{,}571 & 1\,060 & Epistémica \\
        approval       & 0{,}555 & 2\,939 & Positiva social \\
        anger          & 0{,}550 & 1\,567 & Negativa reactiva \\
        caring         & 0{,}542 & 1\,087 & Positiva orientada \\
        desire         & 0{,}535 & \phantom{0\,}641 & Positiva orientada \\
        disapproval    & 0{,}527 & 2\,022 & Negativa reactiva \\
        disgust        & 0{,}521 & \phantom{0\,}793 & Negativa reactiva \\
        confusion      & 0{,}515 & 1\,368 & Epistémica \\
        excitement     & 0{,}454 & \phantom{0\,}853 & Positiva energética \\
        annoyance      & 0{,}411 & 2\,470 & Negativa reactiva \\
        disappointment & 0{,}275 & 1\,269 & Negativa interna \\
        realization    & 0{,}271 & 1\,110 & Epistémica \\
        embarrassment  & 0{,}267 & \phantom{0\,}303 & Negativa interna \\
        \midrule
        \multicolumn{2}{@{}l}{F1 macro promedio} & & 0{,}577 \\
        \bottomrule
    \end{tabular}
    \\[0.15cm]{\small Grupos semánticos propuestos a partir del clustering neuronal
    (Capítulo~\ref{sec:resultados-interpretabilidad}). Fuente: elaboración propia
    sobre GoEmotions \cite{demszky2020goemotions}.}
\end{table}

\subsubsection{Splits y preprocesamiento}

Se utilizaron los splits oficiales de GoEmotions: entrenamiento (43.410 ejemplos),
validación (5.426) y test (5.427), tras el filtrado. El preprocesamiento se limita a
la tokenización con el tokenizer de BERT-base-uncased, con una longitud máxima de
128 tokens y padding dinámico. No se aplican técnicas de aumento de datos.

% -------------------------------------------------------------
\subsection{Modelo: BERT-base}
\label{sec:met-modelo}
% -------------------------------------------------------------

\subsubsection{Arquitectura}

BERT-base-uncased \cite{devlin2019bert} consta de 12 capas de encoder Transformer,
cada una con un bloque de atención multi-cabeza (12 cabezas) y un bloque
feed-forward. El modelo tiene 110 millones de parámetros, de los cuales
84{,}9 millones corresponden a las 72 capas lineales del encoder (6 matrices por
capa $\times$ 12 capas), que constituyen el objetivo de la compresión SVD. La
Tabla~\ref{tab:capas-bert} detalla las dimensiones de cada tipo de componente.

\begin{table}[H]
    \centering
    \caption{Capas lineales del encoder de BERT-base (por capa).}
    \label{tab:capas-bert}
    \small
    \begin{tabular}{@{}l l c r@{}}
        \toprule
        Bloque & Componente & Dimensiones & Parámetros \\
        \midrule
        Atención & Query ($W_Q$)            & $768 \times 768$   & 589.824 \\
        Atención & Key ($W_K$)              & $768 \times 768$   & 589.824 \\
        Atención & Value ($W_V$)            & $768 \times 768$   & 589.824 \\
        Atención & Output ($W_O$)           & $768 \times 768$   & 589.824 \\
        FFN      & Intermediate ($W_{int}$) & $768 \times 3072$  & 2.359.296 \\
        FFN      & Output ($W_{out}$)       & $3072 \times 768$  & 2.359.296 \\
        \midrule
        \multicolumn{3}{l}{Total por capa del encoder} & 7.077.888 \\
        \multicolumn{3}{l}{Total 12 capas} & 84.934.656 \\
        \bottomrule
    \end{tabular}
    \\[0.15cm]{\small Fuente: elaboración propia.}
\end{table}

\subsubsection{\emph{Fine-tuning}}

El modelo se fine-tuneó sobre el dataset GoEmotions (23 emociones) añadiendo una capa
de clasificación lineal sobre el token \texttt{[CLS]} de la última capa del encoder.
La clasificación multi-etiqueta se implementa con una función de pérdida
\texttt{BCEWithLogitsLoss} (binary cross-entropy con logits), que trata cada emoción
como una predicción binaria independiente.

Los hiperparámetros de \emph{fine-tuning} se resumen en la Tabla~\ref{tab:hyperparams}.

\begin{table}[H]
    \centering
    \caption{Hiperparámetros de \emph{fine-tuning}.}
    \label{tab:hyperparams}
    \small
    \begin{tabular}{@{}l l@{}}
        \toprule
        Parámetro & Valor \\
        \midrule
        Optimizador        & AdamW \\
        Learning rate      & $2 \times 10^{-5}$ \\
        Epochs             & 4 \\
        Batch size         & 32 \\
        Warmup             & 10\% de los steps \\
        Weight decay       & 0{,}01 \\
        Max sequence length & 128 tokens \\
        Semilla aleatoria  & 42 \\
        \bottomrule
    \end{tabular}
    \\[0.15cm]{\small Fuente: elaboración propia.}
\end{table}

El modelo resultante alcanza un F1 macro de 0{,}577 y un F1 micro de 0{,}627 sobre
el conjunto de test, lo que constituye el baseline contra el que se mide toda
degradación posterior.

% -------------------------------------------------------------
\subsection{Compresión SVD}
\label{sec:met-svd}
% -------------------------------------------------------------

\subsubsection{Fundamento teórico}

La SVD factoriza cualquier matriz $W \in \mathbb{R}^{m \times n}$ como
$W = U \Sigma V^T$, donde $U \in \mathbb{R}^{m \times m}$ y
$V \in \mathbb{R}^{n \times n}$ son matrices ortogonales y
$\Sigma \in \mathbb{R}^{m \times n}$ es diagonal con los valores singulares
$\sigma_1 \geq \sigma_2 \geq \cdots \geq \sigma_{\min(m,n)} \geq 0$. El
teorema de Eckart-Young \cite{eckart1936approximation} garantiza que la mejor
aproximación de rango $k$ en norma de Frobenius es la SVD truncada:
\begin{equation}
    W_k = \sum_{i=1}^{k} \sigma_i u_i v_i^T = U_k \Sigma_k V_k^T
    \label{eq:svd-truncada}
\end{equation}
con error $\|W - W_k\|_F^2 = \sum_{i=k+1}^{\min(m,n)} \sigma_i^2$. La SVD induce
una descomposición ortogonal del espacio de salida $\mathbb{R}^m$:
\begin{equation}
    \mathbb{R}^m = \mathcal{S}_k \oplus \mathcal{S}_k^\perp,
    \quad \mathcal{S}_k = \text{span}\{u_1, \ldots, u_k\}
    \label{eq:subespacios}
\end{equation}
donde $\mathcal{S}_k$ es el \emph{subespacio conservado} (asociado a los $k$
mayores valores singulares) y $\mathcal{S}_k^\perp$ el \emph{subespacio
eliminado}. El operador de proyección sobre $\mathcal{S}_k^\perp$ es
$P_k^\perp = I - U_k U_k^T$, y la fracción de señal $x \in \mathbb{R}^m$
eliminada por la compresión a rango $k$ es
$\|P_k^\perp x\|_2 / \|x\|_2$. Esta cantidad se utiliza en
\S\ref{sec:neuron-analysis} para evaluar si la SVD discrimina selectivamente
entre direcciones funcionalmente significativas (vectores de selectividad
neuronal) o si elimina señal uniformemente.

En la práctica, una capa lineal $W$ de dimensión $m \times n$ se reemplaza
por dos capas: $V_k^T$ ($k \times n$) y $U_k \Sigma_k$ ($m \times k$),
reduciendo los parámetros de $mn$ a $k(m+n)$.

\subsubsection{Rango de break-even}

La compresión SVD solo reduce parámetros si el rango elegido es menor que
$mn/(m+n)$. Para BERT-base:
\begin{itemize}[nosep]
    \item Componentes de atención ($768 \times 768$): break-even en rango 384.
    \item Componentes FFN ($768 \times 3072$): break-even en rango 614.
\end{itemize}
Rangos superiores a estos umbrales \emph{aumentan} el número total de parámetros.
Este detalle técnico es relevante para interpretar correctamente las estrategias
adaptativas de alta energía (Capítulo~\ref{sec:resultados-compresion}).

\subsubsection{Implementación}

El módulo de compresión (\texttt{src/compression/svd.py}) implementa una clase
\texttt{SVDLinear} que almacena las dos matrices factorizadas y realiza la
multiplicación en dos pasos durante el forward pass. La función
\texttt{apply\_svd\_compression()} acepta un modelo BERT, una especificación de
rangos por capa y componente, y devuelve el modelo comprimido. La SVD se computa
siempre en \texttt{float32} para evitar inestabilidades numéricas, independientemente
de la precisión del modelo.

\subsubsection{Estrategias de compresión evaluadas}

Se evaluaron cuatro familias de estrategias:

\begin{enumerate}[nosep]
    \item Uniforme: mismo rango $k \in \{32, 64, 128, 256, 384, 512\}$ para
    las 72 matrices del encoder.
    \item Adaptativa (energía): rango por matriz determinado por un umbral
    de energía espectral ($E(k) \geq \{80\%, 90\%, 95\%, 99\%\}$), donde
    $E(k) = \sum_{i=1}^k \sigma_i^2 / \sum_{i=1}^n \sigma_i^2$.
    \item Informada heurística: rangos asignados manualmente a partir de
    reglas derivadas de los resultados de interpretabilidad (proteger capas tardías,
    comprimir agresivamente Q/K). Tres niveles de agresividad (light, moderate,
    aggressive).
    \item Greedy (data-driven): algoritmo que construye un menú de movimientos
    de compresión (cada uno del tipo ``comprimir componente $X$ en capas $Y$ al
    rango $Z$'') y los ordena por eficiencia, definida como parámetros ahorrados
    por punto de F1 perdido. El algoritmo selecciona movimientos del más eficiente
    al menos eficiente hasta alcanzar el ratio de parámetros objetivo. El coste
    en F1 de cada movimiento proviene de datos empíricos de sensibilidad por
    componente y profundidad, obtenidos en los notebooks de compresión y
    \emph{activation patching}.
\end{enumerate}

% -------------------------------------------------------------
\subsection{Técnicas de interpretabilidad mecánica}
\label{sec:met-interpretabilidad}
% -------------------------------------------------------------

Se emplearon cuatro técnicas complementarias que cubren escalas progresivas de
granularidad: capa completa (probing), capa y componente (\emph{activation patching}),
cabeza individual (ablación) y neurona individual (análisis de selectividad).

\subsubsection{Probing lineal}
\label{sec:met-probing}

Para cada una de las 13 capas de BERT (embedding + 12 encoder layers) y cada
una de las 23 emociones, se entrena un probe de regresión logística
(\texttt{LogisticRegression} de scikit-learn, $C = 1{,}0$ sin tuning) sobre
la representación del token \texttt{[CLS]}. La decisión de fijar $C = 1{,}0$
sigue la lógica de Hewitt y Liang \cite{hewitt2019designing}: el probe debe
ser deliberadamente simple para medir la calidad de la representación, no la
capacidad del clasificador. Esta elección tiene una contrapartida que conviene
reconocer: información codificada no linealmente en capas tempranas no sería
detectada por un probe lineal, lo que podría infraestimar la presencia de
señal emocional en las primeras capas. La interpretación de los resultados
(\S\ref{sec:probing}) se realiza con esta cautela.

El total nominal de probes es $23 \times 13 = 299$; sin embargo, los probes
sobre la capa de embedding producen F1 = 0 uniformemente para todas las
emociones (el embedding crudo no contiene información emocional separable
linealmente), por lo que se descartan del análisis efectivo. El número
efectivo de probes que entran en las correlaciones y métricas es de
$23 \times 12 = 276$, valor que utilizamos consistentemente en el resto de la
memoria.

La representación utilizada es el vector del token \texttt{[CLS]} de cada
capa, en lugar del mean-pooling sobre todos los tokens. Esta decisión sigue
la convención del propio BERT: durante el \emph{fine-tuning}, el token \texttt{[CLS]}
es el único cuya representación se conecta directamente a la cabeza de
clasificación, por lo que es la posición donde el modelo concentra la
información relevante para la tarea. Utilizar mean-pooling mediría la
información distribuida en toda la secuencia, lo que reflejaría una capacidad
representacional distinta de la que el modelo realmente explota para
clasificar.

La métrica de probing es el F1 por probe. Se define la \emph{capa de
cristalización} de una emoción como la primera capa donde su F1 alcanza el
80\% del máximo observado en cualquier capa. Este umbral es una convención;
el orden de cristalización es estable con umbrales de 0{,}7 y 0{,}9.

\subsubsection{\emph{Logit lens}}
\label{sec:met-logitlens}

El \emph{logit lens} \cite{nostalgebraist2020logit} aplica la cabeza de
clasificación entrenada del modelo (\texttt{BertPooler} + clasificador
lineal) a la representación del token \texttt{[CLS]} de cada capa
intermedia, sin re-entrenamiento. Para cada capa $\ell \in \{0, \ldots, 11\}$
y cada ejemplo, se calcula:
\begin{equation}
    \hat{p}_\ell = \sigma\left(W_{\text{cls}} \cdot \tanh\left(W_{\text{pool}} \cdot h_\ell^{[\text{CLS}]}\right) + b_{\text{cls}}\right) \in \mathbb{R}^{23}
    \label{eq:logitlens}
\end{equation}
donde $h_\ell^{[\text{CLS}]}$ es la representación del token \texttt{[CLS]}
en la capa $\ell$ y $\sigma$ es la sigmoide aplicada componente a
componente. Las predicciones se agregan sobre 2\,300 ejemplos del conjunto
de test (submuestreo equilibrado entre los tres grupos de cristalización del
probing). Las métricas reportadas son la sigmoid media de la emoción
correcta (\emph{gold sigmoid}), la sigmoid media de la emoción más probable
(\emph{top-1 sigmoid}) y la suma total de las 23 sigmoids.

Esta técnica tiene una limitación de calibración conocida que se discute en
\S\ref{sec:logit-lens}: los pesos del pooler y del clasificador fueron
optimizados exclusivamente para representaciones de la última capa tras
LayerNorm, por lo que los valores absolutos en capas intermedias no son
estrictamente comparables al baseline. Las interpretaciones se realizan
sobre tendencias relativas (entre capas o entre grupos emocionales), no
sobre magnitudes absolutas.

\subsubsection{\emph{Activation patching}}
\label{sec:met-patching}

El \emph{activation patching} parte de un modelo completamente colapsado:
comprimido con SVD uniforme a rango 64 (F1 = 0{,}000). Se eligió el rango 64
como punto de colapso porque es el rango más alto que produce un colapso
total y uniforme del modelo (F1 = 0{,}000 en las 23 emociones), garantizando
un estado base completamente destruido desde el cual cualquier recuperación
es atribuible al componente restaurado. Rangos inferiores (32) producen el
mismo colapso, pero con mayor distorsión de los pesos que podría introducir
artefactos en el patching. Rangos superiores (128) permiten recuperación
parcial espontánea en algunas emociones, lo que contaminaría la atribución
causal.

La corrupción mediante SVD truncada a rango 64 no es
\emph{neutral} en sentido estadístico (a diferencia del \emph{causal tracing}
clásico de Meng et al.\ \cite{meng2022locating}, que usa ruido gaussiano
sobre los embeddings), sino una distorsión \emph{estructurada} que respeta
la geometría espectral de los pesos. Las conclusiones extraídas tienen,
por tanto, naturaleza \emph{funcional} (qué componentes son
suficientes para reactivar el modelo desde el colapso) más que estrictamente
causal en el sentido de Pearl. Esta distinción se discute con detalle en
\S\ref{sec:activation-patching}; reconocerla forma parte del rigor
metodológico del trabajo.

Desde este estado, se restauran selectivamente los pesos originales de una
capa o componente y se mide la recuperación del F1 macro. La métrica de
interés es el porcentaje de restauración:
\begin{equation}
    \text{Restauración}(\%) = \frac{F1_{\text{patched}} - F1_{\text{comprimido}}}{F1_{\text{baseline}} - F1_{\text{comprimido}}} \times 100
    \label{eq:restauracion}
\end{equation}

El patching se realiza a dos niveles de granularidad: (a) capa completa,
restaurando las 6 matrices de una capa del encoder; y (b) componente,
restaurando solo el bloque de atención o solo el bloque FFN de una capa.

\subsubsection{Ablación de cabezas de atención}
\label{sec:met-ablacion}

Se anulan individualmente cada una de las 144 cabezas de atención ($12 \times 12$)
estableciendo su salida a cero (\textit{zero-out}). Para cada cabeza se mide la
caída en F1 macro y en F1 por emoción. Se construye una taxonomía bidimensional
basada en dos métricas ortogonales:

\begin{itemize}[nosep]
    \item Importancia: caída absoluta en F1 macro al ablacionar la cabeza.
    \item Selectividad: concentración del impacto en pocas emociones, medida
    mediante el coeficiente de Gini sobre el vector de caídas por emoción.
\end{itemize}

Las cuatro categorías resultantes son: Critical Specialist (alta importancia, alta
selectividad), Critical Generalist (alta importancia, baja selectividad), Minor
Specialist (baja importancia, alta selectividad) y Dispensable (baja importancia,
baja selectividad).

\subsubsection{Análisis de neuronas FFN}
\label{sec:met-neuronas}

Para cada una de las 36.864 neuronas intermedias del modelo ($12 \times 3.072$) se
extrae la activación post-GELU mediante forward hooks, promediando sobre las posiciones
de la secuencia. Se procesa el conjunto de test completo.

La \emph{selectividad neuronal} para una neurona $n$ y una emoción $e$ se define
como un score tipo Cohen's $d$:
\begin{equation}
    s(n, e) = \frac{\bar{a}_{n,\text{presente}} - \bar{a}_{n,\text{ausente}}}{\sigma_{\text{pooled}}}
    \label{eq:selectividad}
\end{equation}
donde $\bar{a}_{n,\text{presente}}$ es la activación media de la neurona $n$ cuando
la emoción $e$ está presente en el texto, y $\bar{a}_{n,\text{ausente}}$ cuando está
ausente. Se consideran significativas las neuronas con $|s| > 2{,}0$.

Para conectar la especialización neuronal con la compresión SVD, se investiga
una hipótesis concreta: \textit{¿las emociones cuyas neuronas especializadas
caen mayoritariamente en el subespacio eliminado por la SVD son las más
vulnerables a la compresión?} Para responder a esta pregunta, se proyecta el
vector de selectividad de cada emoción (un vector de 3.072 dimensiones por capa,
con la selectividad de cada neurona) sobre los subespacios conservado y eliminado
por la SVD a rango 128, y se calcula la fracción de señal emocional que reside
en cada subespacio. Si la hipótesis es correcta, las emociones con mayor fracción
de señal en el subespacio eliminado deberían mostrar mayor caída de F1 bajo
compresión.

% -------------------------------------------------------------
\subsection{Métricas de evaluación}
\label{sec:met-metricas}
% -------------------------------------------------------------

La evaluación combina métricas a tres escalas: rendimiento global (F1 macro y
micro), impacto diferencial por emoción (F1 por emoción, retención de F1) y
calidad representacional (silhouette score).

\begin{itemize}[nosep]
    \item F1 macro: media no ponderada del F1 por emoción. Trata todas las
    emociones por igual, independientemente de su frecuencia, y es sensible a la
    degradación de emociones minoritarias. Es la métrica principal del proyecto.
    \item F1 micro: F1 calculado sobre el total de predicciones. Favorece
    las emociones frecuentes y proporciona una perspectiva complementaria.
    \item F1 por emoción: permite identificar impactos diferenciales de la
    compresión.
    \item Retención de F1: $F1_{\text{comprimido}} / F1_{\text{baseline}}$,
    expresada como porcentaje.
    \item Ratio de parámetros: $\text{params}_{\text{comprimido}} /
    \text{params}_{\text{original}}$, donde un ratio $< 1$ indica compresión
    efectiva.
    \item Silhouette score: calidad de los clusters en proyecciones t-SNE
    de embeddings \texttt{[CLS]}.
\end{itemize}

\subsection{Entorno de ejecución}
\label{sec:met-entorno}

Todos los experimentos se ejecutaron en Google Colab (GPU Tesla T4, 16~GB VRAM) y/o
en entorno local (CPU). Los notebooks son compatibles con ambos entornos. Se fijó la
semilla aleatoria (42) en PyTorch, NumPy y Python para garantizar la reproducibilidad.
El código se organizó en módulos reutilizables (\texttt{src/data},
\texttt{src/models}, \texttt{src/compression}, \texttt{src/utils}), separados de los
notebooks de experimentación (NB01--NB10).

\subsection{Reproducibilidad}
\label{sec:met-reproducibilidad}

La reproducibilidad del proyecto se aborda en cuatro niveles:

\begin{itemize}[nosep]
    \item \textbf{Código.} El repositorio contiene los nueve notebooks
    de experimentación (\texttt{notebooks/01b}--\texttt{09}) más el
    notebook de \emph{logit lens} (\texttt{notebooks/08b}), todos ejecutables
    end-to-end y con dependencias declaradas en \texttt{requirements.txt}.
    Los módulos compartidos viven en \texttt{src/} (carga de datos,
    SVD, métricas, hooks de extracción de activaciones).

    \item \textbf{Datos.} GoEmotions \cite{demszky2020goemotions} es
    público bajo licencia Apache~2.0 y se descarga vía la librería
    \texttt{datasets} de Hugging Face. Los splits oficiales se utilizan
    sin modificación más allá del filtrado de las cinco emociones
    documentado en \S\ref{sec:met-dataset}.

    \item \textbf{Modelo.} El checkpoint baseline
    (\texttt{bert-goemotions-23emo-final}, 110~M parámetros) se obtiene
    aplicando el procedimiento de \emph{fine-tuning} descrito en
    \S\ref{sec:met-modelo} sobre BERT-base-uncased
    \cite{devlin2019bert} con la semilla fija. La ejecución completa
    ocupa unas 25~minutos en una GPU T4 de Colab.

    \item \textbf{Resultados intermedios.} Cada notebook exporta sus
    resultados a \texttt{csvs/notebookN/} y \texttt{results/notebookN/}
    como ficheros CSV planos (más algunos PNG). Esto permite reproducir
    cualquier figura o tabla de la memoria sin tener que re-ejecutar
    los experimentos: basta con cargar los CSV correspondientes. Los
    análisis estadísticos posteriores (correlaciones de Spearman,
    bootstrap por emoción, intervalos de confianza al 95\%) son
    deterministas dada la semilla y operan sobre estos CSV.
\end{itemize}

La principal fuente de no-determinismo residual es el orden de los
\textit{batches} dentro de cada época de \emph{fine-tuning}, que en GPU
introduce variaciones en la cuarta cifra decimal del F1 macro entre
ejecuciones independientes. Las cifras reportadas en la memoria
corresponden a la ejecución con semilla 42, y cualquier divergencia
respecto a una réplica posterior debería caer dentro de ese margen.
% =============================================================
% CAPÍTULO 4 — Resultados de compresión SVD
% Alimentado por: NB02 (espectral, uniforme, adaptativa) y
%                 NB03 (sensibilidad por componente y profundidad)
% =============================================================

\section{Resultados de compresión SVD}
\label{sec:resultados-compresion}

Este capítulo analiza el comportamiento de BERT-base \emph{fine-tuned} sobre GoEmotions
(23 emociones) bajo distintas estrategias de compresión basadas en la descomposición
en valores singulares (SVD). El objetivo es doble: primero, determinar cuánta
información se puede eliminar sin destruir la capacidad emocional del modelo, y segundo,
identificar dónde reside esa información dentro de la arquitectura.

El capítulo se organiza en cuatro secciones. La primera examina la estructura espectral
de las 72 matrices lineales del modelo (\S\ref{sec:analisis-espectral}). La segunda
evalúa la compresión uniforme y documenta una transición de fase abrupta
(\S\ref{sec:compresion-uniforme}). La tercera desglosa la sensibilidad por tipo de
componente y por profundidad (\S\ref{sec:sensibilidad-componente}). La cuarta presenta
la compresión adaptativa como alternativa informada por la estructura espectral
(\S\ref{sec:compresion-adaptativa}).

El modelo de referencia (\emph{baseline}) a lo largo de todo el capítulo es BERT-base
\emph{fine-tuned} sin ninguna compresión, con un F1 macro de 0{,}577 sobre el conjunto de test.


% =============================================================
% §4.1 — Análisis espectral
% =============================================================

\subsection{Análisis espectral de las matrices de peso}
\label{sec:analisis-espectral}

Antes de aplicar cualquier compresión, conviene entender la estructura interna de los
pesos del modelo. Cada capa de BERT-base contiene seis matrices lineales: Query ($W_Q$),
Key ($W_K$), Value ($W_V$) y Attention Output ($W_O$) en el bloque de self-attention, y
FFN Intermediate ($W_{\text{inter}}$) y FFN Output ($W_{\text{out}}$) en el bloque
feed-forward. Con 12 capas, el modelo contiene 72 matrices lineales en total.

Para cada matriz $W \in \mathbb{R}^{m \times n}$, calculamos su descomposición en valores
singulares $W = U \Sigma V^\top$ y definimos la energía acumulada al rango $k$
como la fracción de la norma de Frobenius capturada por los primeros $k$ valores singulares:

\begin{equation}
    E(k) = \frac{\sum_{i=1}^{k} \sigma_i^2}{\sum_{i=1}^{n} \sigma_i^2}
    \label{eq:energia-acumulada}
\end{equation}

El rango mínimo $k_\tau$ para un umbral de energía $\tau$ es el menor $k$ tal
que $E(k) \geq \tau$. Un valor bajo de $k_\tau$ indica que la matriz concentra su
información en pocos componentes principales y es, por tanto, más compresible.

\subsubsection{Asimetría self-attention vs.\ feed-forward}

La Figura~\ref{fig:anatomy} muestra el mapa de rangos mínimos $k_{95}$ para cada una de
las 72 matrices, organizadas por capa (eje vertical) y tipo de componente (eje horizontal).
El resultado más inmediato es una asimetría pronunciada entre los componentes de
self-attention y los de feed-forward.

\begin{figure}[H]
    \centering
    \includegraphics[width=0.85\textwidth]{figures/cap4/cap4_anatomy_heatmap_es.png}
    \caption{Rango mínimo $k_{95}$ por capa y tipo de componente. Colores fríos indican alta compresibilidad; colores cálidos, baja. Se observa una división clara entre los componentes de self-attention (izquierda) y los feed-forward (derecha).}
    \label{fig:anatomy}
\end{figure}

La Tabla~\ref{tab:rangos-medios} cuantifica esta diferencia. Las matrices Query y Key
requieren un rango medio de aproximadamente 395 para retener el 95\% de la energía
espectral, mientras que las matrices FFN Intermediate y Output necesitan rangos superiores
a 618. El gap medio entre Q y FFN Intermediate es de unos 222 rangos, lo que equivale
al 29\% del rango completo (768).

\begin{table}[H]
    \centering
    \caption{Rango medio $k_{95}$ por tipo de componente (media sobre las 12 capas).}
    \label{tab:rangos-medios}
    \small
    \begin{tabular}{@{}lcc@{}}
        \toprule
        \textbf{Componente} & $\boldsymbol{\mu(k_{95})}$ & \textbf{Compresibilidad} \\
        \midrule
        Query (Q)           & $\sim$396 & Alta \\
        Key (K)             & $\sim$394 & Alta \\
        Value (V)           & $\sim$424 & Alta \\
        Attn Output         & $\sim$434 & Media-alta \\
        FFN Intermediate    & $\sim$618 & Baja \\
        FFN Output          & $\sim$625 & Baja \\
        \bottomrule
    \end{tabular}
\end{table}

Esta asimetría es consistente a lo largo de las 12 capas: no se trata de un artefacto
puntual de una o dos capas anómalas, sino de una propiedad estructural del modelo
\emph{fine-tuned}. La interpretación es directa: las matrices Q y K concentran los patrones de
atención en un subespacio de dimensionalidad efectiva mucho menor que 768, mientras que
las matrices FFN distribuyen su información de forma más uniforme a lo largo de todo el
espectro de valores singulares.

\subsubsection{Perfil de decaimiento espectral}

La Figura~\ref{fig:fingerprint} refuerza esta observación desde una perspectiva
complementaria. Muestra las curvas de decaimiento de los valores singulares (media $\pm$
desviación estándar sobre las 12 capas) para cada tipo de componente.

\begin{figure}[H]
    \centering
    \includegraphics[width=0.85\textwidth]{figures/cap4/cap4_spectral_fingerprint_es.png}
    \caption{Curvas de decaimiento espectral por tipo de componente (media $\pm$ std
    sobre 12 capas). Q y K presentan decaimiento rápido; FFN Intermediate y Output
    mantienen valores singulares elevados a lo largo de todo el espectro.}
    \label{fig:fingerprint}
\end{figure}

Las matrices Q y K exhiben un decaimiento pronunciado: los primeros componentes principales
capturan la mayor parte de la varianza y los restantes contribuyen marginalmente. Las
matrices FFN, en cambio, presentan un espectro esencialmente plano, donde cada valor
singular adicional aporta una fracción similar de información. Esta diferencia cualitativa
en la forma del espectro explica por qué un mismo rango de truncamiento tiene efectos
radicalmente distintos sobre ambos tipos de componente.

\subsubsection{Solapamiento de subespacios Q-K-V}

Como observación complementaria, analizamos el solapamiento entre los subespacios
dominantes de las matrices Q, K y V dentro de cada capa, medido como la similitud de
coseno media entre los vectores singulares principales.

\begin{table}[H]
    \centering
    \caption{Rango de solapamiento entre subespacios (top-32 y top-64 vectores singulares).}
    \label{tab:overlap}
    \small
    \begin{tabular}{@{}lcc@{}}
        \toprule
        \textbf{Par} & \textbf{Top-32} & \textbf{Top-64} \\
        \midrule
        Q--K & 0{,}34 -- 0{,}63 & 0{,}31 -- 0{,}47 \\
        Q--V & 0{,}03 -- 0{,}18 & 0{,}06 -- 0{,}25 \\
        K--V & 0{,}05 -- 0{,}21 & 0{,}08 -- 0{,}27 \\
        \bottomrule
    \end{tabular}
\end{table}

Q y K comparten un subespacio significativo (solapamiento de hasta 0{,}63 para los 32
vectores singulares principales), lo cual es esperable dado que su producto $Q \cdot K^\top$
genera los pesos de atención. V, en cambio, opera en un subespacio sustancialmente
ortogonal a los de Q y K. Este resultado es consistente con la diferencia funcional de
ambos bloques: Q y K definen \emph{a qué atender}, mientras que V define \emph{qué
información extraer}. La Figura~\ref{fig:qkv-overlap} visualiza esta asimetría
capa a capa.

\begin{figure}[H]
    \centering
    \includegraphics[width=0.85\textwidth]{figures/cap4/cap4_qkv_overlap_es.png}
    \caption{Solapamiento de subespacios entre Q, K y V por capa, medido como la
    similitud de coseno media entre los \emph{top}-32 vectores singulares de cada par.
    El par Q--K (azul) presenta solapamiento sistemáticamente alto en todas las
    capas, coherente con que $QK^\top$ genere los pesos de atención. Los pares
    Q--V y K--V (naranja, verde) operan en subespacios prácticamente ortogonales.}
    \label{fig:qkv-overlap}
\end{figure}

No explotamos este solapamiento para compresión conjunta, pero lo
documentamos como hallazgo descriptivo de la geometría interna del modelo.

\subsubsection{Estructura espectral de las activaciones internas}
 
El análisis anterior se ha centrado en la geometría de las matrices de
peso. Una pregunta complementaria es cómo evoluciona la estructura
espectral de las \emph{activaciones} a lo largo de la red, es decir, qué
tan concentrada está la información que sale efectivamente de cada capa
durante la inferencia. La Figura~\ref{fig:cap4-internal-compression}
documenta tres métricas calculadas sobre las activaciones del modelo
\emph{fine-tuneado}, evaluadas sobre el conjunto de test.
 
\begin{figure}[H]
    \centering
    \includegraphics[width=0.95\textwidth]{figures/cap4/cap4_internal_compression_es.png}
    \caption{Caracterización espectral de las activaciones internas por
    capa. \textbf{Arriba, eje izquierdo (azul):} norma euclidiana media
    de los vectores de salida, separada en tokens de contenido y token
    \texttt{[CLS]}. \textbf{Arriba, eje derecho (rojo):} rango efectivo
    derivado de la entropía espectral de la matriz de activaciones, y
    umbral $k_{95}$ (número mínimo de valores singulares que captura el
    95\% de la energía). El sombreado de fondo distingue tres regímenes:
    léxico (Emb--L3), mezcla (L4--L7) y colapso (L8--L11).
    \textbf{Abajo:} heatmap de energía espectral acumulada por capa
    (filas) y posición del valor singular ordenado de mayor a menor
    (columnas). La línea de puntos blancos marca $k_{95}$. En L4--L7
    hacen falta $\sim 250$ valores singulares para capturar el 95\% de
    la energía; en L11, menos de 50.}
    \label{fig:cap4-internal-compression}
\end{figure}
 
La estructura es informativa para el problema de compresión. Las capas
intermedias (L4--L7) operan en una geometría de rango efectivo alto
(en torno a 200), lo que indica que la información se distribuye sobre
muchas direcciones independientes; la redundancia es alta y el espacio
está bien aprovechado. Las capas tardías (L8--L11), en cambio, presentan
un colapso espectral progresivo: el rango efectivo cae a aproximadamente
30 en L11, y el $k_{95}$ se reduce desde $\sim 250$ en L7 hasta menos de
50 al final del modelo. Esto significa que las capas tardías han
\emph{comprimido} ya su información en muy pocas direcciones.
 
Esta caracterización proporciona una explicación estructural del gradiente
de fragilidad que se documentará en \S\ref{sec:sensibilidad-componente}:
truncar a rango 128 una capa con rango efectivo 200 elimina información
que el modelo estaba usando, pero conserva la mayor parte de la señal;
truncar a rango 128 una capa con rango efectivo 30 es matemáticamente
inocuo (el modelo no usa más de 30 direcciones de todas formas), pero
la señal residual es tan concentrada que cualquier perturbación destruye
predicciones específicas. La capacidad de las capas tardías está
\emph{ya comprimida internamente}, y eso las hace estructuralmente
imposibles de comprimir más sin romperlas.


\subsubsection{Predicciones a partir del análisis espectral}

Los resultados espectrales generan tres predicciones concretas que las secciones
siguientes contrastarán experimentalmente:

\begin{enumerate}[nosep]
    \item Las matrices Q y K deberían tolerar una reducción de rango agresiva (hasta
    $r \approx 128$) con degradación mínima del rendimiento, dado que su dimensionalidad
    efectiva es muy inferior a 768.

    \item Las matrices FFN deberían ser el cuello de botella de cualquier estrategia de
    compresión, al requerir rangos sustancialmente mayores para retener una fracción
    equivalente de energía.

    \item Una estrategia de compresión uniforme, que aplique el mismo rango a todos los
    componentes, será necesariamente subóptima, al ignorar esta heterogeneidad espectral.
\end{enumerate}

Las secciones \S\ref{sec:compresion-uniforme} y \S\ref{sec:sensibilidad-componente}
confirman las dos primeras predicciones, y la \S\ref{sec:compresion-adaptativa} valida
la tercera.


% =============================================================
% §4.2 — Compresión uniforme
% =============================================================

\subsection{Compresión uniforme y transición de fase}
\label{sec:compresion-uniforme}

La estrategia más simple de compresión SVD consiste en aplicar el mismo rango de
truncamiento $r$ a todas las matrices lineales del modelo simultáneamente. Esta sección
evalúa cinco niveles de rango uniforme ($r \in \{512, 384, 256, 128, 64\}$) y documenta
una transición de fase abrupta en el rendimiento.

\subsubsection{Degradación global}

La Tabla~\ref{tab:uniform} muestra el F1 macro, el ratio de parámetros y
la retención de rendimiento para cada nivel de rango.

\begin{table}[H]
    \centering
    \caption{Resultados de compresión uniforme sobre el conjunto de test.
    El \emph{ratio de parámetros} se define como
    $\text{params}_{\text{comprimido}}/\text{params}_{\text{original}}$:
    valores $<1$ indican compresión real; $=1$, sin reducción; $>1$,
    expansión paramétrica.}
    \label{tab:uniform}
    \small
    \begin{tabular}{@{}lccc@{}}
        \toprule
        \textbf{Estrategia} & \textbf{F1 macro} & \textbf{Ratio params.} & \textbf{Retención} \\
        \midrule
        Baseline           & 0{,}577 & 1{,}00 & 100{,}0\% \\
        Uniform $r$=512    & 0{,}464 & 1{,}00 & 80{,}5\%  \\
        Uniform $r$=384    & 0{,}251 & 0{,}81 & 43{,}5\%  \\
        Uniform $r$=256    & 0{,}025 & 0{,}61 & 4{,}3\%   \\
        Uniform $r$=128    & 0{,}000 & 0{,}42 & 0{,}0\%   \\
        Uniform $r$=64     & 0{,}000 & 0{,}32 & 0{,}0\%   \\
        \bottomrule
    \end{tabular}
\end{table}

El patrón dominante es una \emph{transición de fase} entre $r=384$ y $r=256$: el modelo
pasa de retener el 43\% del rendimiento a retener apenas el 4\%, una caída de casi 40
puntos porcentuales en un solo salto de rango. A partir de $r=128$, el modelo entra en
``muerte clínica'' con F1 exactamente igual a 0{,}0: no es capaz de clasificar
correctamente ni una sola emoción. La Figura~\ref{fig:uniform-transition} ilustra
visualmente este \emph{cliff}.

\begin{figure}[H]
    \centering
    \includegraphics[width=0.78\textwidth]{figures/cap4/cap4_uniform_transition_es.png}
    \caption{Transición de fase bajo compresión uniforme. F1 macro
    en función del rango: la curva se mantiene en niveles funcionales
    para $r \geq 384$ y colapsa abruptamente entre $r=384$ y $r=256$. La
    región sombreada marca la zona de muerte clínica ($F1 = 0$).}
    \label{fig:uniform-transition}
\end{figure}

La zona útil de compresión uniforme se limita a $r=512$ (degradación significativa pero
funcional) y $r=384$ (degradación severa pero no total). Cualquier rango inferior destruye
la capacidad del modelo de forma irrecuperable.

\subsubsection{Vulnerabilidad emocional diferencial}

La degradación no afecta a todas las emociones por igual. La Figura~\ref{fig:emotion-vulnerability}
muestra el F1 por emoción para cada nivel de compresión uniforme, revelando un espectro
amplio de resistencia.

\begin{figure}[H]
    \centering
    \includegraphics[width=0.75\textwidth]{figures/cap4/cap4_emotion_vulnerability_es.png}
    \caption{F1 por emoción para cada nivel de compresión uniforme. Las emociones están
    ordenadas por su F1 en el baseline. Las emociones frecuentes y léxicamente simples
    (arriba) sobreviven más que las infrecuentes y semánticamente complejas (abajo).}
    \label{fig:emotion-vulnerability}
\end{figure}

Las emociones más resistentes comparten características comunes: son frecuentes en el
corpus y están asociadas a marcadores léxicos fuertes. \emph{Gratitude}, por ejemplo,
retiene un F1 de 0{,}92 a $r=384$ y aún mantiene 0{,}55 a $r=256$, apoyada por la presencia
recurrente de expresiones como ``thank you'' y ``thanks''. \emph{Amusement} (0{,}74 a
$r=384$) y \emph{love} (0{,}77 a $r=384$) siguen un patrón similar.

En el extremo opuesto, \emph{embarrassment} colapsa ya a $r=512$ (de 0{,}27 a 0{,}00),
\emph{realization} cae de 0{,}27 a 0{,}16 en el mismo nivel, y \emph{disappointment} pasa
de 0{,}28 a 0{,}06. Estas emociones son infrecuentes en el corpus de entrenamiento,
semánticamente ambiguas y carecen de marcadores léxicos distintivos. Su representación
depende de dimensiones espectrales de cola, precisamente las primeras en ser eliminadas
por la compresión SVD.

Este patrón conecta directamente con los resultados de cristalización del
Capítulo~\ref{sec:resultados-interpretabilidad}: las emociones que cristalizan temprano en
la profundidad del modelo (es decir, que son reconocibles desde las primeras capas) tienden
a ser las más resistentes a la compresión, mientras que las que requieren procesamiento
profundo son las más vulnerables. La \S\ref{sec:sintesis} formaliza esta correlación.


% =============================================================
% §4.3 — Sensibilidad por componente y profundidad
% =============================================================

\subsection{Sensibilidad por componente y profundidad}
\label{sec:sensibilidad-componente}

La sección anterior mostró que la compresión uniforme presenta una transición de fase
abrupta. El análisis espectral de la \S\ref{sec:analisis-espectral} predijo que los
distintos tipos de componente deberían tener sensibilidades radicalmente diferentes.
Esta sección lo verifica experimentalmente mediante un estudio sistemático de 33
configuraciones, aislando el efecto del tipo de componente (8 tipos $\times$ 3 rangos)
y de la profundidad (3 grupos $\times$ 3 rangos).

\subsubsection{Sensibilidad por tipo de componente}

Para aislar el efecto de cada tipo de componente, comprimimos selectivamente solo las 12
matrices de un tipo dado (por ejemplo, las 12 matrices Query, una por capa) a un rango
fijo, manteniendo el resto del modelo intacto. La Tabla~\ref{tab:component-sensitivity}
presenta los resultados al rango $r=128$, que maximiza el contraste entre componentes.

\begin{table}[H]
    \centering
    \caption{Sensibilidad por componente a $r=128$. La sensibilidad normalizada mide
    los puntos porcentuales de F1 perdidos por cada 1\% de parámetros eliminados.}
    \label{tab:component-sensitivity}
    \small
    \begin{tabular}{@{}lcccc@{}}
        \toprule
        \textbf{Componente} & \textbf{F1 macro} & \textbf{Retención} & \textbf{Sensibilidad} \\
        \midrule
        Query (Q)           & 0{,}574 & 99{,}4\% & 0{,}14 \\
        Key (K)             & 0{,}568 & 98{,}4\% & 0{,}38 \\
        Value (V)           & 0{,}395 & 68{,}5\% & 7{,}31 \\
        Attn Output         & 0{,}325 & 56{,}3\% & 10{,}13 \\
        FFN Output          & 0{,}128 & 22{,}2\% & 3{,}80 \\
        FFN Intermediate    & 0{,}040 & 6{,}9\%  & 4{,}55 \\
        Attention (todas)   & 0{,}027 & 4{,}6\%  & --- \\
        FFN (todas)         & 0{,}001 & 0{,}1\%  & --- \\
        \bottomrule
    \end{tabular}
\end{table}

El resultado confirma y amplifica la predicción espectral. A rango 128, Query retiene el
99,4\% del rendimiento mientras que FFN Intermediate retiene apenas el 6,9\%: un factor
de 14$\times$ en retención absoluta. Normalizado por la fracción de parámetros
eliminados, la diferencia en sensibilidad es de 72$\times$ (0,14 frente a 10,13
para Attn Output, el componente de atención más sensible).

Este resultado valida la primera predicción del análisis espectral: Q y K son
efectivamente ``inmunes'' a la compresión agresiva, porque su información reside en un
subespacio de dimensionalidad mucho menor que el rango completo de la matriz.

\subsubsection{Tres regímenes de compresibilidad}

Al examinar la evolución de la retención a lo largo de los tres rangos evaluados
($r \in \{256, 128, 64\}$), los seis tipos de componente se agrupan en tres regímenes
cualitativamente distintos:

\begin{description}[nosep, leftmargin=1.2em]
    \item[Régimen~1 --- Inmunes (Q, K):] Degradación gradual y casi lineal. Query
    pierde apenas 1{,}8 puntos porcentuales entre $r=256$ y $r=64$. Estas matrices
    operan en un subespacio de dimensionalidad efectiva muy inferior a 768.

    \item[Régimen~2 --- Efecto cliff (V, Attn Output):] Sobreviven a $r=256$
    (retención $\sim$93\%) pero colapsan a $r=64$ ($\sim$3--5\%). Existe un umbral
    crítico entre $r=128$ y $r=64$ por debajo del cual la información no se puede
    recuperar. Value, por ejemplo, cae 63 puntos porcentuales solo entre $r=128$ y
    $r=64$.

    \item[Régimen~3 --- Frágiles (FFN Inter, FFN Output):] Ya gravemente dañados a
    $r=256$ (52\% y 75\% de retención respectivamente). El cliff principal se sitúa
    entre $r=256$ y $r=128$. Son matrices de alta dimensionalidad efectiva cuya
    información está distribuida a lo largo de todo el espectro de valores singulares.
\end{description}

\begin{figure}[H]
    \centering
    \includegraphics[width=0.85\textwidth]{figures/cap4/cap4_component_sensitivity_es.png}
    \caption{Curvas de retención de F1 en función del rango para cada uno
    de los seis tipos de componente. Los tres regímenes
    descritos en el texto se separan visualmente: Q y K (azules) decaen
    suavemente; V y Attn Output (naranjas) presentan el \emph{cliff}
    entre $r=128$ y $r=64$; las dos FFN (rojas) caen abruptamente desde
    $r=256$. La línea horizontal punteada marca el 50\% de retención.}
    \label{fig:component-sensitivity}
\end{figure}

Estos tres regímenes corresponden directamente a las tres formas de espectro observadas
en la \S\ref{sec:analisis-espectral}: decaimiento rápido (Q, K), decaimiento intermedio
(V, Attn Output) y espectro plano (FFN). La estructura espectral predice con precisión
el comportamiento bajo compresión.

\subsubsection{Sensibilidad por profundidad}

Para evaluar el efecto de la profundidad, comprimimos simultáneamente todas las matrices
de un grupo de capas (Early: 0--3, Middle: 4--7, Late: 8--11) manteniendo el resto del
modelo intacto.

\begin{table}[H]
    \centering
    \caption{Retención de F1 por grupo de profundidad y rango de compresión.}
    \label{tab:depth-sensitivity}
    \small
    \begin{tabular}{@{}lccc@{}}
        \toprule
        \textbf{Grupo} & $\boldsymbol{r=256}$ & $\boldsymbol{r=128}$ & $\boldsymbol{r=64}$ \\
        \midrule
        Early (0--3)   & 85{,}4\% & 62{,}5\% & 37{,}0\% \\
        Middle (4--7)  & 82{,}5\% & 46{,}5\% & 25{,}3\% \\
        Late (8--11)   & 33{,}3\% &  0{,}0\% &  0{,}0\% \\
        \bottomrule
    \end{tabular}
\end{table}

Las capas tardías (8--11) son \emph{catastróficamente sensibles}: a $r=128$, el modelo
muere por completo (F1 = 0{,}0). En contraste, las capas tempranas retienen más del 60\%
del rendimiento al mismo rango. La Figura~\ref{fig:depth-sensitivity} sintetiza este
gradiente.

\begin{figure}[H]
    \centering
    \includegraphics[width=0.85\textwidth]{figures/cap4/cap4_depth_sensitivity_es.png}
    \caption{Sensibilidad por grupo de profundidad. Cada barra representa
    la retención de F1 macro al comprimir todas las matrices de un grupo
    de capas a un rango dado, manteniendo el resto intacto. Las capas
    tardías (rojo) son las únicas que colapsan por completo a $r=128$;
    las tempranas (verde) y medias (azul) retienen capacidad funcional.}
    \label{fig:depth-sensitivity}
\end{figure}

Este resultado es consistente con la literatura de
interpretabilidad mecánica, que identifica las capas finales de los Transformers como las
responsables de la representación semántica de alto nivel. En el contexto de la clasificación
emocional, las capas 8--11 son las que codifican las distinciones finas entre emociones;
destruir su estructura equivale a destruir la capacidad clasificadora.

Este hallazgo complementa el análisis por componente: la compresión uniforme falla no solo
porque trata igual a Q y a FFN, sino porque trata igual a las capas 0 y a las capas 11.
Una estrategia de compresión verdaderamente informada debería respetar ambos ejes de
asimetría: tipo de componente y profundidad.

\subsubsection{Nota metodológica: silhouette scores en contextos de colapso}
\label{sec:silhouette-caveat}

Como métrica complementaria al F1, evaluamos los silhouette scores sobre
proyecciones t-SNE de los embeddings del modelo comprimido. Los resultados
revelan una limitación importante de esta métrica en contextos de
compresión severa que conviene documentar como aporte metodológico
transferible a otros estudios de compresión.

Los silhouette scores negativos observados para FFN Intermediate a $r=64$
($-0{,}125$) y $r=128$ ($-0{,}051$) indican que la compresión no solo
destruye la estructura de clusters sino que la \emph{desordena}
activamente: las muestras quedan más cerca de clusters ajenos que del
propio.

Sin embargo, las capas tardías a $r=128$ presentan un silhouette score de
$0{,}15$, el más alto de todas las configuraciones, a pesar de un F1
de $0{,}0$. Esto constituye un \emph{artefacto sistemático}: cuando el
modelo colapsa, todos los embeddings convergen a un único cluster trivial,
y la métrica de silhouette, que premia la cohesión intra-cluster y la
separación inter-cluster, interpreta erróneamente esa convergencia
patológica como ``buena estructura''. Documentamos esta observación como
una limitación general de las métricas geométricas de clustering en
contextos de colapso del modelo, y recomendamos que cualquier estudio de
compresión que utilice silhouette o métricas relacionadas (Davies-Bouldin,
Calinski-Harabasz) verifique siempre la validez de la métrica mediante
inspección directa de las proyecciones cuando el F1 sea bajo. La regla
operativa es simple: si el modelo colapsa por F1, las métricas de
clustering geométrico no son interpretables.


% =============================================================
% §4.4 — Compresión adaptativa
% =============================================================

\subsection{Compresión adaptativa y recetas mixtas}
\label{sec:compresion-adaptativa}

Los resultados de las secciones anteriores establecen que la compresión uniforme es
inherentemente subóptima: aplica el mismo rango a componentes con sensibilidades que
difieren en un factor de hasta 72$\times$ (normalizado por parámetros eliminados) y
a capas con vulnerabilidades que van del 0\% al 62\% de retención al mismo rango. La compresión adaptativa propone una alternativa:
asignar a cada matriz un rango proporcional a su complejidad espectral, determinado por
el umbral de energía $\tau$.

\subsubsection{Resultados de compresión adaptativa}

En la compresión adaptativa, cada matriz recibe el rango mínimo $k_\tau$ necesario para
retener una fracción $\tau$ de su energía espectral. Matrices con espectro concentrado
(Q, K) reciben rangos bajos; matrices con espectro plano (FFN) reciben rangos altos. La
Tabla~\ref{tab:adaptive} presenta los resultados para cinco umbrales de energía.

\begin{table}[H]
    \centering
    \caption{Resultados de compresión adaptativa y recetas mixtas.
    \emph{Ratio params.} =
    $\text{params}_{\text{comprimido}}/\text{params}_{\text{original}}$,
    donde $<1$ es compresión real y $>1$ expansión paramétrica
    (ver \S\ref{sec:caveat-compresion}).}
    \label{tab:adaptive}
    \small
    \begin{tabular}{@{}lccc@{}}
        \toprule
        \textbf{Estrategia} & \textbf{F1 macro} & \textbf{Ratio params.} & \textbf{Retención} \\
        \midrule
        Baseline             & 0{,}577 & 1{,}00 & 100{,}0\% \\
        Adaptive 99\%        & 0{,}558 & 1{,}20 & 96{,}7\%  \\
        Mixed Conservative   & 0{,}505 & 1{,}04 & 87{,}5\%  \\
        Adaptive 95\%        & 0{,}491 & 1{,}02 & 85{,}1\%  \\
        Mixed Aggressive     & 0{,}382 & 0{,}91 & 66{,}3\%  \\
        Adaptive 90\%        & 0{,}357 & 0{,}89 & 61{,}8\%  \\
        Adaptive 85\%        & 0{,}260 & 0{,}80 & 45{,}1\%  \\
        Adaptive 80\%        & 0{,}143 & 0{,}72 & 24{,}7\%  \\
        \bottomrule
    \end{tabular}
\end{table}

La compresión adaptativa al 99\% de energía retiene el 96.7\% del rendimiento, con
una degradación de apenas 1{,}9 puntos porcentuales de F1 respecto al baseline. Al 95\%,
la retención es del 85.1\%, lo cual constituye un punto de equilibrio razonable.

Las recetas mixtas combinan compresión adaptativa con reglas específicas por tipo de
componente (por ejemplo, comprimir más agresivamente Q y K mientras se protegen las FFN).
La receta conservadora (87.5\% de retención) se sitúa entre Adaptive 95\% y 99\%,
ofreciendo un compromiso intermedio.

\subsubsection{Frontera de Pareto: adaptativa vs.\ uniforme}

La Figura~\ref{fig:pareto} compara las dos familias de estrategias en el espacio
retención--compresión. La compresión adaptativa \emph{domina} a la uniforme en toda la
frontera de Pareto: para cualquier nivel de retención deseado, la adaptativa lo consigue
con un ratio de compresión igual o superior.

\begin{figure}[H]
    \centering
    \includegraphics[width=0.85\textwidth]{figures/cap4/cap4_pareto_frontier_es.png}
    \caption{Frontera de Pareto: compresión uniforme (cuadrados) vs.\ adaptativa
    (círculos). La adaptativa domina en todo el rango. Los puntos Mixed se sitúan
    sobre la frontera adaptativa.}
    \label{fig:pareto}
\end{figure}

La ventaja de la compresión adaptativa es más pronunciada en la zona de alta retención
(>80\%), donde la uniforme solo dispone de un punto útil ($r=512$) mientras que la
adaptativa ofrece un continuo de opciones entre 85\% y 97\% de retención.

\subsubsection{Distribución de rangos asignados}

La Figura~\ref{fig:rank-distribution} ilustra cómo la compresión adaptativa distribuye
los rangos entre los distintos tipos de componente para cada umbral de energía.

\begin{figure}[H]
    \centering
    \includegraphics[width=0.85\textwidth]{figures/cap4/cap4_adaptive_rank_distribution_es.png}
    \caption{Distribución de rangos asignados por la compresión adaptativa para cada
    umbral de energía, desglosada por tipo de componente. Q y K reciben rangos
    consistentemente menores que FFN.}
    \label{fig:rank-distribution}
\end{figure}

La distribución confirma la lógica espectral: a un umbral del 95\%, Q y K reciben rangos
del orden de 390--400, mientras que FFN Intermediate y Output reciben rangos de 610--625.
La adaptativa simplemente \emph{respeta} la estructura natural de las matrices en lugar
de forzar un rango uniforme.

\subsubsection{Caveat: ratio de compresión y simplificación de rango}
\label{sec:caveat-compresion}

Es necesario documentar una limitación importante de la compresión adaptativa a umbrales
altos. La factorización SVD truncada a rango $r$ reemplaza una matriz $W \in \mathbb{R}^{m
\times n}$ por dos matrices $U_r \in \mathbb{R}^{m \times r}$ y $V_r \in \mathbb{R}^{r
\times n}$ (absorbiendo $\Sigma_r$ en una de las dos). El número de parámetros pasa de
$mn$ a $r(m + n)$. Para que exista compresión real, se requiere $r < \frac{mn}{m+n}$.

En el caso de BERT-base, las matrices de self-attention tienen dimensiones $768 \times 768$,
lo que da un umbral de $r < 384$ para compresión real. Las matrices FFN tienen dimensiones
$768 \times 3072$ (Intermediate) y $3072 \times 768$ (Output), con un umbral de
$r < \frac{768 \times 3072}{768 + 3072} = 614.4$.

Las estrategias Adaptive 95\% y 99\% asignan, por tanto, rangos que para muchas
matrices \emph{superan} el umbral de compresión real. Los ratios de parámetros de
1{,}02 (Adaptive 95\%) y 1{,}20 (Adaptive 99\%) confirman que estas estrategias
producen un modelo con \emph{más} parámetros que el original.

En rigor, estas estrategias no constituyen compresión de parámetros sino
\emph{simplificación de rango}: la representación factorizada tiene menor rango
matricial que la original, lo cual puede ser útil para análisis e interpretabilidad, pero
no reduce el tamaño del modelo. La compresión real solo comienza a partir de Adaptive
90\% (ratio 0{,}89), y los niveles con compresión significativa (85\% y 80\%)
ya implican degradaciones del 55\% y 75\% respectivamente.


% =============================================================
% Cierre del capítulo
% =============================================================

\subsection{Síntesis del capítulo}
\label{sec:sintesis-compresion}

Los resultados de este capítulo establecen cuatro hallazgos principales:

\begin{enumerate}[nosep]
    \item \textbf{Asimetría espectral.} Las matrices de self-attention (Q, K) concentran
    su información en un subespacio de dimensionalidad efectiva $\sim$400, mientras que
    las matrices FFN requieren $\sim$620 dimensiones. Este gap de $\sim$220 rangos es
    una propiedad estructural del modelo \emph{fine-tuned}, consistente en las 12 capas.

    \item \textbf{Transición de fase en compresión uniforme.} La compresión uniforme
    presenta un cliff abrupto entre $r=384$ (43\% de retención) y $r=256$ (4\%).
    La zona útil se limita a $r \geq 384$, y cualquier rango inferior destruye el
    modelo de forma irrecuperable.

    \item \textbf{Sensibilidad heterogénea.} Los componentes presentan sensibilidades
    que difieren en un factor de 72$\times$ (normalizado por parámetros). Q y K son
    inmunes a compresión agresiva; V y Attn Output presentan un cliff entre $r=128$ y
    $r=64$; las FFN son frágiles desde $r=256$. Las capas tardías (8--11) son
    catastróficamente sensibles, con muerte total a $r=128$.

    \item \textbf{Dominancia matizada de la compresión adaptativa.} La
    compresión adaptativa supera a la uniforme \emph{en el rango donde hay
    compresión real} (umbrales del 80\% y 90\%, con ratios de parámetros
    de 0{,}72 y 0{,}89 respectivamente). Por encima del 90\%, los umbrales
    del 95\% y 99\% producen ratios de parámetros $>1$ ($1{,}02$ y
    $1{,}20$ respectivamente): no son técnicamente compresión, sino
    \emph{simplificación de rango con expansión paramétrica}. La aparente
    dominancia de la adaptativa en toda la frontera de Pareto debe
    interpretarse, por tanto, con esta cautela: cuando el ratio de
    parámetros es genuinamente menor que el original ($<1$), la adaptativa
    supera a la uniforme; en la región $\geq 1$, ambas familias se
    encuentran fuera del régimen de compresión efectiva.
\end{enumerate}

Estos hallazgos motivan dos líneas de investigación que se desarrollan en los capítulos
siguientes: la interpretabilidad mecánica (Capítulo~\ref{sec:resultados-interpretabilidad}),
que busca explicar \emph{por qué} ciertos componentes y capas son más críticos, y la
síntesis (\S\ref{sec:sintesis}), que conecta la vulnerabilidad a compresión con la
profundidad de cristalización emocional.
% =============================================================
% CAPÍTULO 5 — Resultados de interpretabilidad mecánica
% Versión: matrícula de honor
% Alimentado por: NB04 (probing), NB05 (activation patching),
%                 NB06 (head ablation), NB07 (neuron analysis),
%                 NB08 (logit lens)
% =============================================================

\section{Resultados de interpretabilidad mecánica}
\label{sec:resultados-interpretabilidad}

El capítulo anterior documentó qué ocurre cuando se comprime BERT:
una transición de fase abrupta entre rango 384 y 256, una asimetría
radical entre componentes (factor 14$\times$ a rango~128 y hasta
72$\times$ a rango~64, normalizado por parámetros eliminados) y un
gradiente de fragilidad creciente con la profundidad. Este capítulo
aborda la pregunta complementaria: por qué ocurre. Para ello se
aplican cinco técnicas de interpretabilidad mecánica que examinan la
estructura interna del modelo desde ángulos diferentes pero convergentes.

La primera técnica, \emph{probing} (\S\ref{sec:probing}), entrena
clasificadores lineales sobre las representaciones intermedias para
determinar en qué capa se forma la señal emocional. La segunda,
\emph{logit lens} (\S\ref{sec:logit-lens}), aplica la cabeza de
clasificación entrenada del modelo a las representaciones de capas
intermedias para distinguir información disponible de información
alineada con la base del clasificador. La tercera, \emph{activation
patching} (\S\ref{sec:activation-patching}), restaura selectivamente
componentes del modelo original dentro de una versión comprimida para
localizar los loci funcionales críticos. La cuarta, \emph{ablación de
cabezas de atención} (\S\ref{sec:head-ablation}), anula individualmente
cada una de las 144 cabezas para construir una taxonomía funcional del
mecanismo atencional. La quinta, \emph{análisis neuronal de las FFN}
(\S\ref{sec:neuron-analysis}), examina la especialización de las 36\,864
neuronas de las redes feed-forward y su relación con la vulnerabilidad a
la compresión.

Cada sección presenta sus resultados de forma autónoma, pero el lector
observará un patrón transversal: las cinco técnicas convergen en una
arquitectura funcional de tres niveles, donde las capas tempranas extraen
señal léxica bruta, las capas intermedias realizan cómputo en una
geometría interna ilegible para el clasificador, y las capas profundas
(8--11) proyectan la representación sobre la base del clasificador. La
síntesis formal de esta convergencia se desarrolla en
\S\ref{sec:interp-sintesis}.


% =============================================================
% 5.1 — PROBING: CRISTALIZACIÓN EMOCIONAL POR CAPAS
% =============================================================

\subsection{Probing: cristalización emocional por capas}
\label{sec:probing}

\subsubsection{Diseño experimental}

Para determinar en qué punto de la jerarquía de capas de BERT emerge la
información emocional, se entrenan clasificadores \emph{probe} sobre las
representaciones intermedias del modelo. Se utiliza regresión logística
con regularización por defecto ($C = 1{,}0$) como clasificador, siguiendo
la metodología de Hewitt y Liang~\cite{hewitt2019designing}: el objetivo
del probe es medir la calidad de la representación, no la capacidad del
clasificador, por lo que se evita deliberadamente el tuning de
hiperparámetros. Un probe no-lineal complejo podría extraer información
de cualquier representación; un probe lineal solo detecta información ya
separada linealmente en el espacio de activaciones.

La representación de entrada es el token \texttt{[CLS]} extraído de cada
uno de los 13 \textit{hidden states} del modelo (la capa de embedding más
las 12 capas del encoder). Se entrena un probe binario independiente para
cada combinación de emoción y capa con una capa de información
significativa, resultando en 276 clasificadores activos (las 12 capas del
encoder por las 23 emociones; los probes sobre embedding producen
F1 = 0 uniformemente y se omiten del análisis efectivo).

Se define la \emph{capa de cristalización} de una emoción como la primera
capa en la que el F1 del probe alcanza el 80\% de su valor máximo a lo
largo de todas las capas. Este umbral es convencional; análisis de
sensibilidad con umbrales de 0{,}7 y 0{,}9 confirman que el orden
relativo de cristalización entre emociones permanece estable.

\subsubsection{Tres regímenes de cristalización}

Las 23 emociones se distribuyen en tres grupos claramente diferenciados
según su profundidad de cristalización
(Tabla~\ref{tab:cristalizacion-grupos},
Figura~\ref{fig:probing-heatmap}).

\begin{table}[ht]
\centering
\caption{Grupos de cristalización emocional. La capa indica la primera
capa donde el F1 del probe alcanza el 80\% de su máximo. La columna
``F1 en cap.'' muestra el F1 del probe en la capa de cristalización.}
\label{tab:cristalizacion-grupos}
\small
\begin{tabular}{llcl}
\toprule
\textbf{Grupo} & \textbf{Capas} & \textbf{N} & \textbf{Emociones representativas} \\
\midrule
Temprano & L0--L3 & 9 & gratitude (L0, 0{,}922), amusement (L0, 0{,}827), \\
         &        &   & love (L0, 0{,}818), curiosity (L0, 0{,}674), \\
         &        &   & optimism (L0, 0{,}569), admiration (L2, 0{,}735), \\
         &        &   & disgust (L2, 0{,}532), fear (L3, 0{,}708), \\
         &        &   & sadness (L3, 0{,}621) \\
\addlinespace
Medio    & L4--L7 & 8 & anger (L4, 0{,}543), remorse (L5, 0{,}655), \\
         &        &   & embarrassment (L5, 0{,}406), \\
         &        &   & disapproval (L6, 0{,}521), caring (L6, 0{,}502), \\
         &        &   & surprise (L7, 0{,}558), confusion (L7, 0{,}502), \\
         &        &   & realization (L7, 0{,}327) \\
\addlinespace
Tardío   & L8--L11 & 6 & joy (L8, 0{,}601), desire (L8, 0{,}582), \\
         &         &   & approval (L8, 0{,}547), excitement (L8, 0{,}472), \\
         &         &   & annoyance (L8, 0{,}416), \\
         &         &   & disappointment (L9, 0{,}361) \\
\bottomrule
\end{tabular}
\end{table}

\begin{figure}[H]
\centering
\includegraphics[width=0.85\textwidth]{figures/cap5/cap5_probing_heatmap_es.png}
\caption{Heatmap de F1 de los probes (23 emociones $\times$ 13 capas). El
marcador indica la capa de cristalización de cada emoción. Las emociones
se ordenan por profundidad de cristalización creciente.}
\label{fig:probing-heatmap}
\end{figure}

La interpretación de estos grupos conecta la profundidad de procesamiento
con la naturaleza lingüística de la señal emocional. Las emociones del
grupo temprano poseen vocabulario altamente distintivo:
\emph{gratitude} se asocia a marcadores léxicos explícitos (``thanks'',
``thank you''), \emph{amusement} a indicadores como ``lol'' o ``haha'',
\emph{love} a un campo semántico bien delimitado. En estos casos, la
primera capa del encoder (L0) ya dispone de suficiente información
léxica para que un clasificador lineal alcance un rendimiento cercano
al máximo.

En el extremo opuesto, las emociones tardías requieren integración
contextual para desambiguarse de categorías vecinas:
\emph{annoyance} debe distinguirse de \emph{anger},
\emph{approval} de \emph{admiration}, \emph{disappointment} de
\emph{sadness}. Estas distinciones no están disponibles en el
vocabulario de superficie y emergen solo tras múltiples capas de
procesamiento atencional.

Las emociones medias ocupan posiciones intermedias y a menudo combinan
señales léxicas con requisitos contextuales. \emph{Embarrassment} es un
caso ilustrativo: tiene cierta señal léxica (``embarrassed'',
``awkward''), pero su uso pragmático (a veces autoirónico, a veces
genuino) demanda contexto, lo que retrasa su cristalización a L5.


La Figura~\ref{fig:cap5-galaxia-evolution} visualiza esta cristalización
progresiva proyectando las representaciones del token \texttt{[CLS]} en
cuatro capas representativas mediante t-SNE.
 
\begin{figure}[H]
    \centering
    \includegraphics[width=\textwidth]{figures/cap5/cap5_galaxia_evolution_es.png}
    \caption{Proyecciones t-SNE de las representaciones del token
    \texttt{[CLS]} en cuatro capas del modelo. Los puntos pequeños son
    ejemplos individuales del conjunto de test; los rombos grandes
    representan los centroides de cada emoción. Los colores agrupan las
    23 emociones en seis clusters semánticos derivados del análisis
    neuronal (\S\ref{sec:neuron-analysis}). El embedding inicial
    apenas diferencia las emociones (los centroides colapsan en una
    región pequeña). En L4 los puntos se distribuyen, pero los
    centroides aún se solapan. En L8 emergen estructuras parciales:
    las positivas de alta energía se separan, las negativas internas
    comienzan a agruparse. En L11 los seis clusters están claramente
    diferenciados.}
    \label{fig:cap5-galaxia-evolution}
\end{figure}
 
La progresión es consistente con los grupos de cristalización que la
Tabla~\ref{tab:cristalizacion-grupos} identifica numéricamente. En el
embedding ninguna emoción es separable, lo que es esperable: el embedding
codifica el significado del token \texttt{[CLS]} antes de procesar la
frase. En L4 se observa el comportamiento medio del cluster ``Positivas
alta energía'' (gratitude, amusement, love), cuyos centroides ya se
desplazan hacia regiones distintivas, anticipando que cristalizan en
L0--L2. En L8 emergen las distinciones contextuales: el cluster de
``Negativas internas'' se concentra en la región superior, el de
``Negativas reactivas'' en la derecha. En L11 los seis clusters ocupan
regiones diferenciadas del plano, lo que cierra la formación
representacional documentada por probing.

\subsubsection{Ganancia informativa por capa}

La contribución incremental de cada capa se cuantifica como la
diferencia en F1 macro medio entre capas consecutivas
(Figura~\ref{fig:probing-info-gain}). Este análisis revela una
arquitectura de procesamiento en dos fases con un detalle
metodológicamente relevante.

\begin{figure}[H]
\centering
\includegraphics[width=0.8\textwidth]{figures/cap5/cap5_information_gain_es.png}
\caption{Ganancia informativa incremental por capa, medida como el
aumento en F1 macro medio respecto a la capa anterior. La capa L0
absorbe el 61\% de la separabilidad final del modelo.}
\label{fig:probing-info-gain}
\end{figure}

La capa de embedding no contiene información emocional (F1 medio = 0
para todas las emociones, valor exacto). La primera capa del encoder
(L0) produce un salto de $+0{,}349$, capturando aproximadamente el 61\%
de la separabilidad lineal final del modelo
($F1_{L0}/F1_{L11} = 0{,}349/0{,}569 \approx 0{,}61$). Este salto
corresponde a la extracción de señal léxica bruta. Las capas L1--L7
contribuyen con ganancias modestas pero positivas
($+0{,}010$ a $+0{,}048$), refinando gradualmente las representaciones.
La capa L8 produce un segundo pico de $+0{,}033$, asociado a la
desambiguación contextual de las emociones tardías. Las capas finales
(L10--L11) aportan ganancias marginales: L10 incluso muestra una
ganancia ligeramente negativa ($-0{,}002$), seguida de un pequeño
$+0{,}004$ en L11.

Esta estructura de dos fases, con captura léxica masiva en L0,
refinamiento contextual con pico en L8 y consolidación marginal en
L10--L11, anticipa un resultado que será central en el resto del
capítulo: la distinción entre crear información emocional y
formatearla para el clasificador. Las capas L10--L11 apenas
crean información nueva según el probing, pero \S\ref{sec:logit-lens}
y \S\ref{sec:activation-patching} demostrarán que son las únicas
capaces de hacer utilizable la información existente para la
cabeza de clasificación. Esta distinción se cuantifica directamente en
la siguiente sección.

\subsubsection{La frecuencia no predice la profundidad}

Un análisis de correlaciones (Tabla~\ref{tab:probing-correlaciones},
Figura~\ref{fig:probing-correlaciones}) descarta la hipótesis de que la
profundidad de cristalización esté determinada por el volumen de datos
de entrenamiento, hipótesis intuitiva pero empíricamente no soportada.

\begin{table}[ht]
\centering
\caption{Correlaciones de Spearman entre propiedades de las emociones
($n = 23$). Solo la correlación entre profundidad de cristalización y
F1 máximo es significativa.}
\label{tab:probing-correlaciones}
\begin{tabular}{lccc}
\toprule
\textbf{Par de variables} & $\rho$ & \textbf{p-valor} & \textbf{Sig.} \\
\midrule
Capa cristalización vs.\ F1 máximo  & $-0{,}673$ & $0{,}0004$ & *** \\
Frecuencia vs.\ capa cristalización & $-0{,}265$ & $0{,}22$   & ns. \\
Frecuencia vs.\ F1 máximo           & $+0{,}33$\phantom{0}  & $0{,}12$   & ns. \\
\bottomrule
\end{tabular}
\end{table}

\begin{figure}[H]
\centering
\includegraphics[width=\textwidth]{figures/cap5/cap5_crystal_correlations_es.png}
\caption{Izquierda: capa de cristalización vs.\ F1 máximo
($\rho = -0{,}673$, $p < 0{,}001$). Derecha: frecuencia en el dataset
vs.\ capa de cristalización ($\rho = -0{,}265$, $p = 0{,}22$).}
\label{fig:probing-correlaciones}
\end{figure}

La correlación entre profundidad de cristalización y F1 máximo es fuerte
y altamente significativa ($\rho = -0{,}673$, $p = 0{,}0004$): las
emociones que cristalizan antes alcanzan mayor rendimiento. Sin embargo,
la frecuencia en el dataset no predice ni la profundidad de
cristalización ($p = 0{,}22$) ni el rendimiento máximo ($p = 0{,}12$).

Un caso ilustrativo: \emph{disgust}, con 793~ejemplos de entrenamiento,
cristaliza en L2; \emph{annoyance}, con 2\,470~ejemplos (un factor
3$\times$), cristaliza en L8. Más datos no adelantan la cristalización.
La profundidad está determinada por la naturaleza semántica de
la emoción (la complejidad de las distinciones lingüísticas que el
modelo necesita construir), no por el volumen de datos disponible.

Esta observación tiene una implicación práctica relevante para la
construcción de datasets de emociones: añadir más ejemplos de emociones
contextualmente complejas no las hace más fáciles de aprender. La
limitación es estructural, no estadística. Si fuera puramente estadística
(falta de datos), \emph{annoyance} cristalizaría más temprano a medida
que aumenta su frecuencia; lo que observamos es lo contrario.

La potencia estadística es limitada por construcción ($n=23$ emociones).
Para acotar la incertidumbre con honestidad, se computa el intervalo de
confianza al 95\% de cada coeficiente mediante \emph{bootstrap por
emoción} (10\,000 réplicas, semilla 42, resampleo de pares emoción
con reemplazo): la correlación cristalización--F1 máximo
($\rho = -0{,}673$) tiene CI95\,$=[-0{,}87, -0{,}33]$, estrictamente
negativo y por tanto significativo; la correlación
frecuencia--cristalización ($\rho = -0{,}265$) tiene
CI95\,$=[-0{,}65, +0{,}19]$, que cruza cero, confirmando estadísticamente
que la no significancia ($p=0{,}22$) refleja ausencia de evidencia, no
evidencia de ausencia. Con $n=23$, el intervalo es necesariamente
amplio; descartar definitivamente la influencia de la frecuencia
requeriría un dataset con más emociones, vía propuesta en
\S\ref{sec:trabajo-futuro}.


% =============================================================
% 5.2 — LOGIT LENS
% =============================================================

\subsection{\emph{Logit lens}: la brecha entre información disponible e información leíble}
\label{sec:logit-lens}

\subsubsection{Diseño experimental y motivación}

El probing del apartado anterior establece cuándo la información
emocional se vuelve linealmente accesible en las representaciones
intermedias del modelo. Sin embargo, el probe es un clasificador
entrenado ad hoc sobre cada capa: detecta información que está
estructurada linealmente en alguna dirección, pero no
necesariamente en la geometría que utiliza la cabeza de clasificación
final del modelo. Una representación puede ser linealmente separable
según el probe y, simultáneamente, ilegible para el clasificador
entrenado, si la dirección de separación no coincide con la base sobre
la que opera la cabeza de clasificación.

Para distinguir información disponible de información
formateada para el clasificador, se aplica la técnica del logit
lens \cite{nostalgebraist2020logit}: tomar la representación del token
\texttt{[CLS]} de cada capa intermedia y aplicarle directamente la
cabeza de clasificación entrenada del modelo, sin re-entrenamiento. Las
probabilidades resultantes indican la predicción que el modelo emitiría
si tuviera que decidir desde esa capa, utilizando la cabeza ya
optimizada para el \texttt{[CLS]} de L11.

\paragraph{Limitación de calibración.}
Esta técnica tiene una limitación conocida que conviene reconocer
explícitamente: la cabeza de clasificación fue entrenada exclusivamente
sobre representaciones de L11 tras LayerNorm, por lo que aplicarla a
capas intermedias produce predicciones cuyas \emph{magnitudes
absolutas} no son directamente comparables al baseline. Sin embargo,
las comparaciones relativas, entre capas para una misma
emoción o entre emociones para una misma capa, preservan información
sobre la alineación geométrica de la representación con la base del
clasificador. Las cifras absolutas funcionan como diagnóstico de tendencias, no como medida calibrada. La forma cualitativa del patrón es robusta a la
descalibración (se corrobora con probing y patching), pero las
magnitudes específicas no deben interpretarse como probabilidades
absolutas. Aproximaciones más rigurosas como \textit{tuned
lens}~\cite{belrose2023eliciting} entrenan una transformación lineal
por capa para corregir el desajuste de distribución; se discute como
trabajo futuro en el Capítulo~\ref{sec:conclusiones}.

El experimento se realiza sobre 2\,300 ejemplos del conjunto de test
(submuestreo equilibrado entre los tres grupos de cristalización),
agregando los resultados según los grupos del probing
(\S\ref{sec:probing}, Tabla~\ref{tab:cristalizacion-grupos}).

\subsubsection{Patrón agregado: tres regímenes geométricos}

\begin{figure}[H]
\centering
\includegraphics[width=0.85\textwidth]{figures/cap5/cap5_iterative_u_es.png}
\caption{\emph{Logit lens} agregado sobre 2\,300 ejemplos del conjunto de test.
Tres curvas: top-1 sigmoid (probabilidad media de la emoción más
probable), gold sigmoid (probabilidad media de la emoción correcta), y
suma de las 23 sigmoids (masa total de probabilidad). El patrón en U
distingue tres regímenes geométricos sucesivos.}
\label{fig:logit-lens-agregado}
\end{figure}

Promediando sobre los 2\,300 ejemplos, la probabilidad asignada a la
emoción correcta (\emph{gold sigmoid}) sigue un patrón en forma de U a
lo largo del modelo (Figura~\ref{fig:logit-lens-agregado}): parte de un
valor inicial moderado ($\approx 0{,}24$ en L0), cae a casi cero en
capas medias ($\approx 0{,}02$ en L7--L8) y resurge bruscamente en
capas tardías hasta alcanzar 0{,}44 en L11. Este patrón distingue tres
regímenes geométricos cualitativamente distintos:

\paragraph{Régimen léxico (Emb--L2).}
El \texttt{[CLS]} permanece próximo al embedding crudo y activa
difusamente múltiples direcciones del clasificador. La probabilidad
media de la top-1 es alta ($\approx 0{,}6$), pero la masa total de las
23 emociones también lo es ($\approx 5{,}4$), reflejando una activación
distribuida sin discriminación. El clasificador detecta señal léxica
sin filtrar, no decisión emocional.

\paragraph{Régimen de transición (L3--L8).}
Las probabilidades colapsan \emph{para todas las emociones
simultáneamente}, incluida la correcta. Este valle universal indica
que el \texttt{[CLS]} reside en una región del espacio
representacional desalineada con la base del clasificador: la
información que el probing detecta como presente en estas capas
(\S\ref{sec:probing}) existe en una \emph{geometría interna del
modelo}, no proyectable directamente sobre las direcciones de
clasificación. Es la fase de cómputo intermedio.

\paragraph{Régimen de cristalización (L9--L11).}
La probabilidad de la emoción correcta aumenta abruptamente, mientras
que la suma total se mantiene moderada ($\approx 1{,}3$ en L11). El
modelo no solo amplifica la señal correcta: también suprime activamente
las incorrectas. La masa total de las 23 probabilidades es de 1{,}3,
indicando concentración de masa en pocas direcciones. Esta es la fase
de proyección sobre la base del clasificador.

La Figura~\ref{fig:logit-lens-grupos-aggregate} desagrega el mismo
patrón en U según los tres grupos de cristalización del probing
(\S\ref{sec:probing}), mostrando que aunque la \emph{forma} es
universal, la \emph{magnitud} de cada régimen depende del tipo de
emoción. Las emociones tempranas alcanzan mayor gold sigmoid en L11
(consolidación más limpia), mientras que las tardías presentan un valle
más profundo en L4--L8 y una recuperación más modesta.

\begin{figure}[H]
\centering
\includegraphics[width=0.85\textwidth]{figures/cap5/cap5_logit_lens_groups_es.png}
\caption{\emph{Logit lens} desagregado por grupo de cristalización emocional.
Las tres curvas reproducen el patrón en U pero con magnitudes y
tiempos distintos: las emociones tempranas (azul) alcanzan los valores
más altos en L11; las tardías (rojo) muestran el valle más profundo
y una recuperación más modesta. La forma universal del patrón sugiere
que el régimen geométrico es propiedad arquitectónica del modelo, no
de la emoción.}
\label{fig:logit-lens-grupos-aggregate}
\end{figure}

\subsubsection{La brecha probing-\emph{logit lens} revela el trabajo de las
capas tardías}

Superponer el F1 del probe (información disponible) con el gold sigmoid
del \emph{logit lens} (información leíble por el clasificador) sobre los
grupos de cristalización emocional revela un patrón sistemático
(Figura~\ref{fig:logit-lens-grupos}): la brecha entre ambas
curvas cuantifica el trabajo de alineación geométrica que realizan las
capas tardías, y su magnitud depende del tipo de emoción.

\begin{figure}[H]
\centering
\includegraphics[width=\textwidth]{figures/cap5/cap5_lens_vs_probe_es.png}
\caption{Comparación entre probe F1 (información disponible) y gold
sigmoid del \emph{logit lens} (información leíble por el clasificador) por
grupo de cristalización emocional. La región sombreada entre ambas
curvas representa información presente en la representación pero no
proyectable sobre la base del clasificador.
Arriba: emociones tempranas ($n = 9$, 900 frases).
Centro: emociones medias ($n = 8$, 800 frases).
Abajo: emociones tardías ($n = 6$, 600 frases).}
\label{fig:logit-lens-grupos}
\end{figure}

\begin{table}[ht]
\centering
\caption{Brecha entre F1 del probe y gold sigmoid del \emph{logit lens}, por
capa y grupo emocional. La brecha mide la información presente en la
representación pero no leíble por la cabeza de clasificación.}
\label{tab:logit-lens-brechas}
\begin{tabular}{lcccccc}
\toprule
\textbf{Grupo} & \textbf{L0} & \textbf{L4} & \textbf{L7} & \textbf{L9} & \textbf{L11}
& \textbf{Brecha máx.} \\
\midrule
Tempranas & 0{,}37 & 0{,}57 & 0{,}63 & 0{,}62 & 0{,}09 & 0{,}63 (en L7) \\
Medias    & 0{,}02 & 0{,}29 & 0{,}50 & 0{,}52 & 0{,}14 & 0{,}52 (en L9) \\
Tardías   & $-0{,}05$ & 0{,}15 & 0{,}35 & 0{,}44 & 0{,}13 & 0{,}44 (en L9) \\
\bottomrule
\end{tabular}
\end{table}

Este patrón demuestra que las capas tardías realizan dos tipos
cualitativamente distintos de trabajo según el contenido:

\paragraph{Para emociones tempranas,}
la información léxica está disponible desde L0 (probe F1 $\approx 0{,}6$
ya en L0, prácticamente plateau hasta L11), pero el clasificador no
puede leerla hasta L9--L11 (\emph{logit lens} en el suelo durante L4--L8 y
resurgir abrupto en L9--L11). La brecha máxima entre ambas curvas
alcanza 0{,}63 en L7 ---justo antes del bloque tardío--- y se cierra
abruptamente en L9--L11. Es decir, las capas L8--L11 realizan, para
estas emociones, un trabajo de alineación geométrica considerable:
rotar información ya formada en capas anteriores hacia la base del
clasificador, sin crear señal nueva.

\paragraph{Para emociones tardías,}
la curva del probe sube progresivamente entre L0 y L9 (de 0{,}20 a
0{,}48, sin plateau temprano), mostrando que la información se
\emph{construye gradualmente} a lo largo de las capas, no está
disponible desde el principio. La curva del \emph{logit lens} sigue el mismo
perfil con cierto retraso, alcanzando el máximo en L11. La brecha
máxima es de solo 0{,}44 en L9. Para estas emociones, formación
y alineación se realizan en paralelo, no en fases separadas.

\paragraph{Las emociones medias}
ocupan una posición intermedia: la información se forma entre L0 y L7
(probe sube de 0{,}32 a 0{,}55), y la alineación se concentra en
L9--L11. Brecha máxima de 0{,}52 en L9.

\subsubsection{Conexión con \emph{activation patching}}

Esta interpretación geométrica explica un resultado central del
\emph{activation patching} que se presenta a continuación
(\S\ref{sec:activation-patching}): por qué la restauración exclusiva de
la capa~11 desde un modelo colapsado recupera el 100\% del rendimiento.
La capa~11 no contiene ``todo el conocimiento emocional'' en sentido
literal (el probing muestra que la información se distribuye a lo
largo de todas las capas), pero es donde culmina la proyección
de esa información sobre la base del clasificador. Restaurar L11
reconstituye la última etapa de alineación geométrica, y la información
residual que sobrevive a las capas comprimidas anteriores resulta
suficiente para que el clasificador produzca predicciones correctas.

Esto reconcilia tres hallazgos aparentemente disonantes que se
desarrollarán en las secciones siguientes:

\begin{enumerate}[nosep]
\item el probing localiza la formación de información en capas
tempranas y medias;
\item el logit lens muestra que esa información no es leíble hasta
capas tardías;
\item el activation patching identifica L11 como suficiente para la
recuperación.
\end{enumerate}

Las tres mediciones son consistentes si se distingue entre tres
conceptos relacionados pero distintos: representación (qué
información está presente, medido por probing), alineación
(esa información está proyectada sobre la base del clasificador, medido
por \emph{logit lens}) y acceso (qué componentes son suficientes
para que el clasificador opere, medido por patching). Esta clarificación
tripartita es, en sí misma, una contribución metodológica del trabajo:
sugiere que cualquier análisis interpretativo que confunda los tres
conceptos producirá conclusiones contradictorias como las que el
probing y el patching parecen producir aisladamente.
% =============================================================
% 5.3 — ACTIVATION PATCHING
% =============================================================

\subsection{\emph{Activation patching}: localización del conocimiento emocional}
\label{sec:activation-patching}

\subsubsection{Diseño experimental}

El probing identifica dónde se forma la información emocional;
el \emph{logit lens} identifica dónde se alinea con la base del
clasificador; el \emph{activation patching} identifica \emph{qué componentes
son funcionalmente necesarios} para que el modelo opere. Las tres
preguntas son complementarias: una capa puede crear poca información
nueva (bajo gain en probing) pero ser imprescindible para que la
información existente llegue al clasificador en formato utilizable.

El experimento parte del modelo comprimido con SVD uniforme a rango 64,
cuyo rendimiento es F1 macro = 0{,}000 (colapso total documentado en
\S\ref{sec:compresion-uniforme}). Desde este estado de destrucción
completa, se restauran los pesos originales de una capa o componente
a la vez y se mide la recuperación del F1 macro. El porcentaje de
restauración se calcula como
$\text{F1}_{\text{patched}} / \text{F1}_{\text{baseline}} \times 100$,
donde $\text{F1}_{\text{baseline}} = 0{,}577$.

\paragraph{Naturaleza funcional vs.\ causal.}
El sentido en el que este experimento establece causalidad merece una matización: la corrupción mediante SVD truncada no es neutral
en sentido estadístico (a diferencia del \emph{causal tracing} clásico de
Meng et al.~\cite{meng2022locating}, que usa noise gaussiano), sino una
\emph{distorsión estructurada}. Esto significa que las conclusiones
que extraemos son de naturaleza funcional (qué componentes
son suficientes para reactivar el modelo desde el colapso) más que
estrictamente causal en el sentido de Pearl. La distinción es
relevante para la replicabilidad: el \emph{activation patching} aquí responde
a la pregunta ``¿qué componente del modelo original es suficiente
para superar la perturbación SVD?'', no a la pregunta ``¿qué
componente codifica causalmente el conocimiento emocional?''. La
distinción no debilita los resultados ---los hace más precisos en su
alcance epistémico.

\subsubsection{Patching por capa completa}

La restauración capa por capa (Tabla~\ref{tab:patching-capas}) revela
una concentración extrema de la capacidad funcional en las capas
finales.

\begin{table}[ht]
\centering
\caption{Restauración del F1 macro al parchear una sola capa del modelo
comprimido (rango 64). Se omiten las capas con restauración nula
(0--7).}
\label{tab:patching-capas}
\begin{tabular}{ccc}
\toprule
\textbf{Capa restaurada} & \textbf{F1 Macro} & \textbf{Restauración (\%)} \\
\midrule
0--7 & 0{,}000 & 0{,}0 \\
8    & 0{,}001 & 0{,}1 \\
9    & 0{,}024 & 4{,}1 \\
10   & 0{,}139 & 24{,}1 \\
11   & 0{,}577 & 100{,}0 \\
\bottomrule
\end{tabular}
\end{table}

Las capas 0--7 no contribuyen en absoluto a la recuperación. La capa~11
sola es \emph{suficiente} para restaurar el 100\% del rendimiento
baseline. La interpretación correcta de este resultado, anticipada por
el \emph{logit lens} (\S\ref{sec:logit-lens}), es que L11 es donde culmina la
alineación geométrica de la representación con la base del
clasificador: restaurar L11 reconstituye la última proyección, y la
señal residual que sobrevive a 10 capas comprimidas resulta suficiente
para que el clasificador produzca predicciones correctas.

Una posible objeción al resultado de patching de L11 conviene anticiparla.
Podría argumentarse que el clasificador está conectado precisamente a
L11 mediante una capa lineal entrenada, y que por tanto la
``suficiencia'' de L11 es trivial. La objeción no aplica, por dos motivos. Primero, el clasificador no está conectado solo a la salida
cruda de L11 sino a L11 \emph{tras LayerNorm} y, en BERT, tras la
cabeza \textit{pooler} (CLS de L11 $\rightarrow$ tanh $\rightarrow$
linear). Restaurar los pesos de L11 incluye las seis matrices del
encoder pero \emph{no} la cabeza pooler ni el clasificador, que ya
estaban no comprimidos en el modelo \textit{corrupted run}. El hallazgo
no es ``el clasificador ve L11'' (eso sería trivial) sino ``L11 tras 10 capas comprimidas todavía produce señal suficiente para que el
clasificador clasifique correctamente'', lo cual no es trivial.
Segundo, el \emph{activation patching} sobre las capas 8--10 muestra
restauración parcial pero no completa, lo que indica que la suficiencia
de L11 no es un artefacto de localización: cuando otras capas
``cercanas'' al clasificador se restauran, la recuperación es parcial.
La especificidad funcional de L11 es genuina.

\subsubsection{La FFN como locus computacional crítico}

El análisis por componente (Tabla~\ref{tab:patching-componentes})
desglosa cada capa en sus dos subcomponentes (atención multi-cabeza y
red feed-forward) y revela una asimetría funcional de primer
orden.

\begin{table}[ht]
\centering
\caption{Restauración del F1 macro al parchear un solo componente
(atención o FFN) de las capas 8--11 del modelo comprimido.}
\label{tab:patching-componentes}
\begin{tabular}{cccc}
\toprule
\textbf{Capa} & \textbf{Componente} & \textbf{F1 Macro} & \textbf{Restauración (\%)} \\
\midrule
11 & FFN       & 0{,}577 & \textbf{100{,}0} \\
11 & Atención  & 0{,}365 & 63{,}3 \\
10 & Atención  & 0{,}156 & 27{,}1 \\
10 & FFN       & 0{,}139 & 24{,}1 \\
9  & Atención  & 0{,}033 & 5{,}6 \\
9  & FFN       & 0{,}024 & 4{,}1 \\
8  & Atención  & 0{,}000 & 0{,}0 \\
8  & FFN       & 0{,}001 & 0{,}1 \\
\bottomrule
\end{tabular}
\end{table}

El resultado más destacado es que restaurar únicamente la FFN
de la capa~11, aproximadamente un tercio de los parámetros de esa
capa, recupera el 100\% del rendimiento. La atención de la misma
capa solo alcanza un 63{,}3\%, y con alta variabilidad entre emociones.
Este hallazgo posiciona a la FFN de la capa~11 como el locus
computacional crítico de la capacidad emocional del modelo.

La dominancia de la FFN desafía productivamente la narrativa centrada
en la atención que ha prevalecido desde \emph{Attention Is All You
Need} \cite{vaswani2017attention}. En este modelo \emph{fine-tuned} para
clasificación emocional, la atención contribuye a enrutar información
entre posiciones, pero es la FFN la que realiza la transformación
computacional que convierte representaciones contextualizadas en
representaciones clasificables. Este resultado es coherente con
trabajos recientes que identifican las FFN como \textit{memorias
asociativas} del modelo \cite{geva2021transformer}, y conecta
directamente con el hallazgo del \emph{logit lens}: la proyección sobre la
base del clasificador que documentamos en \S\ref{sec:logit-lens} es
realizada principalmente por las FFN de las capas tardías, no por la
atención.

\subsubsection{Restauración diferencial por emoción}

La restauración de la atención de la capa~11 presenta una marcada
heterogeneidad entre emociones (Tabla~\ref{tab:patching-emociones}).

\begin{table}[ht]
\centering
\caption{Restauración del F1 por emoción al parchear solo la atención
de la capa~11. Se muestran las 5 emociones mejor y peor restauradas.}
\label{tab:patching-emociones}
\begin{tabular}{lclc}
\toprule
\textbf{Bien restauradas} & \textbf{Rest.\ (\%)} &
\textbf{Mal restauradas} & \textbf{Rest.\ (\%)} \\
\midrule
gratitude   & 100{,}0 & remorse        & 0{,}0 \\
annoyance   & 98{,}8  & fear           & 3{,}6 \\
admiration  & 95{,}4  & desire         & 4{,}4 \\
anger       & 94{,}5  & realization    & 10{,}0 \\
curiosity   & 91{,}1  & disappointment & 18{,}8 \\
\bottomrule
\end{tabular}
\end{table}

Las emociones que la atención restaura bien son aquellas con señales
léxicas claras: \emph{gratitude} se identifica por marcadores
explícitos, \emph{anger} por vocabulario de alta valencia negativa. La
atención puede redirigir estas señales al token \texttt{[CLS]} de
forma eficaz incluso con FFN comprimidas. En cambio, las emociones
que la atención no restaura (\emph{remorse}, \emph{fear},
\emph{desire}) son aquellas que requieren computación contextual
más rica, precisamente la función que desempeña la FFN.

Este patrón refuerza la interpretación de los grupos de
cristalización del probing (\S\ref{sec:probing}) y del \emph{logit lens}
(\S\ref{sec:logit-lens}): las emociones tempranas son léxicas y
dependen menos de la FFN profunda; las emociones tardías son
contextuales y necesitan la FFN intacta para completar tanto la
formación como la alineación de su representación.


% =============================================================
% 5.4 — ABLACIÓN DE CABEZAS DE ATENCIÓN
% =============================================================

\subsection{Ablación de cabezas de atención}
\label{sec:head-ablation}

\subsubsection{Diseño experimental}

Los experimentos anteriores operan a nivel de capa completa o de
subcomponente (atención vs.\ FFN). La ablación de cabezas desciende a
la granularidad individual: para cada una de las 144~cabezas de
atención (12~capas $\times$ 12~cabezas), se anulan sus pesos
(\textit{zero-out}) y se mide la caída en F1 macro y en F1 por
emoción. Para mantener tractable el coste computacional ---144~cabezas
$\times$ una pasada de evaluación cada una--- la ablación se ejecuta
sobre un submuestreo de 2\,000 ejemplos del conjunto de test (en lugar
del test completo de 5\,427). El baseline local sobre este submuestreo
es 0{,}582, ligeramente superior al baseline global de 0{,}577 del
test completo (\S\ref{sec:met-modelo}); todas las caídas reportadas en
esta sección son absolutas respecto a ese baseline local, lo que las
mantiene comparables entre sí.

A partir de estas mediciones se construyen dos métricas ortogonales:
la \emph{importancia}, definida como el impacto agregado en F1 a
través de todas las emociones, y la \emph{selectividad}, que mide la
concentración de dicho impacto en pocas emociones (valor alto = efecto
concentrado en 1--2 emociones; valor bajo = efecto distribuido). La
ortogonalidad de ambas dimensiones se verifica empíricamente
($r = -0{,}055$), lo que habilita una clasificación 2$\times$2 de las
cabezas. La consistencia interna de la métrica de importancia se
valida mediante su correlación con la caída real de F1 por ablación
($r = 0{,}71$).

\subsubsection{Taxonomía funcional}

Las 144~cabezas se clasifican en cuatro categorías
(Tabla~\ref{tab:head-taxonomy}).

\begin{table}[ht]
\centering
\caption{Taxonomía de cabezas de atención basada en importancia y
selectividad. La baja correlación entre ambas dimensiones
($r = -0{,}055$) valida la clasificación.}
\label{tab:head-taxonomy}
\begin{tabular}{lcccl}
\toprule
\textbf{Categoría} & \textbf{N} & \textbf{Import.} & \textbf{Select.} & \textbf{Rol funcional} \\
\midrule
Critical Specialist & 37 & 0{,}282 & 0{,}680 & Especializada en emociones concretas \\
Critical Generalist & 35 & 0{,}265 & 0{,}544 & Infraestructura emocional general \\
Minor Specialist    & 35 & 0{,}101 & 0{,}702 & Nicho funcional (1--2 emociones) \\
Dispensable         & 37 & 0{,}120 & 0{,}556 & Candidata a poda \\
\bottomrule
\end{tabular}
\end{table}

\subsubsection{Gradiente de criticidad por profundidad}

El hallazgo principal de este análisis es que la proporción de cabezas
críticas crece monotónicamente con la profundidad
(Tabla~\ref{tab:head-profundidad}, Figura~\ref{fig:head-taxonomy-depth}).

\begin{table}[ht]
\centering
\caption{Distribución de cabezas críticas y dispensables por bloque de
profundidad. La capa 11 no contiene ninguna cabeza dispensable.}
\label{tab:head-profundidad}
\begin{tabular}{lcc}
\toprule
\textbf{Bloque} & \textbf{Críticas} & \textbf{Dispensables} \\
\midrule
Capas 0--4  & 16/60 (27\%) & 22/60 (37\%) \\
Capas 5--7  & 19/36 (53\%) & 10/36 (28\%) \\
Capas 8--11 & 37/48 (77\%) & \phantom{0}5/48 (10\%) \\
\bottomrule
\end{tabular}
\end{table}

La capa~11 constituye un caso extremo: ninguna de sus 12~cabezas es
dispensable: las 12 son críticas (6~specialists y 6~generalists). La
importancia media por capa presenta un salto abrupto a partir de la
capa~8 ($\sim 0{,}25$ vs.\ $\sim 0{,}15$ en capas anteriores), con la
capa~11 alcanzando una importancia media de 0{,}381. Las 5~cabezas más
importantes del modelo pertenecen a las capas 10--11, encabezadas por
L11-H11 con una caída de F1 macro de 0{,}023.

\begin{figure}[H]
\centering
\includegraphics[width=0.85\textwidth]{figures/cap5/cap5_head_categories_es.png}
\caption{Distribución de categorías de cabezas por capa. La proporción
de cabezas críticas (specialists + generalists) crece con la
profundidad. La capa~11 no contiene ninguna cabeza dispensable.}
\label{fig:head-taxonomy-depth}
\end{figure}

Este resultado proporciona triangulación con cuatro métodos
independientes: el probing (\S\ref{sec:probing}) muestra que la
separabilidad lineal máxima se alcanza en capas finales; el \emph{logit lens}
(\S\ref{sec:logit-lens}) identifica L9--L11 como la fase de
cristalización donde la representación se alinea con el clasificador;
el \emph{activation patching} (\S\ref{sec:activation-patching}) demuestra
que la capa~11 es suficiente para restaurar el rendimiento completo;
y la ablación confirma que cada cabeza individual de esa capa es
funcionalmente necesaria. Cinco técnicas distintas convergen en la
misma conclusión: las capas 8--11 son el núcleo funcional del modelo
emocional.

\subsubsection{Cabezas interferentes}

Un hallazgo secundario es la existencia de 21~cabezas cuya ablación
\emph{mejora} el rendimiento del modelo (caída de F1 más negativa que
$-0{,}003$). Aplicando un umbral más laxo ($\Delta F1 < 0$) el conteo
sube a 79 cabezas, pero la mayoría de esos efectos están en el rango
del ruido de medición; las 21~cabezas que superan el umbral de
$0{,}003$ son las únicas con un beneficio cuantificable. La
Tabla~\ref{tab:cabezas-interferentes} lista las diez con mayor efecto.

\begin{table}[ht]
\centering
\caption{Top 10 cabezas interferentes (ordenadas por mejora del F1
macro tras ablación). Categoría según la taxonomía de
\S\ref{sec:head-ablation}.}
\label{tab:cabezas-interferentes}
\small
\begin{tabular}{c c l c c}
\toprule
\textbf{Capa} & \textbf{Cabeza} & \textbf{Categoría} & \textbf{$\Delta$F1 (mejora)} & \textbf{Selectividad} \\
\midrule
11 &  5 & Critical Specialist  & $+$0{,}0121 & 0{,}635 \\
10 & 11 & Critical Generalist  & $+$0{,}0077 & 0{,}596 \\
 1 & 10 & Critical Specialist  & $+$0{,}0051 & 0{,}674 \\
 2 &  1 & Critical Generalist  & $+$0{,}0047 & 0{,}469 \\
 3 &  7 & Critical Specialist  & $+$0{,}0047 & 0{,}750 \\
 7 &  7 & Minor Specialist     & $+$0{,}0046 & 0{,}829 \\
 7 &  1 & Critical Generalist  & $+$0{,}0045 & 0{,}583 \\
 5 &  1 & Critical Specialist  & $+$0{,}0043 & 0{,}673 \\
 7 &  9 & Critical Specialist  & $+$0{,}0043 & 0{,}703 \\
 5 &  2 & Critical Specialist  & $+$0{,}0043 & 0{,}621 \\
\bottomrule
\end{tabular}
\end{table}

La distribución por profundidad es informativa: 7 interferentes en
capas tempranas (L0--L4), 9 en intermedias (L5--L7) y solo 5 en
tardías (L8--L11). El 76\% se concentra en capas no críticas para la
clasificación final, lo que es coherente con la taxonomía
2$\times$2: las cabezas dispensable y las críticas-generalistas son
las que más contribuyen al ruido cuando se acumulan. La cabeza más
interferente es L11-H5 ($\Delta F1 = +0{,}012$), que además es
nominalmente \emph{Critical Specialist}; este caso aparentemente
paradójico sugiere que su contribución específica entra en conflicto con
otras cabezas de L11, y eliminarla resuelve la interferencia.

Como estimación aditiva, ablacionar simultáneamente las 21~cabezas
interferentes daría una ganancia teórica de $+0{,}095$ en F1 macro
(suma directa de las caídas negativas), lo que llevaría el modelo de
0{,}577 a $\sim$0{,}672. El supuesto de aditividad es probablemente
demasiado optimista debido a interacciones no-lineales entre cabezas
ablacionadas, pero incluso una fracción significativa de esa ganancia
sería notable para una intervención sin re-entrenamiento. El
experimento de ablación coordinada se propone como trabajo futuro
(\S\ref{sec:trabajo-futuro}).

Este hallazgo es, en sí mismo, una observación metodológicamente
relevante: incluso un modelo \emph{fine-tuned} eficaz contiene componentes
contraproducentes, lo que abre la puerta a optimizaciones post-hoc
que no requieren acceso al proceso de entrenamiento ---ni siquiera a
los gradientes--- y se complementan naturalmente con la compresión
SVD del Capítulo~\ref{sec:resultados-compresion}.

\subsubsection{Especialización emocional}

Cada emoción depende de cabezas diferentes, concentradas
mayoritariamente en las capas 8--11
(Tabla~\ref{tab:head-emociones}).

\begin{table}[ht]
\centering
\caption{Cabeza más importante para cada emoción (top 10 por magnitud
de impacto). La caída de F1 indica cuánto pierde esa emoción al
ablacionar su cabeza más crítica.}
\label{tab:head-emociones}
\begin{tabular}{lcc}
\toprule
\textbf{Emoción} & \textbf{Cabeza crítica} & \textbf{$\Delta$F1} \\
\midrule
embarrassment  & L10-H1 & $-0{,}320$ \\
annoyance      & L10-H0 & $-0{,}090$ \\
disappointment & L10-H4 & $-0{,}080$ \\
realization    & L11-H6 & $-0{,}064$ \\
fear           & L0-H8  & $-0{,}056$ \\
confusion      & L9-H2  & $-0{,}044$ \\
sadness        & L11-H6 & $-0{,}041$ \\
disgust        & L9-H7  & $-0{,}040$ \\
joy            & L10-H7 & $-0{,}039$ \\
anger          & L8-H8  & $-0{,}031$ \\
\bottomrule
\end{tabular}
\end{table}

El caso de \emph{embarrassment} es extremo: pierde un 32\% de su F1
al ablacionar cualquiera de cuatro cabezas equivalentes
(L10-H1, L10-H4, L11-H2 y L11-H8, todas con $\Delta\text{F1} = -0{,}320$
por empate exacto en el ranking). Esta concentración funcional en un
puñado pequeño de cabezas redundantes ---no en una sola, pero sí en
muy pocas--- probablemente refleja la baja frecuencia de la emoción
en GoEmotions (303~ejemplos), que hace que el modelo aprenda un
circuito atencional acotado y replicado por seguridad estadística.
Desde la perspectiva de la compresión, esto implica que la eliminación
o degradación coordinada de esas cuatro cabezas destruiría
selectivamente la capacidad del modelo para detectar
\emph{embarrassment}, un ejemplo de sesgo silencioso por compresión
que las métricas agregadas no detectan fácilmente.

Otro patrón relevante es el compartimiento de circuitos:
\emph{realization} y \emph{sadness} dependen de la misma cabeza
(L11-H6), lo que sugiere un mecanismo atencional común
, posiblemente la detección de un patrón contextual de expectativa
frustrada que ambas emociones comparten.


% =============================================================
% 5.5 — ANÁLISIS NEURONAL DE LAS FFN
% =============================================================

\subsection{Análisis neuronal de las redes feed-forward}
\label{sec:neuron-analysis}

\subsubsection{Diseño experimental}

El \emph{activation patching} estableció la FFN de la capa~11 como el locus
computacional crítico. El análisis neuronal desciende al nivel de las
neuronas individuales dentro de las FFN para examinar dos preguntas:
cómo se distribuye la especialización emocional a través de las
36\,864~neuronas (12~capas $\times$ 3\,072~neuronas), y si la
compresión SVD daña selectivamente las neuronas más especializadas.

La selectividad de cada neurona para cada emoción se cuantifica
mediante un estadístico tipo Cohen's $d$: la diferencia entre la
activación media cuando la emoción está presente y cuando está
ausente, normalizada por la desviación típica agrupada. Se extraen
las activaciones post-GELU de la capa intermedia de cada FFN
(\texttt{BertIntermediate}) mediante \textit{forward hooks}, con
mean-pooling sobre posiciones de secuencia. Se definen como
\emph{significativas} las neuronas con $|d| > 2{,}0$.

Para investigar la conexión con la compresión, se proyecta el vector
de selectividad de cada emoción (un vector de 3\,072~dimensiones por
capa) sobre los subespacios SVD conservado y eliminado de la matriz
$W_{\text{int}}$ a rango 128, y se calcula la fracción de señal
emocional que cae en el subespacio eliminado.

\subsubsection{Especialización concentrada en capas profundas}

La distribución de neuronas significativas por profundidad
(Tabla~\ref{tab:neuron-profundidad}) replica el patrón ya observado
en probing, patching y ablación de cabezas.

\begin{table}[ht]
\centering
\caption{Distribución de neuronas especializadas ($|d| > 2{,}0$) por
bloque de profundidad. El 84\% de las neuronas significativas se
concentra en las capas 8--11.}
\label{tab:neuron-profundidad}
\begin{tabular}{lcc}
\toprule
\textbf{Bloque} & \textbf{Neuronas significativas} & \textbf{Selectividad máx.\ media} \\
\midrule
Capas 0--3  & 11 (0{,}3\%)    & 0{,}91--1{,}15 \\
Capas 4--7  & 570 (16\%)      & 1{,}50--2{,}19 \\
Capas 8--11 & 3\,061 (84\%)   & 3{,}10--3{,}92 \\
\bottomrule
\end{tabular}
\end{table}

El 84\% de las neuronas especializadas se concentra en las capas
8--11, con la capa~11 aportando 1\,127~neuronas significativas por sí
sola. Las capas 0--3 contienen prácticamente ninguna neurona con
especialización emocional detectable. Este resultado es coherente con
la función de las capas tempranas como procesadores léxicos genéricos:
contribuyen a la señal emocional (probing), pero lo hacen a través
de \emph{representaciones distribuidas}, no de neuronas individuales
especializadas.

\subsubsection{Asimetría radical entre emociones}

El número de neuronas significativas varía en órdenes de magnitud
entre emociones (Tabla~\ref{tab:neuron-emociones}).

\begin{table}[ht]
\centering
\caption{Neuronas significativas y selectividad máxima por emoción
(selección). Tres emociones no tienen ninguna neurona con
$|d| > 2{,}0$.}
\label{tab:neuron-emociones}
\begin{tabular}{lcc}
\toprule
\textbf{Emoción} & \textbf{Neuronas sig.} & \textbf{Selectividad máx.} \\
\midrule
gratitude      & 818 & 6{,}88 \\
remorse        & 442 & 7{,}75 \\
love           & 399 & 5{,}41 \\
amusement      & 394 & 5{,}90 \\
fear           & 306 & 7{,}71 \\
\midrule
annoyance      & 0   & 1{,}71 \\
disappointment & 0   & 1{,}89 \\
realization    & 0   & 1{,}96 \\
\bottomrule
\end{tabular}
\end{table}

Las emociones mejor codificadas (\emph{gratitude}, \emph{remorse})
disponen de cientos de neuronas dedicadas, mientras que
\emph{annoyance}, \emph{disappointment} y \emph{realization} carecen
de neuronas individuales significativas. Las neuronas con mayor
selectividad de todo el modelo ($d > 7$) pertenecen a \emph{fear} y
\emph{remorse}, ambas en la capa~11.

Esta asimetría sugiere que el modelo emplea estrategias de
codificación cualitativamente diferentes: las emociones ``fuertemente
escritas'' utilizan neuronas altamente especializadas, mientras que
las emociones ``débilmente escritas'' se codifican mediante patrones
distribuidos más sutiles que la métrica de selectividad individual no
captura.

\subsubsection{La hipótesis SVD-neuronas: un resultado nulo informativo}

La hipótesis principal de este análisis era que la compresión SVD
dañaría selectivamente las emociones cuyas neuronas especializadas
cayeran desproporcionadamente en el subespacio eliminado. Esta
hipótesis resulta \emph{falsa}: la correlación entre la fracción de
señal en el subespacio eliminado y la caída de F1 por emoción es
prácticamente nula (Spearman $\rho = -0{,}004$, $p = 0{,}99$).

Además, la fracción de señal emocional que cae en el subespacio
eliminado es prácticamente idéntica para todas las emociones (rango
0{,}53--0{,}56, media 0{,}54). La SVD no discrimina entre señales
emocionales: elimina señal de forma uniforme, independientemente de
la identidad funcional de las neuronas afectadas.

Este resultado nulo es informativo porque revela una disociación clara: la geometría que la SVD optimiza (energía de valores
singulares) es ortogonal a la estructura funcional del modelo
(especialización emocional). La compresión opera en un espacio
matemático que no ``ve'' la función. Esto es teóricamente importante
porque clarifica por qué la compresión SVD informada por
interpretabilidad funciona: no porque la SVD pueda dirigirse a
neuronas específicas, sino porque puede asignar más rango a los
componentes (FFN tardía) cuya destrucción causa más daño.

Este patrón es consistente con la hipótesis de \emph{superposition}
formalizada por Elhage et al.~\cite{elhage2022toy}: cuando una red
necesita representar más \emph{features} funcionales que dimensiones
disponibles, las distribuye en direcciones casi-ortogonales del
espacio de activación, aceptando pequeñas interferencias a cambio de
mayor capacidad. En esta lectura, las features emocionales en BERT
estarían codificadas en direcciones que no coinciden con las
direcciones SVD principales de los pesos: las primeras provienen del
optimizador de la tarea durante el \emph{fine-tuning}; las segundas
provienen de la geometría algebraica de los pesos. No hay razón a
priori para que coincidan, y nuestro resultado empírico (fracción de
señal eliminada uniforme entre emociones) es exactamente lo que
predice superposition. La implicación práctica es que los métodos de
compresión geométrica como la SVD truncada nunca podrán ser
\emph{selectivamente} funcionales por sí solos: la selectividad
funcional debe entrar por la puerta de la sensibilidad empírica
medida ---es decir, justo lo que el algoritmo greedy del
Capítulo~\ref{sec:sintesis} hace al usar las caídas de F1 por
componente como costes de optimización.

\subsubsection{La norma de selectividad como predictor de vulnerabilidad}

Si la geometría SVD no explica la vulnerabilidad diferencial, ¿qué la
explica? Un análisis sistemático testa seis predictores posibles de
la caída de F1 bajo compresión
(Tabla~\ref{tab:neuron-predictores}).

\begin{table}[ht]
\centering
\caption{Predictores de la vulnerabilidad emocional a compresión SVD
($n = 23$). La norma de selectividad, que mide la intensidad total
de codificación neuronal, es el mejor predictor.}
\label{tab:neuron-predictores}
\begin{tabular}{lccc}
\toprule
\textbf{Predictor} & $\rho$ & \textbf{p-valor} & \textbf{Sig.} \\
\midrule
Fracción en subespacio eliminado & $-0{,}004$ & $0{,}986$ & ns. \\
Nº ejemplos de entrenamiento     & $+0{,}174$ & $0{,}427$ & ns. \\
Aislamiento representacional     & $+0{,}452$ & $0{,}031$ & * \\
Selectividad máxima              & $+0{,}517$ & $0{,}012$ & * \\
Nº neuronas significativas       & $+0{,}581$ & $0{,}004$ & ** \\
\textbf{Norma de selectividad}   & $\mathbf{+0{,}639}$ & $\mathbf{0{,}001}$ & *** \\
\bottomrule
\end{tabular}
\end{table}

El mejor predictor es la \emph{norma del vector de selectividad}
($\rho = 0{,}639$, $p = 0{,}001$), que mide la intensidad total con la
que el modelo codifica cada emoción a través de sus neuronas FFN. Las
emociones con codificación más intensa (\emph{gratitude}: norma 140,
\emph{remorse}: 111, \emph{love}: 105) sufren más con la compresión
que las de codificación débil (\emph{disappointment}: 42,
\emph{realization}: 43).

La interpretación es la siguiente: la SVD no ataca selectivamente las
neuronas emocionales: la geometría del subespacio eliminado es
ortogonal a la especialización funcional, como se ha demostrado. Pero
las emociones que el modelo codifica con mayor intensidad neuronal
son más vulnerables a perturbaciones generales de los pesos,
porque su representación depende de configuraciones más específicas y
exigentes. Es una vulnerabilidad por intensidad, no por
geometría. La distinción es útil metodológicamente: sugiere
que la robustez a perturbación general (no solo a SVD) puede
estimarse a partir de la intensidad de codificación, una métrica
intrínseca al modelo que no requiere experimentos de compresión.

\subsubsection{Clusters emocionales emergentes}

La agrupación jerárquica sobre los vectores de selectividad
concatenados produce seis clusters emocionales
(Tabla~\ref{tab:neuron-clusters}) que reflejan categorías psicológicas
coherentes.

\begin{table}[ht]
\centering
\caption{Clusters emocionales emergentes de la estructura neuronal del
modelo. Las agrupaciones se obtienen por clustering jerárquico sobre
los vectores de selectividad de 36\,864~neuronas.}
\label{tab:neuron-clusters}
\begin{tabular}{cl}
\toprule
\textbf{Cluster} & \textbf{Emociones} \\
\midrule
Positivas de alta energía & admiration, amusement, excitement, gratitude, joy, love \\
Negativas reactivas       & anger, annoyance, disappointment, disapproval, disgust \\
Negativas internas        & embarrassment, fear, remorse, sadness \\
Epistémicas               & confusion, curiosity, surprise \\
Orientadas al otro        & caring, desire, optimism \\
Baja especificidad        & approval, realization \\
\bottomrule
\end{tabular}
\end{table}

Estos clusters no se imponen \emph{a priori} sino que emergen
directamente de la estructura de activaciones del modelo. Que
reproduzcan categorías reconocibles de la psicología de las
emociones (positivas vs.\ negativas, reactivas vs.\
internas, epistémicas vs.\ afectivas) sugiere que el \emph{fine-tuning} ha
organizado las representaciones internas de BERT de forma
psicológicamente coherente, no solo estadísticamente eficaz. Esta
correspondencia es no trivial: el modelo no ha sido entrenado con
información sobre la estructura psicológica de las emociones, y sin
embargo la recupera implícitamente.


% =============================================================
% 5.6 — SÍNTESIS DEL CAPÍTULO
% =============================================================

\subsection{Síntesis: una arquitectura funcional convergente}
\label{sec:interp-sintesis}

Las cinco técnicas de este capítulo emplean métodos diferentes
(clasificadores lineales entrenados \textit{ad hoc}, aplicación
directa del clasificador final, restauración selectiva de pesos,
ablación individual de cabezas, estadísticas de activación neuronal)
pero convergen en una descripción unificada de cómo BERT procesa la
información emocional. La Tabla~\ref{tab:convergencia-tecnicas}
resume esta convergencia.

\begin{table}[ht]
\centering
\caption{Convergencia de cinco técnicas de interpretabilidad en las
mismas conclusiones estructurales.}
\label{tab:convergencia-tecnicas}
\small
\begin{tabular}{p{2.8cm}p{3.5cm}p{6.5cm}}
\toprule
\textbf{Técnica} & \textbf{Pregunta} & \textbf{Resultado clave} \\
\midrule
Probing &
¿Dónde se \emph{forma} la señal? &
L0 captura el 61\% (léxica); L8 segundo pico (contextual).
Capas finales aportan $< 1\%$. \\
\addlinespace
\emph{Logit lens} &
¿Dónde es leíble por el clasificador? &
Patrón en U: alta en L0--L2, valle en L3--L8, cristalización
abrupta en L9--L11. La brecha probing--\emph{logit lens} cuantifica el
trabajo de alineación. \\
\addlinespace
\emph{Activation patching} &
¿Dónde \emph{reside} el conocimiento? &
Solo L11 restaura el 100\%. La FFN sola basta;
la atención solo recupera el 63\%. \\
\addlinespace
Ablación de cabezas &
¿Qué cabezas son \emph{necesarias}? &
77\% críticas en capas 8--11; 0 dispensables en capa~11.
Gradiente monotónico con profundidad. \\
\addlinespace
Análisis neuronal &
¿Qué neuronas están \emph{especializadas}? &
84\% de neuronas significativas en capas 8--11.
La norma de selectividad ($\rho = 0{,}64$) predice vulnerabilidad. \\
\bottomrule
\end{tabular}
\end{table}

De esta convergencia se extraen cuatro conclusiones principales que
constituyen la base sobre la que el Capítulo~\ref{sec:sintesis}
construirá la estrategia de compresión informada.

\paragraph{Arquitectura de tres niveles.}
BERT procesa la información emocional en tres fases cualitativamente
distintas. Las capas tempranas (Emb--L2) extraen señal léxica bruta:
información linealmente separable pero distribuida sobre muchas
direcciones simultáneamente. Las capas intermedias (L3--L8) realizan
cómputo de transición en una geometría \emph{interna} del modelo: la
información se refina y reorganiza, pero permanece desalineada con la
base del clasificador (valle del \emph{logit lens}). Las capas profundas
(L9--L11) ejecutan la cristalización: alinean la representación con
la cabeza de clasificación y suprimen activamente las opciones
incorrectas. Este último paso es el cuello de botella geométrico del modelo, y su criticidad bajo compresión SVD se
explica por su baja redundancia espectral en las matrices que
realizan la alineación.

\paragraph{La paradoja información-dependencia, resuelta.}
El probing muestra que las capas finales (L10--L11) apenas crean
información nueva (gain $< 0{,}004$ del total), pero el activation
patching y la ablación demuestran que son las más críticas para el
rendimiento del modelo. El \emph{logit lens} resuelve esta aparente paradoja:
las capas finales no \emph{crean} señal emocional, sino que la rotan
a la geometría del clasificador. La brecha entre probe F1 (información
presente) y gold sigmoid del \emph{logit lens} (información leíble)
cuantifica directamente este trabajo de alineación, que varía
sistemáticamente entre 0{,}63 (emociones tempranas) y 0{,}44 (tardías)
en su valor máximo. Las capas tardías son el cuello de botella entre representación y predicción, y por eso resultan frágiles bajo
compresión.

\paragraph{Las capas tardías hacen dos tipos de trabajo distintos.}
El análisis por grupo emocional revela que las capas L9--L11 no
realizan una única operación, sino dos cualitativamente distintas
según el contenido. Para información \emph{léxica ya formada}
(emociones tempranas como \emph{gratitude}), realizan
alineación pura: rotar una representación ya construida a la
base del clasificador. Para información \emph{contextual} que se
construye a lo largo del modelo (emociones tardías como
\emph{annoyance}), realizan formación y alineación coordinadas. La brecha entre probing y \emph{logit lens} distingue ambos
regímenes cuantitativamente.

\paragraph{Disociación entre geometría SVD y función.}
La compresión SVD opera sobre la estructura de valores singulares de
las matrices de pesos, una geometría que es ortogonal a la
organización funcional del modelo. La SVD no ``ve'' qué neuronas
codifican qué emociones: elimina señal uniformemente. Sin embargo, la
vulnerabilidad diferencial existe y se explica por la intensidad de
codificación: las emociones que el modelo escribe con más fuerza en
sus neuronas son más sensibles a cualquier perturbación de los pesos,
incluida la compresión.

Estas cuatro conclusiones fundamentan la estrategia de compresión
informada que se desarrolla en el Capítulo~\ref{sec:sintesis}, donde
la profundidad de cristalización y la intensidad de codificación
neuronal se utilizan como predictores para guiar decisiones de
compresión selectiva. El conjunto del capítulo establece que la
interpretabilidad mecánica no produce una colección de hallazgos
disjuntos sobre el modelo, sino una descripción \emph{unificada} de su
arquitectura funcional, descripción que cinco métodos
metodológicamente independientes corroboran.
% =============================================================
% CAPÍTULO 6 — Síntesis: de la interpretabilidad a la compresión informada
% Versión: matrícula de honor
% =============================================================

\section{Síntesis: compresión informada por interpretabilidad}
\label{sec:sintesis}

Los dos capítulos anteriores han recorrido caminos complementarios. El
Capítulo~\ref{sec:resultados-compresion} midió qué se rompe al
comprimir BERT con SVD: la sensibilidad varía radicalmente entre
componentes (factor 14$\times$ a rango~128 en retención absoluta y hasta
72$\times$ a rango~64 normalizado por parámetros), entre profundidades
(transición de fase entre rango~384 y 256) y entre emociones (desde
\emph{gratitude}, que retiene el 99\% del F1 a rango~384, hasta
\emph{embarrassment}, que colapsa al 0\% al mismo rango). El
Capítulo~\ref{sec:resultados-interpretabilidad} investigó por qué
mediante cinco técnicas de interpretabilidad mecánica que convergen en una
arquitectura funcional de tres niveles: información léxica formada en
capas tempranas, cómputo de transición en capas medias, y alineación
geométrica con la base del clasificador en capas tardías, donde reside el
cuello de botella computacional del modelo.

Este capítulo cierra el bucle: utiliza los hallazgos de interpretabilidad
mecánica para diseñar estrategias de compresión \emph{informadas} y evalúa
si ese conocimiento se traduce en ventaja cuantitativa sobre las
estrategias ciegas presentadas en el
Capítulo~\ref{sec:resultados-compresion}. La estructura del capítulo
refleja una progresión metodológica (intuición, fracaso de la
intuición, pivot a datos, dominancia) que constituye en sí misma una
observación replicable sobre cómo trasladar conocimiento mecanicista a
ingeniería aplicada.

El recorrido se organiza en cinco secciones. La primera establece la
conexión cuantitativa entre cristalización y vulnerabilidad
(\S\ref{sec:cristalizacion-vulnerabilidad}). La segunda describe la
estrategia heurística informada y documenta su fracaso
(\S\ref{sec:informed-heuristic}). La tercera presenta el algoritmo greedy
alimentado con datos empíricos de sensibilidad y su dominancia de la
frontera de Pareto (\S\ref{sec:greedy-algorithm}). La cuarta evalúa la
recuperación mediante \emph{fine-tuning} y reporta una observación inesperada
sobre regularización implícita (\S\ref{sec:finetuning-recovery}). La quinta
sintetiza las implicaciones del conjunto de resultados
(\S\ref{sec:sintesis-discusion}).


% -------------------------------------------------------------
\subsection{La cristalización predice la vulnerabilidad}
\label{sec:cristalizacion-vulnerabilidad}
% -------------------------------------------------------------

La primera pregunta de síntesis es si los dos ejes del proyecto
(compresión e interpretabilidad) convergen cuantitativamente, no
solo metafóricamente. El Capítulo~\ref{sec:resultados-interpretabilidad}
mostró que cada emoción cristaliza a una profundidad distinta:
\emph{gratitude} alcanza el 80\% de su F1 máximo en la capa~0, mientras que
\emph{disappointment} no lo hace hasta la capa~9. El
Capítulo~\ref{sec:resultados-compresion} mostró que la compresión SVD
uniforme afecta a cada emoción de forma diferencial.

\subsubsection{Correlación cristalización-vulnerabilidad}

La correlación entre la capa de cristalización y la caída de F1 bajo
compresión SVD uniforme a rango~128 es positiva y significativa (Spearman
$\rho \approx 0{,}48$, $p < 0{,}05$, $n=23$): las emociones que cristalizan
en capas más profundas son sistemáticamente más vulnerables a la
compresión. Este resultado es coherente con la estructura espectral
documentada en \S\ref{sec:analisis-espectral}: las matrices de las capas
tardías presentan distribuciones de valores singulares menos concentradas
(menor redundancia), lo que las hace más frágiles bajo truncamiento SVD.

La conexión, sin embargo, no se agota en una correlación. Tres predictores
independientes de la vulnerabilidad emergen del análisis combinado:

\begin{itemize}[nosep]
    \item Profundidad de cristalización por probing (\S\ref{sec:probing}):
    $\rho \approx 0{,}48$.
    \item Aislamiento representacional, distancia media a otras emociones
    en el espacio de \texttt{[CLS]} de L11 (\S\ref{sec:probing}):
    $\rho = 0{,}45$, $p = 0{,}031$.
    \item Norma del vector de selectividad neuronal
    (\S\ref{sec:neuron-analysis}): $\rho = 0{,}64$, $p = 0{,}001$.
\end{itemize}

El predictor más fuerte es la \emph{intensidad de codificación neuronal}
, no la geometría SVD ni el tamaño del dataset, ambos no significativos.
Este resultado es teóricamente importante: la SVD no ataca selectivamente
las neuronas emocionales (la fracción de señal en el subespacio eliminado
es prácticamente idéntica para las 23 emociones, rango 0{,}53--0{,}56),
pero las emociones que el modelo escribe con mayor intensidad neuronal son
más sensibles a perturbaciones generales de los pesos. Es una
vulnerabilidad por intensidad, no por geometría.

La Figura~\ref{fig:cap6-emotional-landscape} sintetiza visualmente esta
estructura cruzando los dos predictores principales: la capa de
cristalización (eje horizontal) y la norma de selectividad neuronal
(eje vertical).
 
\begin{figure}[H]
    \centering
    \includegraphics[width=0.95\textwidth]{figures/cap6/cap6_emotional_landscape_es.png}
    \caption{Mapa emocional del modelo. Cada punto representa una de las
    23 emociones, situada según su capa de cristalización por probing
    (eje horizontal, derivado del Capítulo~\ref{sec:resultados-interpretabilidad})
    y su norma de selectividad neuronal (eje vertical, intensidad total
    con la que el modelo codifica esa emoción a través de sus neuronas
    FFN). Los colores agrupan las emociones por cluster semántico. Las
    bandas verticales sombreadas marcan los tres bloques de profundidad:
    temprano (Emb--L3), medio (L4--L7) y tardío (L8--L11). Las emociones
    léxicas (\emph{gratitude}, \emph{love}, \emph{amusement}) ocupan la
    región superior izquierda; las contextuales (\emph{annoyance},
    \emph{disappointment}) la inferior derecha.}
    \label{fig:cap6-emotional-landscape}
\end{figure}
 
La distribución no es uniforme. Las emociones léxicas se concentran en
la región superior izquierda del mapa: cristalización temprana y alta
intensidad neuronal. \emph{Gratitude} es el caso extremo, con
cristalización en L0 y la mayor intensidad de codificación neuronal del
dataset. Las emociones contextuales ocupan la región inferior derecha:
cristalización tardía y baja intensidad. \emph{Disappointment} y
\emph{annoyance} muestran una intensidad de codificación
aproximadamente un orden de magnitud inferior a pesar de cristalizar
en L9--L10. Los valores absolutos del eje vertical de la
Figura~\ref{fig:cap6-emotional-landscape} corresponden a la \emph{máxima}
norma capa a capa del vector de selectividad (resumen visual), mientras
que la norma agregada $\ell_2$ a través de las 12~capas, reportada en
\S\ref{sec:neuron-analysis} y utilizada como predictor cuantitativo
($\rho = 0{,}639$), produce la misma ordenación entre emociones pero
en una escala distinta (\emph{gratitude}: 140; \emph{remorse}: 111;
\emph{disappointment}: 42).
 
Esta correspondencia entre dos magnitudes derivadas de experimentos
independientes ---probing por un lado, análisis de selectividad por
otro--- no es trivial. Probing mide la separabilidad lineal de la
representación; selectividad mide la especialización de las neuronas
individuales. Que ambas correlacionen sugiere que la facilidad de
aprendizaje y la intensidad de codificación son manifestaciones de una
misma propiedad subyacente: la disponibilidad de señal léxica robusta.
Las emociones contextuales no son simplemente más difíciles ---son
representacionalmente más débiles---, y esa debilidad estructural es
la que las hace particularmente vulnerables a la compresión.

\subsubsection{Tres principios de compresión informada}

De estos hallazgos se derivan tres principios que orientan el diseño de
estrategias de compresión informadas:

\begin{enumerate}
    \item Proteger las capas tardías. Concentran la
    especialización funcional (84\% de neuronas significativas, 77\% de
    cabezas críticas) y tienen menor redundancia espectral. Comprimir
    capas 8--11 al mismo rango que capas 0--3 es estructuralmente
    incorrecto.

    \item Distinguir componentes dentro de cada capa. La FFN
    Intermediate es el cuello de botella computacional confirmado por
    \emph{activation patching} (\S\ref{sec:activation-patching}); Q y K toleran
    compresión agresiva con coste despreciable. Aplicar el mismo rango a
    Q y a FFN Intermediate ignora un factor 14$\times$ de diferencia en
    retención de F1 a rango~128 (que se amplía a 72$\times$ a rango~64
    normalizado por parámetros eliminados).

    \item Aceptar que la compresión introduce sesgo emocional.
    Las emociones contextualmente complejas (cristalización tardía,
    codificación intensa) son penalizadas desproporcionadamente. Este
    sesgo es \emph{inevitable} sin intervención selectiva: cualquier
    estrategia uniforme lo introduce de forma silenciosa.
\end{enumerate}

Los principios son cualitativos. La pregunta abierta, y que motiva el
resto del capítulo, es si pueden traducirse directamente en reglas
heurísticas o si requieren cuantificación empírica.


% -------------------------------------------------------------
\subsection{Estrategia heurística: un resultado negativo informativo}
\label{sec:informed-heuristic}
% -------------------------------------------------------------

El primer intento de traducir la comprensión mecánica en una estrategia de
compresión adoptó un enfoque heurístico: asignar rangos SVD a cada
componente y capa según reglas cualitativas derivadas de los hallazgos de
interpretabilidad. La premisa de partida era razonable: si comprender el
modelo permite comprimirlo mejor, las reglas explícitas derivadas de esa
comprensión deberían superar a las estrategias ciegas.

\subsubsection{Diseño de las tres variantes}

Se diseñaron tres niveles de agresividad respetando los principios de
\S\ref{sec:cristalizacion-vulnerabilidad}:

\begin{itemize}[nosep]
    \item \textbf{Informed-light}: compresión conservadora que protege las
    capas tardías con rangos altos (700 en capa~11) y aplica reducciones
    moderadas en capas tempranas. Ratio de parámetros: 0{,}97.

    \item \textbf{Informed-moderate}: compresión intermedia con protección
    de la capa~11 (rango 512) y rangos entre 170 y 430 en el resto. Ratio:
    0{,}72.

    \item \textbf{Informed-aggressive}: compresión fuerte con rangos entre
    86 y 384. Ratio: 0{,}47.
\end{itemize}

Las tres variantes respetaban estrictamente el gradiente de criticidad por
profundidad (más rango en capas tardías que en tempranas) y la asimetría
entre componentes (más rango en FFN que en atención). La intuición
mecanicista estaba codificada explícitamente en las reglas.

\subsubsection{Resultados}

Los resultados, sin embargo, son reveladores
(Tabla~\ref{tab:informed-heuristic}):

\begin{table}[ht]
\centering
\caption{Rendimiento de las estrategias heurísticas informadas frente a
baselines de complejidad comparable. Las tres variantes \emph{igualan
exactamente} al baseline uniforme con el mismo ratio de parámetros
(\texttt{informed\_aggressive} = \texttt{uniform\_r256};
\texttt{informed\_moderate} = \texttt{uniform\_r512}), o requieren más
parámetros que el original sin alcanzar la retención de
\texttt{adaptive\_e99} (\texttt{informed\_light}). Ninguna domina la
frontera de Pareto.}
\label{tab:informed-heuristic}
\small
\begin{tabular}{llccc}
\toprule
\textbf{Estrategia} & \textbf{Tipo} & \textbf{Ratio params.} & \textbf{F1 macro} & \textbf{Retención} \\
\midrule
uniform\_r256        & uniforme   & 0{,}612 & 0{,}025 & 4{,}3\% \\
informed\_aggressive & informada  & 0{,}612 & 0{,}025 & 4{,}3\% \\
\midrule
adaptive\_e80        & adaptativa & 0{,}721 & 0{,}143 & 24{,}8\% \\
\midrule
uniform\_r512        & uniforme   & 1{,}000 & 0{,}464 & 80{,}4\% \\
informed\_moderate   & informada  & 1{,}000 & 0{,}464 & 80{,}4\% \\
\midrule
adaptive\_e95        & adaptativa & 1{,}022 & 0{,}491 & 85{,}0\% \\
adaptive\_e99        & adaptativa & 1{,}202 & 0{,}558 & 96{,}7\% \\
informed\_light      & informada  & 1{,}285 & 0{,}561 & 97{,}2\% \\
\bottomrule
\end{tabular}
\end{table}

Las tres variantes presentan el mismo patrón estructural: \emph{igualan
exactamente} al baseline uniforme de complejidad equivalente, o lo superan
solo a costa de un incremento sustancial de parámetros. La variante
agresiva, con ratio 0{,}612 y F1 = 0{,}025, reproduce número a número los
resultados de \texttt{uniform\_r256}: misma asignación efectiva de rangos
agregada, mismo colapso del modelo. La moderada, con ratio 1{,}000 y
F1 = 0{,}464, iguala a \texttt{uniform\_r512} con idéntica precisión.
La variante conservadora alcanza un F1 de 0{,}561 ---ligeramente
superior al baseline \texttt{adaptive\_e99}--- pero requiere un 28{,}5\%
\emph{más} parámetros que el modelo original (ratio 1{,}285), lo que la
sitúa fuera del régimen de compresión efectiva.

Ninguna de las tres variantes informadas aparece en la frontera de Pareto:
todas son \emph{dominadas} por las baselines de complejidad equivalente o
por las estrategias adaptativas energéticas.

\subsubsection{Diagnóstico del fracaso}

Este resultado negativo es \emph{instructivo}, no anecdótico. La intuición
cualitativa derivada de la interpretabilidad (``proteger capas tardías,
comprimir agresivamente Q y K, respetar la asimetría componente'') es
correcta en su dirección ---el algoritmo greedy
(\S\ref{sec:greedy-algorithm}) confirma exactamente esa secuencia de
decisiones cuando opera con datos cuantitativos---, pero traducirla a
reglas heurísticas explícitas no produce ganancia alguna sobre las
baselines: las informed simplemente reproducen la misma frontera que las
estrategias ciegas alcanzan sin invocar interpretabilidad.

La razón es la \emph{magnitud relativa} de las sensibilidades. Saber que
Q es más compresible que FFN Intermediate no indica cuánto más, ni en
qué profundidad la asimetría es mayor, ni cuál es el orden exacto de
prioridad entre Q en capas tardías vs.\ FFN Output en capas tempranas. La
regla heurística asigna, por ejemplo, ``rango 256 a Q en capas medias'';
los datos cuantitativos revelan que Q a rango 128 retiene 99{,}4\% del
F1 (\S\ref{sec:sensibilidad-componente}), por lo que el rango 256 es
innecesariamente conservador. Multiplicado por las 72 matrices del modelo,
el conservadurismo acumulado lleva a la heurística a converger
exactamente sobre los puntos que la compresión uniforme alcanza sin
ninguna información estructural: \texttt{informed\_aggressive} reproduce
\texttt{uniform\_r256}, \texttt{informed\_moderate} reproduce
\texttt{uniform\_r512}. La heurística \emph{no extrae valor adicional}
del conocimiento mecanicista que pretende traducir.

La conclusión metodológica es precisa: \emph{traducir interpretabilidad
mecánica en compresión efectiva requiere datos empíricos de sensibilidad,
no solo comprensión arquitectónica}. La interpretabilidad cualitativa
identifica qué investigar; la sensibilidad cuantitativa determina cuánto
comprimir. Que las tres variantes informadas converjan sobre la frontera
de las baselines ciegas, en lugar de superarla, es la evidencia más
limpia posible de esta separación: el conocimiento cualitativo de la
estructura funcional es necesario pero no suficiente para guiar la
compresión. Esta distinción motiva el cambio de enfoque hacia el
algoritmo greedy data-driven.


% -------------------------------------------------------------
\subsection{Algoritmo greedy: compresión guiada por datos}
\label{sec:greedy-algorithm}
% -------------------------------------------------------------

El fracaso de la estrategia heurística motivó un cambio de enfoque: en
lugar de traducir la intuición mecánica en reglas, utilizar directamente
los datos cuantitativos de sensibilidad por componente
(\S\ref{sec:sensibilidad-componente}) y por profundidad como entrada para
un algoritmo de optimización.

\subsubsection{Por qué greedy y no otra estrategia}

El problema es seleccionar simultáneamente 18~rangos (6~componentes
$\times$ 3~bandas de profundidad) en un menú discreto de 3~valores
candidatos cada uno, optimizando el F1 macro bajo una restricción de
parámetros máximos. Tres aproximaciones aparecen como candidatas
naturales y todas presentan inconvenientes serios.

La búsqueda exhaustiva es inviable: $3^{18} \approx 3{,}9\cdot 10^8$
configuraciones, $\sim$30~s por evaluación en GPU T4, dan unos
370~años de cómputo. La optimización bayesiana sobre el espacio de
rangos discreto es viable conceptualmente, pero exige cientos de
evaluaciones del modelo comprimido para construir el \emph{surrogate}
(varias horas de GPU) y el espacio discreto degrada la calidad del GP.
Los métodos diferenciables (\textit{straight-through estimators},
Gumbel-softmax sobre rangos) requieren relajaciones que cambian la
naturaleza del problema y no aprovechan los datos de sensibilidad ya
medidos.

El greedy data-driven evita estas tres limitaciones. Es interpretable
---la secuencia de movimientos se puede leer y se corresponde con una
narrativa funcional que valida cruzadamente los hallazgos del
Capítulo~\ref{sec:resultados-interpretabilidad}---, lineal en el
número de movimientos ($O(|moves|)$) y resuelve el problema en menos
de un segundo, sin requerir ninguna evaluación adicional del modelo
porque los costes provienen de los datos pre-calculados de
sensibilidad. Como contrapartida es sub-óptimo: el algoritmo no
explora el espacio combinatorio completo y no garantiza alcanzar el
máximo global. Su validación es empírica: 8~de los 9~puntos
Pareto-óptimos sobre 21~configuraciones evaluadas son greedy
(\S\ref{sec:greedy-algorithm}), lo que sugiere que la sub-optimalidad
real es modesta dentro de la región funcional del espacio
compresión--rendimiento.

\subsubsection{Diseño del algoritmo}

El algoritmo greedy opera sobre un menú finito de \emph{movimientos de
compresión}, donde cada movimiento consiste en reducir el rango SVD de un
componente (Q, K, V, Attn Output, FFN Intermediate, FFN Output) en una
banda de profundidad (Early: capas 0--3, Middle: 4--7, Late: 8--11) desde
su rango actual a un rango inferior. Los rangos candidatos no son
uniformes para todos los componentes: se restringen a aquellos donde la
pérdida de F1 medida en \S\ref{sec:sensibilidad-componente} es asumible.
En concreto, Q y K admiten los tres rangos $\{256, 128, 64\}$; V, Attn
Output y FFN Output sólo $\{256, 128\}$; FFN Intermediate sólo $\{256\}$,
porque a rango~128 ya colapsa al 6{,}9\% de retención y a rango~64 al
0\%. Esta restricción es un \emph{filtro previo} derivado de los
hallazgos de sensibilidad y reduce el menú a 14~movimientos por banda
($\times$\,3 bandas = 42 movimientos posibles). Para llegar al rango más
agresivo el algoritmo debe pasar por los intermedios: los movimientos son
incrementales.

Cada movimiento tiene dos magnitudes asociadas:

\begin{itemize}
    \item \textbf{Coste incremental} ($\Delta$F1): la caída de F1 macro
    al bajar del rango previo al actual (no respecto al baseline), de
    modo que el algoritmo trata la compresión como una secuencia de
    decisiones independientes. El coste se modula por la banda de
    profundidad según las proporciones empíricas de impacto observadas
    en \S\ref{sec:sensibilidad-componente}: las capas tempranas absorben
    el 14{,}8\% del coste total, las medias el 17{,}7\% y las tardías el
    67{,}5\%. Como filtro adicional, se descartan los movimientos cuyo
    coste incremental excede 15~puntos porcentuales (\textit{cap} de
    seguridad \texttt{max\_cost\_per\_move=15.0} en la implementación),
    lo que evita decisiones catastróficas como comprimir FFN Intermediate
    en capas tardías.

    \item \textbf{Ahorro}: la reducción en número de parámetros al bajar el
    rango. Un detalle técnico crítico es el \emph{rango de equilibrio}
    (\textit{break-even rank}): la descomposición SVD $W \approx U \Sigma V^T$
    solo reduce parámetros cuando $r < mn/(m+n)$. Para las matrices de
    atención ($768 \times 768$), el equilibrio está en rango~384; para las
    matrices FFN ($768 \times 3072$), en rango~614{,}4. Rangos superiores a
    estos umbrales \emph{incrementan} el número de parámetros, lo que el
    algoritmo evita por construcción.
\end{itemize}

La eficiencia de cada movimiento se define como el cociente entre el
ahorro de parámetros y el coste incremental en F1:
\begin{equation}
    \mathcal{E}(\text{movimiento}) =
    \begin{cases}
        +\infty & \text{si } \Delta F1 \leq 0 \\[4pt]
        \dfrac{\text{params\_ahorrados}}{\Delta F1 \cdot 100} & \text{en otro caso}
    \end{cases}
    \label{eq:eficiencia}
\end{equation}
El algoritmo selecciona movimientos en orden decreciente de eficiencia
hasta alcanzar un ratio de parámetros objetivo. Los movimientos con coste
incremental nulo o negativo (componentes que mejoran ligeramente al
comprimirse, lo que ocurre, por ejemplo, con Q a rango~256 con retención
del 100{,}4\%) tienen eficiencia infinita y se aplican primero. La
formalización completa se recoge en el Algoritmo~\ref{alg:greedy}.

\begin{algorithm}[H]
\caption{Greedy SVD Compression (informed by sensitivity)}
\label{alg:greedy}
\begin{algorithmic}[1]
\Require Modelo $M$, ratio objetivo de parámetros $r_{\text{tgt}}$,
         tabla de retención $S[c][k]$, fracciones de profundidad $f[d]$,
         cap de coste $\tau$
\Ensure Asignación de rangos $\sigma:(c,d)\!\to\!\text{rank}\cup\{\text{full}\}$
\State $\text{moves}\gets\emptyset$
\For{cada componente $c$}
  \For{cada par consecutivo $(k_{\text{prev}},k)$ en $S[c]$}
    \For{cada banda de profundidad $d\in\{\text{early},\text{mid},\text{late}\}$}
      \State $\Delta F1 \gets \big(S[c][k_{\text{prev}}]-S[c][k]\big)\cdot f[d]$
      \State $\text{saved}\gets |d|\cdot\big(k_{\text{prev}}(m{+}n)-k(m{+}n)\big)$
      \State $\text{eff} \gets \begin{cases}+\infty & \Delta F1\le 0\\ \text{saved}/\Delta F1 & \text{otro caso}\end{cases}$
      \State $\text{moves}.\text{add}(\langle c,d,k_{\text{prev}}{\to}k,\Delta F1,\text{saved},\text{eff}\rangle)$
    \EndFor
  \EndFor
\EndFor
\State ordenar $\text{moves}$ por $\text{eff}$ descendente
\State $\sigma[c,d]\gets\text{full}$ para todo $(c,d)$
\For{$m\in\text{moves}$}
  \If{$\sigma[m.c,m.d]\ne m.k_{\text{prev}}$} \textbf{continue} \EndIf
  \If{$m.\Delta F1>\tau$} \textbf{continue} \EndIf
  \State $\sigma[m.c,m.d]\gets m.k$
  \If{$\text{params\_ratio}(\sigma)\le r_{\text{tgt}}$} \textbf{break} \EndIf
\EndFor
\State \Return $\sigma$
\end{algorithmic}
\end{algorithm}


\subsubsection{Secuencia de decisiones}

La secuencia con que el algoritmo asigna compresión revela una jerarquía
de importancia que valida \emph{independientemente} los hallazgos de
interpretabilidad (Tabla~\ref{tab:greedy-sequence}):

\begin{table}[ht]
\centering
\caption{Secuencia de decisiones del algoritmo greedy. Cada fila muestra
los componentes comprimidos al alcanzar el ratio indicado. El algoritmo
opera exclusivamente con datos numéricos de sensibilidad, sin acceso a
los hallazgos cualitativos de interpretabilidad.}
\label{tab:greedy-sequence}
\small
\begin{tabular}{ccp{8.0cm}}
\toprule
\textbf{Config.} & \textbf{Ratio} & \textbf{Componentes comprimidos} \\
\midrule
greedy\_95 & 0{,}950 & Solo Q y K \\
greedy\_90 & 0{,}864 & Q, K, FFN Output (capas tempranas) \\
greedy\_85 & 0{,}813 & Q, K, FFN Output (tempranas + medias) \\
greedy\_80 & 0{,}799 & Q, K, V (tempranas), FFN Output \\
greedy\_75 & 0{,}741 & Q, K, V, FFN Inter+Out (tempranas) \\
greedy\_70 & 0{,}670 & Todos los componentes, mayoritariamente
                       tempranos y medios \\
\bottomrule
\end{tabular}
\end{table}

Tres observaciones emergen de esta secuencia, cada una con su contraparte
en interpretabilidad:

\begin{enumerate}
    \item Q y K son esencialmente gratuitas de comprimir. El
    primer 5\% de reducción de parámetros proviene íntegramente de
    comprimir las proyecciones de Query y Key, con un coste en F1 inferior
    al 0{,}5\%. Esto es coherente con el análisis de sensibilidad por
    componente (\S\ref{sec:sensibilidad-componente}), que documentó
    retenciones del 99{,}4\% para Q y 98{,}4\% para K a rango~128, y con
    el solapamiento espectral elevado entre subespacios Q-K
    (hasta 0{,}63 para los 32 vectores singulares principales,
    \S\ref{sec:analisis-espectral}). Q y K codifican esencialmente
    redundancia, no información distintiva.

    \item La FFN Intermediate no se comprime en capas tardías.
    Ninguna configuración greedy con retención superior al 30\% toca la
    FFN Intermediate de las capas 8--11. El algoritmo descubre de forma
    independiente el cuello de botella identificado por activation
    patching (\S\ref{sec:activation-patching}): restaurar únicamente la
    FFN de la capa~11 recuperaba el 100\% del rendimiento del modelo. La
    convergencia entre dos métodos completamente disjuntos (uno mide
    sensibilidad puntual, otro mide suficiencia causal) en el mismo
    componente como \emph{único} crítico es una validación cruzada fuerte.

    \item Las capas tardías se protegen sistemáticamente. Incluso
    en las configuraciones más agresivas (\texttt{greedy\_70}, ratio
    0{,}67), las capas 8--11 retienen rangos sustancialmente mayores que
    las tempranas. Esto valida el gradiente de criticidad documentado por
    la ablación de cabezas (\S\ref{sec:head-ablation}): el 77\% de las
    cabezas en capas 8--11 son críticas frente al 27\% en capas 0--4.
\end{enumerate}

La convergencia entre la secuencia de decisiones del algoritmo (que opera
solo con datos numéricos de sensibilidad por componente y profundidad) y
los hallazgos cualitativos de interpretabilidad (que explican \emph{por
qué} esos componentes son sensibles) constituye una validación cruzada de
ambas líneas de análisis. Si interpretabilidad y compresión llegan a las
mismas conclusiones por caminos disjuntos, la probabilidad de que ambas
estén sistemáticamente sesgadas en la misma dirección es baja.


\subsubsection{Dominancia en la frontera de Pareto}

La Figura~\ref{fig:pareto-all-strategies} presenta la comparación completa
de las 21 configuraciones evaluadas (6 uniformes, 4 adaptativas, 3
heurísticas informadas y 8 greedy) en el espacio ratio de parámetros
vs.\ F1 macro.

\begin{figure}[ht]
\centering
\includegraphics[width=0.95\textwidth]{figures/cap6/cap6_pareto_all_strategies_es.png}
\caption{Frontera de Pareto para las cuatro familias de estrategias de
compresión. Los puntos sobre la curva envolvente representan
configuraciones Pareto-óptimas. 8 de los 9 puntos Pareto-óptimos
corresponden a estrategias greedy. La única estrategia no-greedy en la
frontera es \texttt{uniform\_r32}, que aparece solo por tener el menor
ratio absoluto pero con F1 = 0 (el modelo está colapsado).}
\label{fig:pareto-all-strategies}
\end{figure}

De los 9 puntos Pareto-óptimos en sentido estricto, 8 son estrategias
greedy. El algoritmo basado en datos empíricos domina la frontera
completa en la región funcional ($F1 > 0$), desde un ratio de
0{,}58 hasta 0{,}95. La Tabla~\ref{tab:head-to-head} cuantifica las
ventajas en comparaciones directas a ratios de parámetros similares:

\begin{table}[ht]
\centering
\caption{Comparaciones directas entre estrategias greedy y alternativas a
ratios de parámetros comparables. Los intervalos de confianza al 95\% se
obtienen mediante \emph{bootstrap por emoción} (10\,000 réplicas,
$n=23$); el \emph{p}-valor en la columna $\Delta$F1 corresponde a un
test \emph{paired bootstrap} sobre la diferencia por emoción cuando se
dispone de datos por emoción para ambas estrategias.}
\label{tab:head-to-head}
\small
\begin{tabular}{lcccc}
\toprule
\textbf{Comparación a ratio similar} & \textbf{Greedy F1 [CI95]} & \textbf{Alt.\ F1 [CI95]} & $\Delta$\textbf{F1} & \textbf{Sig.} \\
\midrule
greedy\_80 ($\rho=0{,}80$) vs uniform\_r384 ($\rho=0{,}81$)
  & 0{,}502\textsuperscript{$\dagger$}
  & 0{,}251 [0{,}151; 0{,}365]
  & $+$0{,}251 & --- \\
greedy\_90 ($\rho=0{,}86$) vs adaptive\_e90 ($\rho=0{,}89$)
  & 0{,}539 [0{,}467; 0{,}613]
  & 0{,}357 [0{,}255; 0{,}466]
  & $+$0{,}181 [$+$0{,}130; $+$0{,}241]
  & $p<0{,}001$ \\
greedy\_95 ($\rho=0{,}95$) vs adaptive\_e95 ($\rho=1{,}02$)
  & 0{,}575\textsuperscript{$\dagger$}
  & 0{,}491 [0{,}400; 0{,}581]
  & $+$0{,}084 & --- \\
\bottomrule
\multicolumn{5}{@{}l}{\footnotesize \textsuperscript{$\dagger$}Se reporta
únicamente el F1 macro puntual: los datos por emoción para
\texttt{greedy\_80} y \texttt{greedy\_95} no se exportaron a CSV en NB09. El
test \emph{paired bootstrap} para greedy\_90 vs adaptive\_e90 produce un
$\Delta$F1 = $+$0{,}181 con $p<0{,}001$ y CI estrictamente positivo, lo}\\
\multicolumn{5}{@{}l}{\footnotesize que es la evidencia estadística más
limpia disponible. Para los otros dos comparativos, los puntos del greedy
caen \emph{fuera} del intervalo de confianza al 95\% de la alternativa, lo
que es consistente con significancia direccional.}
\end{tabular}
\end{table}

El resultado más llamativo se da en el rango intermedio: a un ratio de
$\sim$0{,}80, la estrategia greedy retiene el 87{,}0\% del F1 original,
mientras que la uniforme a rango~384 solo retiene el 43{,}5\%: una
ventaja superior al doble en retención. La comparación
\texttt{greedy\_90} vs \texttt{adaptive\_e90}, donde sí se dispone de
datos por emoción para ambas, permite cuantificar formalmente la ventaja:
$\Delta$F1 = $+$0{,}181 con intervalo de confianza al 95\%
$[+0{,}130; +0{,}241]$ (\emph{paired bootstrap} sobre las 23 emociones,
10\,000 réplicas, $p<0{,}001$). El intervalo está estrictamente alejado
de cero, descartando la posibilidad de que la ventaja del greedy sea
atribuible a variabilidad muestral entre emociones. Nótese que la
comparación con \texttt{adaptive\_e95} es particularmente reveladora: el
algoritmo greedy obtiene un F1 superior con un 7\% \emph{menos} de
parámetros, dado que la estrategia adaptativa al 95\% de energía produce
un ratio de 1{,}02 (es decir, \emph{incrementa} los parámetros, como se
discutió en \S\ref{sec:compresion-adaptativa}).

A ratio 0{,}95, la estrategia \texttt{greedy\_95} alcanza una retención
del 99{,}6\% del F1 baseline (sin pérdida medible) con
una reducción del 5\% en parámetros que viene íntegramente de comprimir Q
y K, lo que ilustra que existe un margen de compresión genuinamente
gratuito que las estrategias ciegas son incapaces de explotar.

\subsubsection{Comparación cualitativa con destilación}
\label{sec:greedy-vs-distilbert}

Las cifras del greedy se entienden mejor situándolas frente a la
familia ortogonal de compresión más popular: la destilación de
conocimiento.
DistilBERT~\cite{sanh2019distilbert}, la destilación canónica de
BERT-base, reduce el modelo a 6~capas (66~M parámetros, ratio
0{,}60 sobre los 110~M de BERT-base) reteniendo, según los autores,
el 97\% del rendimiento de BERT-base en GLUE. Esta comparación no es
estrictamente trasladable a GoEmotions multi-etiqueta sin re-evaluar
DistilBERT en la misma tarea ---ejercicio que se propone como
trabajo futuro (\S\ref{sec:trabajo-futuro})---, pero permite tres
observaciones de calibración.

A nivel de compresión agresiva ($\rho \approx 0{,}60$), DistilBERT
domina cómodamente: nuestra \texttt{greedy\_60} retiene apenas el
5{,}5\% del F1, mientras que la destilación re-entrena el modelo
pequeño desde cero con supervisión de un \emph{teacher} y puede
reorganizar libremente sus pesos. La SVD opera sobre los pesos
existentes sin re-entrenamiento, una restricción que pasa factura en
el régimen de alta reducción.

En la zona de compresión moderada ($\rho \in [0{,}80, 0{,}90]$) la
situación se invierte. La SVD greedy es competitiva o superior a las
técnicas ciegas convencionales sin necesidad de re-entrenar. Y si se
admite \emph{fine-tuning} posterior, la combinación greedy + FT
(\S\ref{sec:finetuning-recovery}) supera incluso al baseline sin
comprimir, algo que ninguna técnica ciega consigue a este ratio.

Las dos familias son complementarias más que competidoras. La
destilación es la opción correcta cuando se tolera el coste de
pre-procesamiento y la latencia de desarrollo larga; la SVD informada
lo es cuando se requiere intervenir directamente sobre un modelo ya
desplegado, o cuando se buscan reducciones moderadas
($<25\%$ params) sin tocar el ciclo de entrenamiento. La combinación
de ambas ---destilación seguida de SVD sobre el destilado--- está
poco explorada y constituye un experimento natural derivado de este
trabajo.

La elección del SVD selectivo en este TFG no pretende superar a la
destilación en su propio terreno, sino habilitar una zona del espacio
compresión--rendimiento que la destilación deja sin cubrir: ratios
0{,}80--0{,}95 con intervención post-hoc sobre el modelo
ya \emph{fine-tuneado}.


% -------------------------------------------------------------
\subsection{Recuperación mediante \emph{fine-tuning}}
\label{sec:finetuning-recovery}
% -------------------------------------------------------------

La compresión SVD, incluso con asignación de rangos óptima, introduce una
perturbación en los pesos del modelo. Una pregunta natural es si un breve
re-entrenamiento puede recuperar, o incluso superar, el rendimiento
original.

Se seleccionó la configuración \texttt{greedy\_90} (86{,}4\% de los
parámetros, F1 macro = 0{,}539) como candidata para \emph{fine-tuning}, por
ofrecer un equilibrio atractivo entre compresión y rendimiento. El
re-entrenamiento consistió en 3~épocas sobre GoEmotions con los
hiperparámetros estándar (tasa de aprendizaje $2 \times 10^{-5}$,
\textit{warmup} 10\%, \textit{batch size} 32, idénticos a los del
\emph{fine-tuning} original).

\subsubsection{Resultado agregado}

\begin{table}[ht]
\centering
\caption{Rendimiento del modelo comprimido (\texttt{greedy\_90}) antes y
después de \emph{fine-tuning}, comparado con el baseline sin comprimir.}
\label{tab:finetuning-summary}
\small
\begin{tabular}{lccc}
\toprule
\textbf{Modelo} & \textbf{Parámetros} & \textbf{F1 macro} & \textbf{vs.\ baseline} \\
\midrule
Baseline (sin compresión) & 100{,}0\% & 0{,}577 & --- \\
greedy\_90 (comprimido)   & 86{,}4\%  & 0{,}539 & $-$6{,}7\% \\
greedy\_90 (\emph{fine-tuned})   & 86{,}4\%  & \textbf{0{,}591} & \textbf{$+$2{,}4\%} \\
\bottomrule
\end{tabular}
\end{table}

El modelo comprimido y re-entrenado supera al baseline original con un
13{,}6\% menos de parámetros. Tres épocas de \emph{fine-tuning} no solo recuperan
la caída de 6{,}7 puntos porcentuales, sino que añaden 1{,}4 puntos
adicionales (todos los porcentajes calculados sobre F1 macro = 0{,}577 del
baseline). Este resultado, observado consistentemente al repetir el
experimento con distintas semillas aleatorias, sugiere que la compresión
SVD, al eliminar componentes de bajo valor singular, actúa durante el
re-entrenamiento como una forma de regularización implícita.

Esta interpretación tiene una limitación metodológica: el experimento no incluye un grupo de control consistente
en el baseline original con tres épocas adicionales de \emph{fine-tuning}. Es
posible, aunque improbable dada la saturación de F1 macro observada
durante las 4 épocas iniciales, que parte de la mejora provenga
simplemente de más entrenamiento y no del efecto regularizador de la
compresión. El experimento de control que verificaría la hipótesis se propone en el Capítulo~\ref{sec:conclusiones}.


\subsubsection{Análisis por emoción}

El patrón de recuperación no es uniforme: las emociones con F1 baseline
bajo son las que más se benefician de la compresión seguida de
\emph{fine-tuning}, y este patrón es sistemático
(Tabla~\ref{tab:finetuning-emotions}).

\begin{table}[ht]
\centering
\caption{Emociones con mayor ganancia y mayor pérdida tras compresión +
\emph{fine-tuning}, respecto al baseline original. La columna $\Delta$ es la
diferencia entre F1 \emph{fine-tuned} y F1 baseline.}
\label{tab:finetuning-emotions}
\small
\begin{tabular}{lcccc}
\toprule
\textbf{Emoción} & \textbf{F1 base} & \textbf{F1 comprim.} & \textbf{F1 \emph{fine-tuned}} & \textbf{$\Delta$ vs.\ base} \\
\midrule
\multicolumn{5}{l}{\emph{Mayores ganancias}} \\
embarrassment & 0{,}267 & 0{,}273 & 0{,}509 & \textbf{$+$0{,}242} \\
desire        & 0{,}535 & 0{,}554 & 0{,}621 & $+$0{,}086 \\
excitement    & 0{,}454 & 0{,}431 & 0{,}497 & $+$0{,}043 \\
realization   & 0{,}271 & 0{,}195 & 0{,}298 & $+$0{,}027 \\
disappointment & 0{,}275 & 0{,}178 & 0{,}301 & $+$0{,}026 \\
\midrule
\multicolumn{5}{l}{\emph{Mayores pérdidas}} \\
caring        & 0{,}542 & 0{,}479 & 0{,}504 & $-$0{,}038 \\
admiration    & 0{,}757 & 0{,}695 & 0{,}727 & $-$0{,}029 \\
fear          & 0{,}710 & 0{,}680 & 0{,}691 & $-$0{,}018 \\
confusion     & 0{,}515 & 0{,}443 & 0{,}497 & $-$0{,}018 \\
sadness       & 0{,}632 & 0{,}575 & 0{,}621 & $-$0{,}011 \\
\bottomrule
\end{tabular}
\end{table}

El caso de \emph{embarrassment} es el más extremo: el modelo comprimido
mantiene una capacidad similar al baseline (F1 = 0{,}273 frente a 0{,}267),
pero el \emph{fine-tuning} la eleva hasta F1 = 0{,}509, un 90{,}6\% superior al
baseline original. Una observación parecida ocurre con \emph{desire}: el
modelo comprimido apenas pierde rendimiento (0{,}554 vs 0{,}535), pero
tras tres épocas de \emph{fine-tuning} alcanza 0{,}621, un $+16\%$ relativo.
Este patrón, en el que la compresión apenas afecta y el \emph{fine-tuning}
amplifica, encaja con la interpretación de la SVD como
regularizador estructural: al eliminar direcciones de baja energía, la
compresión proporciona al re-entrenamiento un punto de partida sin
sobreajuste a patrones espurios. Las emociones infrarrepresentadas
(\emph{embarrassment}, \emph{desire}, \emph{realization},
\emph{disappointment}) son las que más se benefician de este efecto.

Las pérdidas, por contraste, son pequeñas (máximo 0{,}038 en
\emph{caring}) y afectan a emociones que ya tenían un F1 razonablemente
alto. La compresión seguida de \emph{fine-tuning} actúa como un ecualizador:
recorta ligeramente las emociones bien representadas y amplifica las
infrarrepresentadas. Este efecto es consistente con la eliminación de
componentes de alta varianza y baja contribución que caracteriza al
truncamiento SVD: el modelo, al perder direcciones que codificaban
sobreajuste a patrones frecuentes, recupera capacidad para las clases
minoritarias.


\subsubsection{Plasticidad asimétrica: qué cambia durante el \emph{fine-tuning}}
 
Una pregunta complementaria es \emph{dónde} reorganiza el modelo sus pesos
durante el \emph{fine-tuning} de recuperación. Si el efecto regularizador
hipotetizado en el apartado anterior es real, cabría esperar que el
re-entrenamiento concentrara su esfuerzo en los componentes funcionalmente
más relevantes para la tarea. La Figura~\ref{fig:cap6-finetuning-diff}
mide esta plasticidad como la magnitud relativa del cambio de cada
matriz, $\|\Delta W\| / \|W_0\|$, donde $\|\cdot\|$ es la norma de
Frobenius y $W_0$ el peso del modelo comprimido antes del \emph{fine-tuning}.
 
\begin{figure}[H]
    \centering
    \includegraphics[width=0.85\textwidth]{figures/cap6/cap6_finetuning_diff_es.png}
    \caption{Magnitud relativa del cambio de pesos durante las tres
    épocas de \emph{fine-tuning} posterior a la compresión, calculada como
    $\|\Delta W\| / \|W_0\|$ por capa (filas) y componente (columnas).
    La FFN de la capa~11 es el componente más reorganizado (FFN-i:
    5{,}1\%, FFN-o: 5{,}4\%), seguido de la atención de la misma capa
    (Attn-O: 4{,}6\%). Las proyecciones Q y K son las menos modificadas
    en todas las capas (3{,}0\%--3{,}4\%). El patrón refleja una
    asignación selectiva del esfuerzo de re-entrenamiento.}
    \label{fig:cap6-finetuning-diff}
\end{figure}
 
El resultado proporciona una validación cruzada con los hallazgos de
sensibilidad del Capítulo~\ref{sec:resultados-compresion}: los
componentes que más se reorganizan durante el \emph{fine-tuning} son
precisamente los más sensibles a la compresión. La FFN Intermediate de
la capa~11, cuyo colapso bajo compresión a rango 128 reducía el F1 al
7\% (\S\ref{sec:sensibilidad-componente}), es también el componente con
mayor cambio durante el re-entrenamiento. Las proyecciones Q y K, que
toleraban compresión a rango 128 con retención del 99\% del F1, son
también las que menos cambian al re-entrenarse.
 
Este paralelismo sugiere que sensibilidad a la compresión y plasticidad
bajo \emph{fine-tuning} miden, desde ángulos distintos, una misma propiedad
subyacente: la importancia funcional del componente para la tarea. Los
componentes que la tarea necesita más son los que más se reorganizan al
re-entrenarse, y los que más se rompen al comprimirse. Esta convergencia
añade una tercera evidencia ---junto a la sensibilidad por componente y
la convergencia del algoritmo greedy--- de que la jerarquía funcional
documentada por interpretabilidad mecánica no es un artefacto de un
método particular, sino una propiedad estructural del modelo.


% -------------------------------------------------------------
\subsection{Discusión y síntesis de hallazgos}
\label{sec:sintesis-discusion}
% -------------------------------------------------------------

Los resultados de este capítulo permiten extraer tres conclusiones
principales que integran las líneas de compresión e interpretabilidad del
proyecto.

\subsubsection{La interpretabilidad mecánica tiene valor aplicado, pero
requiere cuantificación}

El resultado negativo de las estrategias heurísticas
(\S\ref{sec:informed-heuristic}) y el éxito del algoritmo greedy
(\S\ref{sec:greedy-algorithm}) establecen un matiz importante para la
literatura de compresión informada: la interpretabilidad mecánica no
genera valor mediante la intuición cualitativa, sino a través de la
medición sistemática de sensibilidades. Las heurísticas informed
convergen exactamente sobre la frontera de las baselines ciegas
(\texttt{informed\_aggressive} $\equiv$ \texttt{uniform\_r256},
\texttt{informed\_moderate} $\equiv$ \texttt{uniform\_r512}), evidencia
limpia de que el conocimiento cualitativo no extrae ganancia adicional
sobre lo que la compresión uniforme alcanza sin invocar interpretabilidad.
Saber que ``la FFN es más importante que la atención'' no basta; lo que
permite diseñar estrategias superiores es cuantificar que Q a rango~256
retiene el 100{,}4\% del F1 mientras que FFN Intermediate a rango~64
colapsa al 0\%.

Esto no invalida la interpretabilidad cualitativa: los hallazgos de
probing, \emph{logit lens}, \emph{activation patching} y ablación siguen siendo
esenciales para explicar el comportamiento del modelo, pero
delimita su utilidad para la ingeniería de compresión: la interpretabilidad
identifica qué investigar, y la sensibilidad empírica determina
cuánto comprimir. La distinción es metodológicamente importante
porque permite predecir cuándo otros trabajos de ``compresión informada
por interpretabilidad'' tendrán éxito o fracasarán, en función de si se
basan en reglas heurísticas o en datos cuantitativos.

\subsubsection{Convergencia de cuatro métodos independientes}

La Tabla~\ref{tab:triangulation} resume la convergencia entre los
diferentes métodos experimentales del proyecto respecto a las tres
conclusiones estructurales principales:

\begin{table}[ht]
\centering
\caption{Triangulación de hallazgos: cuatro métodos independientes
convergen en las mismas conclusiones estructurales.}
\label{tab:triangulation}
\small
\begin{tabular}{p{3.0cm}p{3.3cm}p{3.5cm}p{3.5cm}}
\toprule
\textbf{Conclusión} & \textbf{Compresión} & \textbf{Interpretabilidad} & \textbf{Síntesis} \\
\midrule
Q/K son redundantes &
  Retención $>$99\% a rango 128 (\S\ref{sec:sensibilidad-componente}) &
  Solapamiento espectral 0{,}63 con K (\S\ref{sec:analisis-espectral});
  bajo impacto en ablación (\S\ref{sec:head-ablation}) &
  Primeros en comprimirse en greedy (\S\ref{sec:greedy-algorithm}) \\
\addlinespace
FFN es el cuello de botella &
  Asimetría 14$\times$ Q vs.\ FFN Intermediate (\S\ref{sec:sensibilidad-componente}) &
  Restaurar FFN L11 = 100\% (\S\ref{sec:activation-patching});
  84\% neuronas significativas en L8--11 (\S\ref{sec:neuron-analysis}) &
  Nunca comprimida en capas tardías (\S\ref{sec:greedy-algorithm}) \\
\addlinespace
Capas tardías son críticas &
  50{,}5\% del impacto total (\S\ref{sec:sensibilidad-componente}) &
  77\% cabezas críticas en L8--11 (\S\ref{sec:head-ablation});
  cristalización tardía predice F1 (\S\ref{sec:probing});
  cristalización geométrica L9--L11 (\S\ref{sec:logit-lens}) &
  Protegidas sistemáticamente (\S\ref{sec:greedy-algorithm}) \\
\bottomrule
\end{tabular}
\end{table}

Esta convergencia es significativa porque los métodos son
\emph{independientes en su diseño experimental}: el análisis de
sensibilidad mide F1 bajo perturbación de pesos; la ablación de cabezas
mide F1 bajo anulación de componentes; el probing mide separabilidad
lineal de representaciones; el \emph{logit lens} mide alineación con la base del
clasificador; y el \emph{activation patching} mide restauración funcional. Que
los cinco coincidan en las mismas conclusiones refuerza la validez de cada
método individual y, simultáneamente, la robustez de las conclusiones
estructurales. La probabilidad de que cinco métodos disjuntos compartan
exactamente el mismo sesgo sistemático es despreciable.

\subsubsection{La compresión como técnica de regularización: una hipótesis
con apoyo preliminar}

El resultado de \emph{fine-tuning} post-compresión (un modelo con 13{,}6\%
menos de parámetros que supera al original en F1 macro) invita a
reinterpretar la compresión SVD no solo como reducción de tamaño, sino
como una técnica de regularización implícita. El truncamiento de valores
singulares bajos elimina componentes de los pesos que contribuyen poco a
la varianza explicada pero pueden codificar ruido o sobreajuste. El
re-entrenamiento posterior permite al modelo reorganizar su
representación sobre una base más compacta.

Este efecto es particularmente visible en las emociones infrarrepresentadas
en el dataset (\emph{embarrassment}, \emph{realization}, \emph{desire}),
donde el baseline original sufre de sobreajuste a patrones espurios que la
compresión elimina. La mejora de 90{,}6\% en F1 de \emph{embarrassment}
(de 0{,}267 a 0{,}509) es el ejemplo más extremo, pero el patrón se
observa consistentemente en emociones de baja frecuencia.

La interpretación causal de este patrón (compresión $\rightarrow$
regularización $\rightarrow$ mejora) requiere experimentos de control
que no se han realizado en este trabajo y que se proponen en
\S\ref{sec:trabajo-futuro}. La observación que sí se sostiene es: \emph{la compresión greedy seguida de fine-tuning
produce un modelo simultáneamente más pequeño y más preciso que el
baseline, con ganancias concentradas en clases minoritarias}. Si la
hipótesis de regularización implícita se confirma con grupos de control
adecuados, la compresión SVD podría incorporarse como paso estándar en
pipelines de entrenamiento, no solo como técnica de reducción de tamaño.

La implicación práctica es clara independientemente de la causalidad
subyacente: un pipeline de \emph{fine-tuning $\rightarrow$ compresión
greedy $\rightarrow$ re-\emph{fine-tuning}} produce modelos que son
simultáneamente más pequeños y al menos tan buenos como el baseline
original, siempre que la compresión respete la estructura funcional del
modelo ---lo cual requiere, a su vez, el tipo de análisis de
interpretabilidad mecánica presentado en esta memoria.
% =============================================================
% CAPÍTULO 7 — Conclusiones y trabajo futuro
% Versión: matrícula de honor
% =============================================================

\section{Conclusiones y trabajo futuro}
\label{sec:conclusiones}

Este capítulo recapitula las contribuciones del trabajo, discute con
honestidad sus limitaciones y propone líneas concretas de investigación
futura. Frente a la práctica habitual de cerrar memorias con una lista de
hallazgos, optamos por una estructura que distingue tres niveles: lo que el
trabajo describe sobre la organización interna de BERT, lo que el
trabajo diseña como contribución metodológica, y lo que el trabajo
abre como agenda de investigación.


% -------------------------------------------------------------
\subsection{Contribuciones principales}
\label{sec:contribuciones-finales}
% -------------------------------------------------------------

Este trabajo ha producido un mapa funcional de cómo BERT-base
\emph{fine-tuneado} organiza internamente la información emocional sobre
GoEmotions (23~emociones), utilizando la compresión SVD selectiva como
instrumento experimental y triangulando los hallazgos con cinco técnicas
de interpretabilidad mecánica. Las contribuciones se organizan en tres
bloques: la cartografía funcional propiamente dicha, el uso de la
compresión como diagnóstico estructural y la traducción del mapa en una
estrategia de compresión informada.

La Tabla~\ref{tab:objetivos-cumplimiento} resume el cumplimiento de los
12~sub-objetivos planteados en el Anexo de gestión (\S\ref{sec:gep-alcance}),
indicando para cada uno la sección de la memoria donde se desarrolla el
resultado y la evidencia cuantitativa más relevante. De los 11~objetivos
del núcleo, los 11~se han alcanzado; el O12 (recuperación por
\emph{fine-tuning}, originalmente opcional) se promovió a tarea principal
tras observar el efecto regularizador documentado en
\S\ref{sec:finetuning-recovery}.

\begin{table}[ht]
\centering
\caption{Cumplimiento de los sub-objetivos planteados en el plan inicial.
Sección indica dónde se desarrolla el resultado; evidencia, la métrica
clave que documenta el cumplimiento.}
\label{tab:objetivos-cumplimiento}
\footnotesize
\begin{tabular}{@{}p{0.45cm} p{4.6cm} p{2.9cm} p{4.6cm}@{}}
\toprule
\textbf{Obj.} & \textbf{Descripción} & \textbf{Sección} & \textbf{Evidencia} \\
\midrule
O1  & Baseline competitivo (BERT 23 emociones) & \S\ref{sec:met-modelo}     & F1 macro 0{,}577 sobre test \\
O2  & Módulo SVD reutilizable                 & \S\ref{sec:met-svd}        & \texttt{src/compression/svd.py} \\
O3  & Trade-off compresión--rendimiento       & \S\ref{sec:compresion-uniforme} & Cliff entre $r{=}384$ y $r{=}256$ \\
O4  & Visualización degradación               & \S\ref{sec:probing}        & Fig.~\ref{fig:cap5-galaxia-evolution} \\
O5  & Sensibilidad por componente             & \S\ref{sec:sensibilidad-componente} & Asimetría 14$\times$ a $r{=}128$ \\
O6  & Sensibilidad por profundidad            & \S\ref{sec:sensibilidad-componente} & Late catastrófica a $r{=}128$ \\
O7  & Compresión adaptativa por energía       & \S\ref{sec:compresion-adaptativa} & Adaptive 99\%: 96{,}7\% retención \\
O8  & Localizar cristalización emocional      & \S\ref{sec:probing}        & 276 probes; 9/8/6 emociones por banda \\
O9  & Evidencia funcional de localización     & \S\ref{sec:logit-lens}, \S\ref{sec:activation-patching} & FFN L11 sola = 100\% restauración \\
O10 & Mapear cabezas y neuronas               & \S\ref{sec:head-ablation}, \S\ref{sec:neuron-analysis} & 144 cabezas + 36\,864 neuronas \\
O11 & Compresión informada                    & \S\ref{sec:greedy-algorithm} & 8/9 puntos Pareto-óptimos greedy \\
O12 & Recuperación por \emph{fine-tuning}      & \S\ref{sec:finetuning-recovery} & F1 macro $\geq$ baseline con 86{,}4\% params (preliminar) \\
\bottomrule
\end{tabular}
\end{table}

\subsubsection{Cartografía funcional del modelo}

La aplicación coordinada de cinco técnicas de interpretabilidad (probing
lineal, \emph{logit lens}, \emph{activation patching}, ablación de cabezas y análisis de
selectividad neuronal) ha producido un mapa funcional coherente de cómo
BERT codifica y procesa la información emocional. Las observaciones se
articulan en cuatro hallazgos:

\begin{itemize}
  \item Arquitectura funcional de tres niveles. Las capas tempranas
        (Emb--L2) extraen señal léxica bruta y capturan aproximadamente el
        61\% de la separabilidad lineal final del modelo
        ($F1_{\text{probe}}^{L0} = 0{,}349$ frente al máximo
        $F1_{\text{probe}}^{L11} = 0{,}569$, \S\ref{sec:probing}). Las
        capas intermedias (L3--L7) realizan cómputo de transición en una
        geometría desalineada con la base del clasificador, lo que se
        manifiesta en el valle del \emph{logit lens} (\S\ref{sec:logit-lens}). Las
        capas profundas (L8--L11) ejecutan la cristalización: alinean la
        representación con la cabeza de clasificación entrenada y suprimen
        activamente las opciones incorrectas. Este último paso constituye el
        cuello de botella geométrico del modelo y explica su criticidad
        bajo compresión.

  \item Cristalización emocional progresiva determinada por la
        semántica, no por la frecuencia. Las 23 emociones se distribuyen en
        tres bandas de cristalización con 9, 8 y 6 emociones respectivamente.
        La frecuencia en el dataset no predice la profundidad de
        cristalización (Spearman $\rho = -0{,}26$, $p = 0{,}22$,
        \S\ref{sec:probing}), mientras que la profundidad sí predice
        fuertemente el rendimiento alcanzable
        ($\rho = -0{,}67$, $p = 0{,}0004$). Emociones con señales léxicas
        claras (\emph{gratitude}, \emph{love}, \emph{amusement}) cristalizan
        en L0; emociones que requieren desambiguación pragmática
        (\emph{annoyance}, \emph{disappointment}) cristalizan en L8--L9. Esta
        observación tiene implicaciones que van más allá de este trabajo:
        sugiere que añadir datos a las emociones contextualmente complejas no
        adelanta su cristalización.

  \item La FFN de la capa~11 como locus computacional crítico.
        Restaurar exclusivamente este componente (aproximadamente un tercio
        de los parámetros de una capa) en un modelo previamente colapsado
        por compresión recupera el 100\% del rendimiento baseline
        (\S\ref{sec:activation-patching}). Este hallazgo desafía
        productivamente la narrativa centrada en la atención que ha
        prevalecido desde \emph{Attention Is All You Need}
        \cite{vaswani2017attention}. La interpretación mecanicista, articulada
        a través del \emph{logit lens}, es que la FFN de las capas tardías realiza
        la \emph{rotación geométrica} que lleva la representación
        contextualizada a la base sobre la que opera la cabeza de
        clasificación.

  \item Gradiente monotónico de criticidad atencional. La proporción
        de cabezas críticas crece de forma estricta con la profundidad: 27\%
        en capas 0--4, 53\% en capas 5--7, 77\% en capas 8--11. La capa 11
        constituye un caso extremo: ninguna de sus 12 cabezas es prescindible
        (\S\ref{sec:head-ablation}). Identificamos 21~cabezas con
        \emph{efecto interferente} (su ablación mejora el rendimiento), que
        constituyen candidatas directas a poda sin re-entrenamiento.
\end{itemize}


\subsubsection{Compresión como instrumento de diagnóstico}

El análisis sistemático de compresión SVD, más allá de su finalidad aplicada,
ha funcionado como un instrumento de diagnóstico que revela la organización
interna del modelo de un modo complementario a las técnicas de
interpretabilidad mecánica:

\begin{itemize}
  \item Asimetría radical entre componentes. La compresión a
        rango~128 produce retenciones de F1 que difieren en un factor 14$\times$
        entre el extremo más robusto (Query: 99{,}4\%) y el más frágil (FFN
        Intermediate: 6{,}9\%). En el régimen severo de rango~64 la asimetría
        se amplía hasta el factor 72$\times$ normalizado por parámetros
        eliminados: Query mantiene una retención del 98{,}6\%, mientras que
        FFN Intermediate colapsa al 0\% (\S\ref{sec:sensibilidad-componente}). Esta asimetría es consistente
        con la estructura espectral de los pesos: las matrices Q y K presentan
        decaimiento espectral pronunciado, mientras que las FFN exhiben un
        espectro plano.

  \item Transición de fase bajo compresión uniforme. La compresión
        uniforme presenta un \emph{cliff} abrupto entre rango~384 (43{,}5\%
        de retención) y rango~256 (4{,}3\%). A rangos inferiores el modelo
        entra en colapso clínico (F1 = 0). La existencia de esta transición
        de fase, junto con la asimetría entre componentes, demuestra que la
        compresión uniforme es estructuralmente incapaz de explorar la
        región interesante del espacio compresión--rendimiento.

  \item Predictores de la vulnerabilidad emocional. Tres
        propiedades del modelo predicen significativamente la caída de F1
        bajo compresión: la profundidad de cristalización por probing
        ($\rho \approx 0{,}48$), el aislamiento representacional
        ($\rho = 0{,}45$) y, como predictor más fuerte, la norma del
        vector de selectividad neuronal ($\rho = 0{,}64$, $p = 0{,}001$,
        \S\ref{sec:neuron-analysis}). Las emociones cuyo modelo escribe con
        mayor intensidad neuronal son las más vulnerables a perturbaciones
        generales de los pesos, aunque la SVD no las dañe selectivamente.

  \item Disociación geometría SVD--función. La fracción de señal
        emocional que cae en el subespacio eliminado por SVD es
        prácticamente idéntica para las 23 emociones (rango 0{,}53--0{,}56),
        y su correlación con la caída de F1 es nula
        ($\rho = -0{,}004$, $p = 0{,}99$). La SVD opera en un espacio
        matemático que no ``ve'' la función computacional del modelo: elimina
        señal uniformemente. Este resultado, aparentemente negativo, es
        teóricamente importante: la geometría que la norma de Frobenius
        optimiza es ortogonal a la geometría de la especialización
        neuronal.
\end{itemize}


\subsubsection{Compresión informada por interpretabilidad}

La conexión entre los dos ejes del trabajo se materializa en una estrategia
de compresión data-driven que constituye, en términos prácticos, la
contribución más directamente aplicable de la memoria:

\begin{itemize}
  \item Dominancia de la frontera de Pareto. Un algoritmo greedy
        alimentado con datos empíricos de sensibilidad por componente y
        profundidad produce 8 de los 9 puntos Pareto-óptimos sobre 21
        configuraciones evaluadas (6 uniformes, 4 adaptativas, 3 heurísticas
        informadas y 8 greedy). Al ratio de parámetros del 80\%, la
        estrategia greedy retiene el 87\% del F1 macro, frente al 43\% de
        la compresión uniforme y el 25\% de la adaptativa energética
        (\S\ref{sec:greedy-algorithm}).

  \item Validación cruzada por convergencia algorítmica. La
        secuencia de decisiones del algoritmo greedy, que opera
        exclusivamente con datos numéricos de sensibilidad, reproduce
        independientemente los hallazgos cualitativos de interpretabilidad:
        comprime primero Q y K (casi gratuitas), después FFN Output en
        capas tempranas, y nunca toca la FFN Intermediate de las capas
        tardías (\S\ref{sec:greedy-algorithm}). Esta convergencia, entre dos
        líneas metodológicamente independientes, refuerza la validez de
        ambas: si interpretabilidad y compresión llegan a las mismas
        conclusiones por caminos disjuntos, la probabilidad de que ambas
        estén sistemáticamente sesgadas en la misma dirección disminuye.

  \item Resultado negativo informativo: la heurística no basta.
        Las tres estrategias informadas basadas en reglas cualitativas
        (\textit{informed-light}, \textit{informed-moderate},
        \textit{informed-aggressive}) convergen exactamente sobre la
        frontera de las baselines ciegas: \texttt{informed\_aggressive}
        reproduce \texttt{uniform\_r256}, \texttt{informed\_moderate}
        reproduce \texttt{uniform\_r512}, e \texttt{informed\_light}
        requiere un 28{,}5\% más parámetros que el original
        (\S\ref{sec:informed-heuristic}). Ninguna domina la frontera de
        Pareto. Este resultado delimita con precisión el valor aplicado
        de la interpretabilidad mecánica: el conocimiento cualitativo
        identifica qué investigar, pero solo la medición empírica de
        sensibilidad permite determinar cuánto comprimir. La diferencia
        entre saber que la FFN es más importante que la atención y saber
        que Q a rango~256 retiene 100{,}4\% del F1 mientras que FFN
        Intermediate a rango~64 colapsa al 0\% es la diferencia entre
        intuición y diseño.

  \item Recuperación por fine-tuning y observación de regularización
        implícita. Un modelo comprimido al 86{,}4\% de los parámetros, tras
        tres épocas de \emph{fine-tuning}, alcanza un F1 macro superior al baseline
        sin comprimir. Las ganancias se concentran sistemáticamente en las
        emociones infrarrepresentadas, con \emph{embarrassment} pasando de
        F1 = 0{,}267 a F1 = 0{,}509 ($+90{,}6\%$ relativo,
        \S\ref{sec:finetuning-recovery}). Este patrón es consistente con la
        hipótesis de que la SVD actúa como regularización estructural,
        al eliminar direcciones de baja energía que pueden codificar
        sobreajuste a patrones espurios. Su verificación
        causal requeriría experimentos de control que se proponen en
        \S\ref{sec:trabajo-futuro}.
\end{itemize}


% -------------------------------------------------------------
\subsection{Limitaciones}
\label{sec:limitaciones}
% -------------------------------------------------------------

Los resultados de este trabajo deben interpretarse dentro de un alcance
experimental delimitado, cuyas restricciones se enumeran a continuación con
honestidad metodológica.

\subsubsection{Generalización: un modelo, una tarea}

Todos los experimentos se realizan sobre un único modelo (BERT-base, 110\,M
parámetros) y una única tarea (clasificación emocional multi-etiqueta sobre
GoEmotions). No es posible afirmar que los patrones observados (la
jerarquía de sensibilidad Q/K $\ll$ V $\ll$ FFN, la concentración funcional
en capas tardías, la correlación cristalización-vulnerabilidad, la
dominancia de la FFN sobre la atención en \emph{activation patching}) se
mantengan en otros modelos o en otras tareas. La arquitectura encoder-only de
BERT, con su cabeza de clasificación sobre el token \texttt{[CLS]}, podría
favorecer la concentración de información en capas finales de un modo que no
se reproduzca en arquitecturas decoder-only o encoder-decoder. La extensión
sistemática a BERT-large, RoBERTa, GPT-2 y LLaMA constituye la línea de
trabajo futuro de mayor prioridad (\S\ref{sec:trabajo-futuro}).

\subsubsection{Sesgos del dataset}

GoEmotions presenta desequilibrios notables entre clases. Cinco emociones
fueron excluidas del análisis (neutral, \emph{grief},
\emph{nervousness}, \emph{pride}, \emph{relief}) por dos razones distintas
que conviene no confundir: neutral se excluyó por no constituir una
emoción propiamente dicha y por dominar el dataset (más de 14\,000 ejemplos),
distorsionando las métricas macro-promediadas; las cuatro restantes se
excluyeron por producir F1 = 0 en el baseline tras \emph{fine-tuning}, debido a su
bajísima frecuencia (menos de 100 ejemplos en entrenamiento). Las 23
emociones restantes varían en frecuencia desde 303 ejemplos
(\emph{embarrassment}) hasta 4\,130 (\emph{admiration}), un factor 13$\times$.
Aunque la frecuencia no predice la profundidad de cristalización, sí podría
afectar la robustez estadística de las métricas de especialización neuronal y
la sensibilidad emocional a la compresión, especialmente en emociones con
menos de 500 ejemplos.

\subsubsection{Compresión post-hoc, no compresión durante entrenamiento}

La compresión se ha limitado a SVD aplicada después del \emph{fine-tuning}, sin
re-entrenamiento durante la compresión y sin combinación con otras técnicas.
Otras aproximaciones (poda estructurada, cuantización, destilación de
conocimiento o entrenamiento consciente de la compresión) podrían interactuar de forma diferente con
la estructura funcional del modelo y posiblemente alcanzar mejores
trade-offs. En particular, el pruning estructurado de cabezas de atención
podría beneficiarse directamente de la taxonomía presentada en
\S\ref{sec:head-ablation}, donde las 38 cabezas dispensable y las 21
con efecto interferente son candidatas naturales a eliminación sin
re-entrenamiento.

\subsubsection{Limitaciones de las técnicas de interpretabilidad}

\paragraph{Probes lineales.} Los \textit{probing classifiers} utilizados son
regresiones logísticas con regularización por defecto ($C = 1{,}0$),
deliberadamente simples siguiendo la metodología de Hewitt y
Liang~\cite{hewitt2019designing}. Esta elección tiene ventajas (mide la
geometría de la representación, no la capacidad del clasificador) pero
también limitaciones: información emocional codificada de forma no lineal en
capas tempranas no sería detectada. Es posible que las capas Emb--L2
contengan más información emocional de la que el probing lineal documenta.

\paragraph{\emph{Logit lens} y calibración.} La aplicación del clasificador final a
representaciones de capas intermedias produce predicciones cuyas magnitudes
absolutas no son directamente comparables al baseline, ya que la cabeza de
clasificación se entrenó exclusivamente sobre representaciones de L11 tras
LayerNorm. Reportamos los resultados con esta cautela: las cifras absolutas
funcionan como diagnóstico de tendencias, no como medida calibrada. La
forma cualitativa del patrón en U es robusta y se corrobora con probing y
patching, pero las magnitudes específicas (gold sigmoid de 0{,}02 en L7,
0{,}44 en L11) deben interpretarse en términos relativos. Una versión más
rigurosa con \emph{tuned lens} \cite{belrose2023eliciting} se discute en
\S\ref{sec:trabajo-futuro}.

\paragraph{\emph{Activation patching} y efecto cadena.} El \emph{activation patching}
opera restaurando pesos originales en un modelo previamente comprimido, lo
que introduce un acoplamiento que conviene reconocer: la contribución medida
de una capa depende del estado de las capas no restauradas. Esto explica por
qué las capas tempranas muestran recuperación nula a pesar de ser
informativas según probing: su señal limpia se destruye al atravesar 10
capas comprimidas posteriores. La corrupción mediante SVD truncada a
rango~64 tampoco es neutral en sentido estadístico (a diferencia del
\emph{causal tracing} clásico de Meng et al.\ \cite{meng2022locating},
que usa ruido gaussiano), sino una distorsión estructurada. Las
conclusiones de patching tienen por tanto naturaleza funcional (qué
componentes son suficientes para reactivar el modelo desde el colapso)
más que estrictamente causal en sentido de
Pearl. La distinción es importante para futuras réplicas.

\subsubsection{Potencia estadística}

Varias correlaciones centrales del trabajo (cristalización vs.\
vulnerabilidad, $\rho \approx 0{,}48$; norma de selectividad vs.\ caída de
F1, $\rho = 0{,}64$) se calculan sobre $n = 23$ puntos (uno por emoción).
Aunque las correlaciones más relevantes son altamente significativas
($p \leq 0{,}001$), la potencia estadística es limitada y no permite
detectar efectos moderados con confianza.

Cuantitativamente, mediante la transformación $z$ de Fisher con
$n=23$, $\alpha=0{,}05$ y potencia $1-\beta = 0{,}80$, la mínima
correlación de Spearman detectable es
$|\rho_{\min}| \approx 0{,}556$. Las correlaciones fuertes
documentadas en la memoria (cristalización vs.\ F1 máximo:
$\rho = -0{,}673$; norma de selectividad neuronal vs.\ caída de F1:
$\rho = +0{,}639$) están \emph{justo dentro} del rango detectable, lo
que es consistente con sus $p$-valores por debajo de $0{,}001$. En
contraste, la correlación frecuencia--cristalización ($\rho = -0{,}265$)
no es detectable a nuestro $n$: para detectarla con $\alpha=0{,}05$ y
potencia $0{,}80$ se necesitarían aproximadamente $n \geq 85$ emociones.
La conclusión metodológica es directa: con $n=23$, los efectos
\emph{grandes} ($|\rho| > 0{,}55$) son robustamente detectables, los
efectos \emph{moderados} ($0{,}3 \leq |\rho| \leq 0{,}55$) requieren
muestras mayores, y los efectos \emph{pequeños} ($|\rho| < 0{,}3$) son
inobservables. Los intervalos de confianza bootstrap al 95\% reportados
en \S\ref{sec:probing} reflejan esta limitación con honestidad: la
amplitud de los intervalos para correlaciones moderadas
(p.~ej.\ frecuencia--cristalización $[-0{,}65, +0{,}19]$) cruza cero,
señalando ausencia de evidencia, no evidencia de ausencia.

\subsubsection{Recuperación por \emph{fine-tuning}: ausencia de grupo de control}

La afirmación de que la compresión actúa como regularización implícita se
basa en la observación de que un modelo comprimido + \emph{fine-tuned} supera al
baseline sin comprimir. Sin embargo, no se ha entrenado un grupo de control
consistente en el baseline original con tres épocas adicionales de
\emph{fine-tuning}, lo que impide aislar el efecto de la compresión del efecto del
re-entrenamiento adicional. Es posible, aunque improbable dada la
saturación de F1 macro observada en epochs 3-4, que parte de la mejora
provenga simplemente de más entrenamiento. Este experimento de control se
propone en \S\ref{sec:trabajo-futuro} como verificación causal de la
hipótesis de regularización.


% -------------------------------------------------------------
\subsection{Trabajo futuro}
\label{sec:trabajo-futuro}
% -------------------------------------------------------------

Las limitaciones identificadas y los hallazgos inesperados del trabajo
sugieren seis líneas concretas de investigación futura, ordenadas por
prioridad metodológica.

\subsubsection{Generalización a otros modelos y tareas}

La pregunta más inmediata es si la jerarquía de sensibilidad por componente
y la correlación cristalización--vulnerabilidad se reproducen en modelos de
mayor escala (BERT-large, RoBERTa-large) y en arquitecturas decoder-only
(GPT-2, LLaMA). Un estudio comparativo permitiría determinar qué hallazgos
son propiedades universales de la arquitectura Transformer y cuáles son
específicos de BERT-base o de la tarea emocional. De particular interés sería
verificar si la dominancia funcional de la FFN sobre la atención
documentada en \S\ref{sec:activation-patching} se mantiene en modelos con más
capas y dimensiones ocultas mayores. La hipótesis de Geva et
al.~\cite{geva2021transformer} sobre las FFN como memorias clave-valor
predice que sí, pero requiere verificación empírica directa con la
metodología de este trabajo.

\subsubsection{Verificación causal del efecto regularizador}

La hipótesis de que la compresión SVD funciona como regularización
implícita merece verificación experimental rigurosa. El diseño propuesto
incluiría tres condiciones de comparación: (i) baseline sin compresión + 3
epochs adicionales de \emph{fine-tuning}, (ii) baseline con compresión greedy +
3 epochs de \emph{fine-tuning} (configuración actual), y (iii) baseline + 3 epochs
con dropout/weight decay incrementado (regularización explícita). Si la
condición (ii) supera a (i) y (iii), la hipótesis de regularización
estructural por SVD obtiene apoyo causal. Si (i) iguala a (ii), el efecto
observado se explica por re-entrenamiento adicional. Este experimento es
factible con los recursos disponibles y constituye la verificación más
inmediata pendiente.

\subsubsection{Compresión a granularidad de cabeza individual}

La estrategia greedy actual opera a granularidad de componente y banda de
profundidad. Sin embargo, los resultados de ablación muestran que dentro de
una misma capa coexisten cabezas críticas y dispensables
(\S\ref{sec:head-ablation}). Una extensión natural sería un algoritmo de
compresión que asigne rangos SVD individuales a cada cabeza de atención en
función de su importancia medida, lo que podría mejorar la frontera de
Pareto especialmente en las capas intermedias (L5--L7), donde la
composición de cabezas críticas y dispensables es más heterogénea. La
combinación con la categoría \emph{Dispensable} de la taxonomía
(38 cabezas) eliminada directamente sin SVD constituiría un híbrido
poda-SVD que podría alcanzar ratios de compresión sustancialmente mayores.

\subsubsection{\emph{Tuned lens} para corrección del sesgo de calibración}

La aplicación del \emph{logit lens} que documentamos sufre de un sesgo conocido:
las representaciones de capas intermedias no están calibradas con la cabeza
de clasificación entrenada exclusivamente sobre L11. Belrose et
al.~\cite{belrose2023eliciting} proponen \textit{tuned lens}: entrenar una
transformación lineal $T_l: \mathbb{R}^d \to \mathbb{R}^d$ por capa que
minimice la divergencia KL entre la predicción del clasificador aplicado a
$T_l(h_l)$ y la predicción real del modelo. Aplicar \emph{tuned lens} a este caso
permitiría calibrar el patrón en U y verificar si el valle observado
(L4--L8) es robusto a la corrección de calibración o si se atenúa, lo que
afinaría la interpretación geométrica de las capas medias.

\subsubsection{Combinación de técnicas de compresión}

Los resultados sugieren una estrategia de compresión compuesta que combine
SVD con poda estructurada y cuantización. La taxonomía de cabezas del
\S\ref{sec:head-ablation} identifica 38 cabezas dispensable y 21 con
efecto interferente, candidatas a eliminación directa sin re-entrenamiento.
Un pipeline en cascada \textit{poda-de-cabezas-interferentes
$\rightarrow$ poda-de-cabezas-dispensable $\rightarrow$ SVD greedy
$\rightarrow$ cuantización post-SVD $\rightarrow$ \emph{fine-tuning}} permitiría
componer las reducciones de los cuatro métodos. La pregunta científica
abierta es si las reducciones se acumulan multiplicativamente (cada técnica
explora dimensiones de redundancia ortogonales) o sub-multiplicativamente
(las técnicas compiten por la misma redundancia).

\subsubsection{Interpretabilidad dinámica durante el entrenamiento}

Todos los análisis de este trabajo operan sobre un modelo ya entrenado. Una
extensión ambiciosa consistiría en monitorizar la cristalización y la
especialización neuronal durante el proceso de \emph{fine-tuning}, no solo al
final. Esto permitiría observar cómo emergen los patrones documentados
(cuándo cristaliza cada emoción, cuándo se establece la dominancia de la
FFN, cuándo aparecen las cabezas interferentes) y eventualmente informar
decisiones de \textit{early stopping} o adaptación dinámica del learning
rate por capa. La integración con frameworks de interpretabilidad como
TransformerLens permitiría implementar este análisis sin coste computacional
prohibitivo.

\subsubsection{Validación en entorno de producción}

Los tiempos de inferencia y el consumo de memoria de los modelos
comprimidos no se han medido en este trabajo, que se ha centrado en la
retención de F1 macro. Una evaluación de latencia y throughput en hardware
real (CPU, GPU, dispositivos edge) complementaría el análisis y determinaría
si las ventajas teóricas de la compresión informada se traducen en
beneficios prácticos de despliegue. Esta extensión es relativamente directa y
debería preceder a cualquier propuesta de adopción industrial de la
metodología.


% -------------------------------------------------------------
\subsection{Predicciones falsables del framework}
\label{sec:predicciones-falsables}
% -------------------------------------------------------------

Un trabajo empírico trasciende el caso estudiado cuando genera
predicciones contrastables sobre escenarios fuera del alcance
experimental original. Las cinco predicciones que siguen se derivan
del framework descrito en esta memoria; cualquier réplica futura sobre
otros modelos o tareas puede confirmarlas o refutarlas, delimitando
empíricamente qué hallazgos son específicos de BERT-base sobre
GoEmotions y cuáles son propiedades estructurales más generales.

La primera predicción concierne a la \emph{asimetría espectral entre
Q/K y FFN en BERT-large}. Si esta asimetría
(\S\ref{sec:analisis-espectral}) es una propiedad arquitectónica del
Transformer encoder y no un artefacto de la pequeña escala de
BERT-base, BERT-large debería exhibir un cociente similar entre el
rango efectivo de Q/K y FFN. Cuantitativamente, esperamos
$k_{95}(\text{Q})/k_{95}(\text{FFN}) \in [0{,}55, 0{,}75]$
(BERT-base presenta $0{,}64$). Un cociente cercano a $1$ en BERT-large
indicaría que la asimetría es específica del régimen de baja capacidad
y obligaría a revisar la generalización.

La segunda predicción distingue \emph{encoder-only de decoder-only}.
El hallazgo de \S\ref{sec:activation-patching} (restaurar L11 sola
recupera el 100\%) descansa sobre la geometría de clasificación a
través del token \texttt{[CLS]}. En modelos decoder-only sin cabeza de
clasificación dedicada (GPT-2, LLaMA), el patching debería distribuirse
más uniformemente entre las capas tardías; replicar la misma
concentración extrema en una sola capa obligaría a revisar la
explicación geométrica que articula el logit lens.

La tercera concierne a la \emph{independencia de la cristalización
respecto a la arquitectura}. La distribución 9/8/6 emociones por banda
de cristalización (\S\ref{sec:probing}) y la correlación
cristalización--F1 máximo ($\rho = -0{,}67$) deberían replicarse
cualitativamente en RoBERTa-base sobre GoEmotions, dada la similitud
arquitectónica. Si el orden de cristalización de emociones
\emph{individuales} cambia marcadamente entre modelos, la profundidad
de procesamiento dependería del corpus de pre-entrenamiento
(Wikipedia + BookCorpus para BERT vs.\ corpus extendido para RoBERTa)
y no únicamente de la complejidad lingüística de la emoción.

La cuarta predice que el \emph{efecto regularizador de la SVD se
reproduce en otras tareas con clases desequilibradas}. La observación
de \S\ref{sec:finetuning-recovery} (compresión + \emph{fine-tuning}
beneficia desproporcionadamente a las emociones infrarrepresentadas)
debería aparecer en NER con etiquetas raras o en clasificación
multi-etiqueta sobre datasets desbalanceados como Reuters-21578 o
AAPD. La ausencia del efecto refutaría la hipótesis de que la
regularización implícita procede de la eliminación de direcciones de
baja energía, o exigiría condiciones adicionales no identificadas en
este trabajo.

La quinta es la más útil prácticamente: \emph{la heurística
cualitativa converge sobre la frontera ciega también en otros
dominios}. El resultado negativo de \S\ref{sec:informed-heuristic}
(las tres heurísticas informadas igualan a las baselines uniformes)
predice un fenómeno general: cualquier traducción cualitativa de
hallazgos interpretativos a reglas de compresión convergerá sobre lo
que la compresión ciega alcanza, salvo que se incorporen datos
cuantitativos de sensibilidad. Convertir esta observación en una guía
para futuros trabajos de ``compresión informada por interpretabilidad''
ahorraría meses de trabajo en proyectos similares.

Las cinco predicciones son ortogonales: cada una falla bajo condiciones
distintas. La confirmación conjunta supondría un test estricto del
framework; la refutación individual delimitaría con precisión qué parte
del análisis es generalizable y cuál pertenece al caso particular.


% -------------------------------------------------------------
\subsection{Reflexión final}
\label{sec:reflexion-final}
% -------------------------------------------------------------

Este trabajo partió de una pregunta aparentemente técnica (cuánto se puede
comprimir un modelo de lenguaje sin degradar su rendimiento) y derivó en
una investigación sobre la estructura funcional interna del modelo. El
camino recorrido sugiere una observación metodológica que excede el caso
estudiado: la compresión y la interpretabilidad no son disciplinas
disjuntas, sino dos perspectivas sobre el mismo objeto.

La compresión es interpretabilidad implícita: cuando comprimimos un
componente y observamos qué se rompe, estamos midiendo su contribución
funcional con la herramienta más directa posible: su ausencia. La
interpretabilidad es compresión racional: cuando entendemos qué hace cada
componente, sabemos cuáles podemos eliminar sin perder función. El
resultado negativo de la heurística (\S\ref{sec:informed-heuristic}) y el
éxito del algoritmo greedy (\S\ref{sec:greedy-algorithm}) precisan esta
afirmación: la conexión entre ambas disciplinas no es cualitativa sino
cuantitativa. Los hallazgos de probing, patching y ablación describen la
arquitectura funcional, pero solo los datos numéricos de sensibilidad
permiten traducir esa descripción en una estrategia de compresión efectiva.

El hallazgo central de la memoria, en esta lectura, no es ninguno de los
resultados individuales (la asimetría 72$\times$ entre componentes, la FFN
de la capa 11 como cuello de botella, la dominancia de la frontera de
Pareto, el efecto regularizador del \emph{fine-tuning}), sino su convergencia.
Cinco técnicas de interpretabilidad, dos familias de compresión y un
algoritmo de optimización ciego apuntan al mismo eje funcional: las capas
8--11, y especialmente sus FFN, son donde reside la capacidad emocional del
modelo. Que cinco metodologías independientes lleguen a la misma conclusión
hace la conclusión más robusta que cualquiera de ellas individualmente.

Una segunda observación, más especulativa, concierne al carácter de los
modelos de lenguaje pre-entrenados como objetos de estudio. BERT no fue
diseñado para clasificar emociones; emergió de la competencia general
adquirida durante el pre-entrenamiento masked-language-modeling sobre Wikipedia
y BookCorpus. El \emph{fine-tuning} sobre GoEmotions no creó la arquitectura
funcional documentada aquí: la organizó. Las capas tardías como cuello de
botella geométrico, la dominancia de la FFN, la cristalización progresiva,
el patrón en U del \emph{logit lens}: ninguno de estos fenómenos fue explícitamente
optimizado durante el entrenamiento. Emergieron. Esto sugiere que las
estructuras funcionales de los Transformers fine-tuneados son, en sentido
fuerte, \emph{descubrimientos empíricos}, no propiedades de diseño. La
interpretabilidad mecánica, vista así, es la disciplina que documenta lo
que la optimización por gradiente decidió, no lo que sus diseñadores
prescribieron.

Si este trabajo cumple su función, debería contribuir modestamente a esa
disciplina, ofreciendo un mapa anatómico parcial pero coherente de cómo un
Transformer organiza el procesamiento emocional, y demostrando que ese
mapa puede informar la compresión más allá de lo que las heurísticas
permiten. Que el resultado se obtenga sobre un modelo y una tarea concretos
limita su alcance; que la metodología sea replicable amplifica su utilidad.
Las extensiones a otros modelos y tareas, propuestas en
\S\ref{sec:trabajo-futuro}, son la prueba pendiente.

% --- Bibliografía ---
\newpage
\printbibliography[title={Referencias}]

% --- Anexo de gestión del proyecto (GEP) ---
\appendix
% =============================================================
% ANEXO — Gestión del proyecto (GEP)
% Adaptado a partir de la Entrega 4 GEP (Fita Inicial) con las
% reconciliaciones necesarias respecto a la ejecución real:
%   - 28 → 23 emociones (filtrado documentado en cap 3 §3.1.1)
%   - 364 → 276 probes activos (cap 5 §5.1)
%   - G6 incorpora logit lens como 5ª técnica (cap 5 §5.2)
%   - T26 fine-tuning post-SVD: ejecutado, no opcional (cap 6 §6.4)
%   - Sección A.8 nueva: ejecución real vs planificación
% =============================================================

\section{Gestión del proyecto (GEP)}
\label{sec:gep}

Este anexo recoge la documentación de gestión del proyecto exigida por
el módulo GEP de la FIB. Se redactó originalmente como Entrega 4
(\textit{Fita Inicial}) durante el primer mes del TFG y se actualiza
aquí con las reconciliaciones necesarias respecto a la ejecución real
del proyecto. La sección final (\S\ref{sec:gep-ejecucion-real}) compara
la planificación inicial con lo efectivamente ejecutado.

% -------------------------------------------------------------
\subsection{Contextualización}
\label{sec:gep-contexto}
% -------------------------------------------------------------

La arquitectura Transformer~\cite{vaswani2017attention} ha redefinido el
procesamiento del lenguaje natural. BERT~\cite{devlin2019bert}, con
110~millones de parámetros, fue uno de los primeros modelos en demostrar
que una misma red pre-entrenada podía alcanzar resultados de referencia
en tareas tan distintas como clasificación de texto, extracción de
información o respuesta a preguntas. El problema es que 110~millones de
parámetros es hoy la talla \emph{pequeña}: modelos más recientes escalan
a cientos de miles de millones, y desplegarlos en entornos con recursos
limitados ---dispositivos móviles, sistemas embebidos, infraestructuras
con restricciones de latencia o energía--- sigue siendo un reto
abierto~\cite{strubell2019energy}.

Este TFG se centra en una técnica concreta de compresión ---la
descomposición en valores singulares (SVD)--- aplicada de forma
selectiva a las capas lineales de un modelo Transformer, y la combina
con cinco técnicas de interpretabilidad mecánica para entender
\emph{por qué} ciertos componentes toleran la compresión y otros no.

\subsubsection{Marco del proyecto en la FIB}

El proyecto se enmarca en el Grado en Ingeniería Informática de la
Facultat d'Informàtica de Barcelona (FIB-UPC), especialidad de
\textbf{Computación}, e integra contenidos de varias asignaturas:

\begin{itemize}[nosep]
    \item \textbf{Matemáticas (M1, M2)}: álgebra lineal y cálculo
    matricial, base teórica de la SVD.
    \item \textbf{Computación Numérica (CN)}: estabilidad numérica y
    aproximación matricial.
    \item \textbf{Inteligencia Artificial (IA)}: fundamentos de
    aprendizaje automático y evaluación de modelos.
    \item \textbf{Algorítmica (A) y Estructuras de Datos (EDA)}: diseño
    y análisis de complejidad, base del algoritmo greedy del
    Capítulo~\ref{sec:sintesis}.
    \item \textbf{Paralelismo (PAR)}: ejecución eficiente en GPU.
    \item \textbf{Projectes de Programació (PROP)}: metodologías de
    desarrollo de software y gestión incremental.
    \item \textbf{Intercambio en DTU (Dinamarca, Q1 2024--25)}:
    asignaturas de nivel máster en \emph{Machine Learning},
    \emph{Deep Learning} y \emph{Graph Neural Networks}, base técnica
    directa para el trabajo con Transformers e interpretabilidad.
\end{itemize}

Adicionalmente, durante el Q2 2025--26 se cursa simultáneamente la
asignatura \textbf{Compresión de Datos e Imagen (CDI)}, complementaria
al tema del TFG.

\subsubsection{Formulación del problema}

\begin{quote}
    \emph{¿Cómo afecta la compresión SVD selectiva ---aplicada por tipo
    de componente y por profundidad de capa--- al rendimiento y a las
    representaciones internas de un BERT fine-tuneado para
    clasificación multi-etiqueta de emociones? ¿Y qué revelan las
    técnicas de interpretabilidad mecánica sobre la organización
    interna de la información emocional?}
\end{quote}

\subsubsection{Actores implicados}

\begin{itemize}[nosep]
    \item \textbf{Investigadores e ingenieros de NLP}: beneficiarios
    primarios de los mapas de sensibilidad y del algoritmo greedy.
    \item \textbf{Comunidad académica y de código abierto}: el código
    modular basado en Hugging Face Transformers~\cite{wolf2020transformers}
    permite reutilizar la metodología para otros modelos y tareas.
    \item \textbf{Empresas que despliegan modelos de lenguaje}:
    organizaciones que integran modelos en sus productos (moderación
    de contenido, análisis de opinión, clasificación de tickets) y
    buscan reducir costes de infraestructura.
\end{itemize}


% -------------------------------------------------------------
\subsection{Justificación del enfoque}
\label{sec:gep-justificacion}
% -------------------------------------------------------------

La compresión de modelos Transformer cuenta con cinco familias
principales de técnicas: poda, cuantización, destilación, factorización
de bajo rango (SVD) y adaptación de bajo rango (LoRA). El
Capítulo~\ref{sec:estado-arte} revisa cada una con detalle. De entre
ellas, la SVD encaja con los objetivos del proyecto por cuatro razones:
(1)~no requiere re-entrenamiento, lo que permite evaluar decenas de
configuraciones a coste bajo; (2)~ofrece granularidad máxima por capa y
componente; (3)~el teorema de Eckart-Young~\cite{eckart1936approximation}
garantiza optimalidad respecto a la norma de Frobenius; y (4)~el
espectro de valores singulares aporta interpretabilidad directa.

\subsubsection{Motivación desde la sostenibilidad}

Strubell et al.~\cite{strubell2019energy} estimaron que entrenar ciertos
modelos de NLP puede emitir cantidades de CO$_2$ comparables a las de
varios coches a lo largo de su vida útil; Schwartz et
al.~\cite{schwartz2020green} acuñaron \emph{Green AI} para abogar por
priorizar la eficiencia. La compresión SVD post-entrenamiento es
afín a esta filosofía: reduce el coste de inferencia ---donde se
concentra el gasto energético a largo plazo--- sin requerir
entrenamiento adicional.


% -------------------------------------------------------------
\subsection{Alcance: objetivos y requisitos}
\label{sec:gep-alcance}
% -------------------------------------------------------------

\subsubsection{Objetivo general}

\begin{quote}
    Caracterizar cómo la compresión SVD selectiva afecta a las
    representaciones internas y al rendimiento de un BERT fine-tuneado
    para clasificación de emociones, combinar análisis de sensibilidad
    con interpretabilidad mecánica, y derivar una estrategia de
    compresión optimizada que maximice la reducción de parámetros con
    mínima pérdida de rendimiento.
\end{quote}

\subsubsection{Sub-objetivos}

\begin{enumerate}[nosep]
    \item \textbf{O1 --- Baseline competitivo}: Demostrar que
    BERT-base-uncased alcanza clasificación multi-etiqueta funcional
    sobre las 23~emociones de GoEmotions tras filtrar las 5~que
    producen $F1=0$ o no son emociones (\emph{neutral}, \emph{grief},
    \emph{nervousness}, \emph{pride}, \emph{relief}; ver
    Cap.~\ref{sec:metodologia-tfg}).
    \item \textbf{O2 --- Módulo SVD reutilizable}: Implementar
    \texttt{SVDLinear} con soporte para compresión uniforme y selectiva
    por componente, capa o ambos.
    \item \textbf{O3 --- Trade-off compresión--rendimiento}: Cuantificar
    F1 macro frente a rangos uniformes $\{512, 384, 256, 128, 64\}$.
    \item \textbf{O4 --- Visualización de la degradación}: Producir
    proyecciones t-SNE de embeddings \texttt{[CLS]}.
    \item \textbf{O5 --- Sensibilidad por componente}: Revelar qué
    tipos de capas (Q, K, V, Output, FFN-Inter, FFN-Out) son más
    vulnerables.
    \item \textbf{O6 --- Sensibilidad por profundidad}: Identificar si
    capas tempranas, medias o tardías son más críticas.
    \item \textbf{O7 --- Compresión adaptativa por energía espectral}:
    Asignar rangos no uniformes según el perfil de energía y evaluar
    en la frontera de Pareto.
    \item \textbf{O8 --- Localizar la cristalización emocional}:
    Mediante probing lineal capa a capa, determinar a qué profundidad
    cada emoción se vuelve linealmente separable. Se entrenan
    23~emociones~$\times$~12~capas~=~276~probes activos (los probes
    sobre la capa de embedding producen $F1=0$ uniformemente y se
    descartan; ver Cap.~\ref{sec:resultados-interpretabilidad}).
    \item \textbf{O9 --- Evidencia funcional de localización}: Usar
    activation patching y logit lens para distinguir información
    \emph{presente}, información \emph{leíble} por el clasificador y
    componentes \emph{suficientes} para la recuperación.
    \item \textbf{O10 --- Mapear especialización a nivel de cabeza y
    neurona}: Ablación sistemática de las 144~cabezas y análisis de
    selectividad de las 36\,864 neuronas FFN.
    \item \textbf{O11 --- Compresión informada}: Integrar O5--O10 en
    una asignación de rangos por capa y componente que supere a las
    estrategias ciegas en la frontera de Pareto.
    \item \textbf{O12 --- Recuperación por fine-tuning post-SVD}:
    Evaluar si tres épocas de re-entrenamiento sobre el modelo
    comprimido recuperan o superan al baseline original
    (Cap.~\ref{sec:sintesis}). \emph{Originalmente clasificado como
    objetivo opcional; promovido a objetivo principal tras observar el
    efecto regularizador.}
\end{enumerate}

\subsubsection{Requisitos funcionales}

\begin{itemize}[nosep]
    \item \textbf{RF1}: Compresión SVD aplicable a cualquier subconjunto
    de capas lineales del encoder.
    \item \textbf{RF2}: Evaluación automática con F1 macro y micro.
    \item \textbf{RF3}: Extracción de embeddings \texttt{[CLS]} y
    proyecciones t-SNE conjuntas.
    \item \textbf{RF4}: Cálculo de silhouette score sobre las
    proyecciones.
    \item \textbf{RF5}: Extracción de hidden states de todas las capas
    para probing y logit lens.
    \item \textbf{RF6}: Mecanismo de activation patching por capa y
    componente.
    \item \textbf{RF7}: Ablación individual de las 144 cabezas con
    medición de F1 por emoción.
    \item \textbf{RF8}: Generación de figuras y tablas exportables.
    \item \textbf{RF9}: Rangos diferenciados por componente y capa.
    \item \textbf{RF10}: Fine-tuning parcial del modelo comprimido.
\end{itemize}

\subsubsection{Requisitos no funcionales}

\begin{itemize}[nosep]
    \item \textbf{RNF1 (Reproducibilidad)}: Semilla fija (42),
    versiones de dependencias documentadas.
    \item \textbf{RNF2 (Modularidad)}: Código en módulos (\texttt{src/data},
    \texttt{src/models}, \texttt{src/compression}, \texttt{src/utils})
    separados de los notebooks.
    \item \textbf{RNF3 (Compatibilidad)}: Notebooks ejecutables en local
    y en Google Colab.
    \item \textbf{RNF4 (Eficiencia de memoria)}: Compresión sobre copias
    del modelo, liberando memoria tras cada evaluación.
    \item \textbf{RNF5 (Estabilidad numérica)}: SVD computada en
    \texttt{float32}.
\end{itemize}

\subsubsection{Obstáculos y riesgos}

\begin{table}[H]
    \centering
    \caption{Principales obstáculos y riesgos del proyecto.}
    \label{tab:gep-riesgos}
    \small
    \begin{tabular}{@{}p{3.8cm} c c p{5.5cm}@{}}
        \toprule
        \textbf{Riesgo} & \textbf{Prob.} & \textbf{Imp.} & \textbf{Mitigación} \\
        \midrule
        R1 --- Recursos computacionales insuficientes (GPU)
        & Media & Alto
        & Google Colab Pro; subsampleo para t-SNE; evaluaciones en CPU si es necesario. \\[0.2cm]
        R2 --- Inestabilidad numérica con rangos muy bajos
        & Baja & Medio
        & SVD en \texttt{float32}; validación dimensional pre-evaluación. \\[0.2cm]
        R3 --- Tiempo de ejecución excesivo
        & Alta & Medio
        & Ejecución secuencial con liberación de memoria; priorización. \\[0.2cm]
        R4 --- Métricas no capturan degradación real
        & Baja & Alto
        & F1 + silhouette + visualización cualitativa. \\[0.2cm]
        R5 --- Bug en selección de capas
        & Media & Medio
        & Tests del módulo de filtrado; alternativas documentadas. \\[0.2cm]
        R6 --- Fine-tuning post-SVD no recupera rendimiento
        & Media & Medio
        & Documentar como hallazgo válido; explorar varias configuraciones. \\[0.2cm]
        R7 --- Desviación del calendario por complejidad imprevista
        & Media & Alto
        & O1--O11 como núcleo; O12 y benchmarks como opcionales. \\
        \bottomrule
    \end{tabular}
    \\[0.15cm]{\small Fuente: elaboración propia.}
\end{table}


% -------------------------------------------------------------
\subsection{Metodología y rigor}
\label{sec:gep-metodologia}
% -------------------------------------------------------------

Dado el carácter experimental e iterativo del proyecto, se adopta una
metodología \emph{incremental iterativa}, tomando elementos de marcos
ágiles como Scrum ---estudiados en PROP--- adaptados a un contexto
unipersonal. Cada fase sigue el ciclo:
\emph{diseño}~$\to$~\emph{implementación}~$\to$~\emph{ejecución}~$\to$~%
\emph{análisis}~$\to$~\emph{validación}, con reuniones quincenales con
el director del TFG como puntos formales de revisión.

Esta elección se justifica por la naturaleza exploratoria del proyecto:
los resultados de compresión uniforme (G3) motivaron el análisis
desagregado por componente (G4); los estudios de lesión (G5) revelaron
la necesidad de técnicas causales como activation patching (G6); y la
brecha entre probing y patching motivó la incorporación del logit lens
(no prevista inicialmente) como cuarta técnica de interpretabilidad.

\subsubsection{Herramientas}

\begin{itemize}[nosep]
    \item \textbf{Control de versiones}: Git + GitHub.
    \item \textbf{Seguimiento}: GitHub Issues como backlog.
    \item \textbf{Ecosistema}: Python 3.10+, PyTorch, Hugging Face
    Transformers~\cite{wolf2020transformers}, scikit-learn,
    matplotlib, seaborn.
    \item \textbf{Dependencias fijadas}: \texttt{requirements.txt}.
    \item \textbf{Documentación}: Overleaf (\LaTeX{}).
    \item \textbf{Cómputo}: PC personal + Google Colab (GPU Tesla T4).
\end{itemize}

\subsubsection{Métodos de validación}

\begin{itemize}[nosep]
    \item \textbf{Cuantitativa}: F1 macro y micro sobre el test
    completo + ratio de parámetros.
    \item \textbf{Representacional}: silhouette score sobre
    proyecciones t-SNE.
    \item \textbf{Visual}: scatter + KDE para inspección cualitativa.
    \item \textbf{Funcional/causal}: activation patching, logit lens
    y ablación de cabezas para distinguir computación, alineación y
    transmisión.
    \item \textbf{Reproducibilidad}: semilla fija (42), código modular,
    notebooks autocontenidos.
\end{itemize}


% -------------------------------------------------------------
\subsection{Planificación temporal}
\label{sec:gep-planificacion}
% -------------------------------------------------------------

\subsubsection{Datos globales}

\begin{table}[H]
\centering
\caption{Parámetros temporales del proyecto.}
\label{tab:gep-globales}
\small
\begin{tabular}{@{}ll@{\quad}ll@{}}
\toprule
\textbf{Parámetro}          & \textbf{Valor}        & \textbf{Parámetro}                  & \textbf{Valor} \\
\midrule
Fecha de inicio              & 2~feb.\ 2026          & Horas diarias de trabajo            & 4~h/día (lectivos) \\
Fecha de fin prevista        & 26~jun.\ 2026         & Total horas núcleo                  & 457~h \\
Duración total               & 22~semanas            & Total horas opcionales              & 40~h \\
Fecha prevista de lectura    & Julio~2026            & \textbf{Total general}              & \textbf{497~h ($\approx$20 ECTS)} \\
\bottomrule
\end{tabular}
\\[0.15cm]{\small Fuente: elaboración propia.}
\end{table}

Aunque todas las tareas son ejecutadas por el autor (Guido Biosca
Lasa), se asigna a cada tarea el \emph{perfil profesional} que mejor
se ajusta al tipo de trabajo: Jefe de Proyecto (JP), Científico de
Datos (CD) o Desarrollador (Dev). Esta diferenciación refleja las
competencias requeridas y permite una estimación presupuestaria
realista (\S\ref{sec:gep-presupuesto}).

\subsubsection{Descripción de las fases}

\paragraph{G0 --- Gestión del proyecto (130~h, concurrente).}
T01--T03 (50~h, JP): tres entregas GEP. T04 (20~h, JP): seguimiento
con director (10 reuniones quincenales). T05 (60~h, JP): redacción
progresiva de la memoria.

\paragraph{G1 --- Configuración y baseline (39~h).}
T06 (8~h, Dev): repositorio + entorno. T07 (25~h, CD): fine-tuning de
BERT-base-uncased sobre GoEmotions, con filtrado a 23~emociones tras
descartar \emph{neutral} (no es emoción) y las cuatro de baja
frecuencia con $F1=0$. T08 (6~h, CD): evaluación baseline (F1 macro
0,577).

\paragraph{G2 --- Módulo de compresión SVD (30~h).}
T09 (20~h, Dev): \texttt{SVDLinear} y \texttt{apply\_svd\_compression()}
en \texttt{src/compression}. T10 (10~h, Dev): tests de validación.

\paragraph{G3 --- Compresión uniforme y análisis espectral (50~h).}
T11 (15~h, CD): experimentos uniformes en 5~rangos. T12 (10~h, CD):
perfil de energía espectral de las 72 matrices. T13 (10~h, CD):
estrategia adaptativa con 5 umbrales (80\%, 85\%, 90\%, 95\%, 99\%).
T14 (15~h, CD): frontera de Pareto y vulnerabilidad por emoción.

\paragraph{G4 --- Análisis por componente y profundidad (45~h).}
T15 (20~h, CD): 24 configuraciones (8~componentes $\times$ 3~rangos)
con extracción de embeddings. T16 (15~h, CD): 9 configuraciones
(3~bandas de profundidad $\times$ 3~rangos). T17 (10~h, CD): heatmaps,
t-SNE conjuntos, silhouette.

\paragraph{G5 --- Lesion study (23~h).}
T18 (15~h, CD): 12 lesiones por capa + 6 por componente. T19 (8~h, CD):
mapa emocional, clustering jerárquico, escáner por emoción.

\paragraph{G6 --- Interpretabilidad mecánica (65~h).}
T20 (15~h, CD): probing lineal sobre 13 hidden states; en términos
efectivos quedan 23~emociones~$\times$~12~capas~=~276~probes activos
tras descartar el embedding (F1 nulo uniforme). T21 (20~h, CD):
activation patching causal a granularidad capa y componente. T22
(10~h, CD): ablación de las 144 cabezas con taxonomía
2$\times$2 (importancia, selectividad). T23 (10~h, CD): selectividad de
las 36\,864 neuronas FFN. \emph{T20bis (10~h, CD): logit lens ---
incorporación no prevista en la planificación inicial; necesaria para
distinguir información presente, leíble y suficiente
(\S\ref{sec:logit-lens}).}

\paragraph{G7 --- Síntesis informada (55~h).}
T24 (20~h, CD): estrategia heurística informada (light, moderate,
aggressive). T25 (10~h, CD): algoritmo greedy data-driven a partir de
las sensibilidades empíricas + comparativa Pareto. T26 (25~h, CD):
fine-tuning post-SVD del modelo \texttt{greedy\_90} (3~epochs);
\emph{originalmente opcional, promovido a tarea principal de G7 tras
observar el efecto regularizador documentado en
\S\ref{sec:finetuning-recovery}}.

\paragraph{G8 --- Objetivos secundarios opcionales (15~h, no
ejecutados).}
T27 (15~h, CD, opcional): benchmarks de latencia y throughput. No
ejecutado en este TFG; documentado como trabajo futuro
(\S\ref{sec:trabajo-futuro}).

\paragraph{G9 --- Cierre y síntesis (55~h).}
T28 (40~h, JP): integración final de la memoria. T29 (15~h, JP):
preparación de la presentación oral.

\subsubsection{Tabla resumen de tareas}

\begin{longtable}{@{}p{0.55cm} p{4.6cm} c c p{2.3cm} p{1.5cm}@{}}
    \caption{Tabla resumen de tareas.}
    \label{tab:gep-tareas}\\
    \toprule
    \textbf{Cód.} & \textbf{Tarea} & \textbf{Perfil} & \textbf{h} & \textbf{Dependencias} & \textbf{Rec.\ mat.} \\
    \midrule
    \endfirsthead
    \multicolumn{6}{c}{\footnotesize(Tabla~\ref{tab:gep-tareas}, continuación)}\\
    \toprule
    \textbf{Cód.} & \textbf{Tarea} & \textbf{Perfil} & \textbf{h} & \textbf{Dependencias} & \textbf{Rec.\ mat.} \\
    \midrule
    \endhead
    \midrule\endfoot
    \bottomrule\\[-4pt]
    \multicolumn{6}{@{}l}{\footnotesize [opt]=opcional; [conc.]=concurrente. Rec.\ humano en todas: Guido Biosca Lasa. Fuente: elab.\ propia.}
    \endlastfoot
    \rowcolor{ganttG0!40}\multicolumn{6}{@{}l}{\footnotesize\textbf{G0 --- Gestión (130~h)}}\\
    T01 & Contextualización (E1)               & JP  & 20 & ---     & Overleaf \\
    T02 & Planificación (E2)                    & JP  & 15 & T01     & Overleaf \\
    T03 & Presupuesto + sostenibilidad (E3)     & JP  & 15 & T01     & Overleaf \\
    T04 & Reuniones director [conc.]            & JP  & 20 & T01     & Videoconf. \\
    T05 & Redacción memoria [conc.]             & JP  & 60 & T01     & Overleaf \\
    \midrule
    \rowcolor{ganttG1!40}\multicolumn{6}{@{}l}{\footnotesize\textbf{G1 --- Baseline (39~h)}}\\
    T06 & Config.\ entorno + repo               & Dev &  8 & T01     & PC, GitHub \\
    T07 & Fine-tuning BERT (23 emociones)       & CD  & 25 & T06     & GPU \\
    T08 & Evaluación baseline                   & CD  &  6 & T07     & PC \\
    \midrule
    \rowcolor{ganttG2!40}\multicolumn{6}{@{}l}{\footnotesize\textbf{G2 --- Módulo SVD (30~h)}}\\
    T09 & Impl.\ \texttt{SVDLinear}              & Dev & 20 & T08     & PC \\
    T10 & Tests validación                      & Dev & 10 & T09     & PC \\
    \midrule
    \rowcolor{ganttGE!30}\multicolumn{6}{@{}l}{\footnotesize\textbf{G3 --- Uniforme + espectral (50~h)}}\\
    T11 & Comp.\ uniforme (5 rangos)            & CD  & 15 & T10     & GPU \\
    T12 & Análisis espectral 72 matrices        & CD  & 10 & T10     & PC \\
    T13 & Adaptativa (5 umbrales energéticos)   & CD  & 10 & T12     & PC \\
    T14 & Pareto + vulnerabilidad emocional     & CD  & 15 & T11,T13 & GPU \\
    \midrule
    \rowcolor{ganttGE!45}\multicolumn{6}{@{}l}{\footnotesize\textbf{G4 --- Comp.\ + profundidad (45~h)}}\\
    T15 & Componente (24 configs)               & CD  & 20 & T11     & GPU \\
    T16 & Profundidad (9 configs)               & CD  & 15 & T15     & GPU \\
    T17 & Heatmaps + t-SNE + silhouette         & CD  & 10 & T16     & PC \\
    \midrule
    \rowcolor{ganttGE!55}\multicolumn{6}{@{}l}{\footnotesize\textbf{G5 --- Lesion study (23~h)}}\\
    T18 & Lesiones capa (12) + componente (6)   & CD  & 15 & T10     & GPU \\
    T19 & Mapa emocional + clustering           & CD  &  8 & T18     & PC \\
    \midrule
    \rowcolor{ganttGE!70}\multicolumn{6}{@{}l}{\footnotesize\textbf{G6 --- Interpretabilidad (65~h)}}\\
    T20 & Probing lineal (276 probes activos)    & CD  & 15 & T08     & PC \\
    T20bis & Logit lens [no previsto inicialmente] & CD & 10 & T20  & PC \\
    T21 & Activation patching                    & CD  & 20 & T11     & GPU \\
    T22 & Ablación cabezas (144)                 & CD  & 10 & T08     & GPU \\
    T23 & Análisis neuronal FFN (36\,864)        & CD  & 10 & T08     & GPU \\
    \midrule
    \rowcolor{ganttGE!85}\multicolumn{6}{@{}l}{\footnotesize\textbf{G7 --- Síntesis informada (55~h)}}\\
    T24 & Heurística informada (3 variantes)     & CD  & 20 & T17,T19--T23 & GPU \\
    T25 & Greedy data-driven + Pareto comp.      & CD  & 10 & T24     & PC \\
    T26 & Fine-tuning post-SVD [promov.\ a obligatoria] & CD & 25 & T25 & GPU \\
    \midrule
    \rowcolor{ganttGO!50}\multicolumn{6}{@{}l}{\footnotesize\textbf{G8 --- Opcional no ejecutado (15~h)}}\\
    T27 & Benchmarks latencia/throughput [opt]   & CD  & 15 & T25     & PC+GPU \\
    \midrule
    \rowcolor{ganttG9!50}\multicolumn{6}{@{}l}{\footnotesize\textbf{G9 --- Cierre (55~h)}}\\
    T28 & Integración + revisión memoria         & JP  & 40 & T26     & Overleaf \\
    T29 & Preparación presentación oral          & JP  & 15 & T28     & PC \\
    \midrule
    \multicolumn{3}{@{}l}{\textbf{Total núcleo (T01--T26, T28--T29)}} & \textbf{492} & & \\
    \multicolumn{3}{@{}l}{\textbf{Total opcional (T27)}} & \textbf{15} & & \\
    \multicolumn{3}{@{}l}{\textbf{Total general planificado}} & \textbf{507} & & \\
\end{longtable}

\subsubsection{Gestión del riesgo: planes alternativos}

Los planes B se activan únicamente si el riesgo se materializa
(Tabla~\ref{tab:gep-planesb}).

\begin{table}[H]
\centering
\caption{Planes alternativos asociados a los riesgos R1--R7.}
\label{tab:gep-planesb}
\footnotesize
\begin{tabular}{@{}p{2.1cm} c c p{5.3cm} p{1.5cm} p{1.2cm}@{}}
\toprule
\textbf{Riesgo} & \textbf{Prob.} & \textbf{Imp.} & \textbf{Plan B} & \textbf{Rec.\ adic.} & \textbf{$\Delta$h} \\
\midrule
R1 GPU insuf.       & Media & Alto  & Reducir T15 a 4 componentes; subsamplear t-SNE a 1\,500 puntos. & Ninguno & $-$13~h \\
\midrule
R2 Inest.\ num.     & Baja  & Medio & Validación numérica adicional (5~h) en \texttt{float64}. & 5~h & $+$5~h \\
\midrule
R3 Tiempo excesivo  & Alta  & Medio & Reducir T15--T16; paralelizar T11/T18. & Ninguno & $-$15~h \\
\midrule
R4 Métricas insuf.  & Baja  & Alto  & Análisis cualitativo (8~h): inspección de errores. & 8~h & $+$8~h \\
\midrule
R5 Bug capas        & Media & Medio & T10 incluye tests; indexación manual de respaldo (4~h). & 4~h & $+$4~h \\
\midrule
R6 FT no recupera   & Media & Medio & Documentar como resultado válido. T26 originalmente opcional. & Ninguno & 0~h \\
\midrule
R7 Desviación gral. & Media & Alto  & Eliminar T27 ($-$15~h); reducir T22+T23; documentar como trabajo futuro. & Ninguno & $-$15--40~h \\
\bottomrule
\end{tabular}
\\[0.15cm]{\small En el peor escenario (R1+R3+R7), el núcleo se reduce a $\approx$370~h, garantizando O1--O9. Fuente: elab.\ propia.}
\end{table}

\subsubsection{Diagrama de Gantt}

La cadena de dependencias crítica es:
\[
\text{T06} \to \text{T07} \to \text{T08} \to \text{T09} \to \text{T10}
\to \text{T11} \to \text{T15} \to \text{T16} \to \text{T17} \to \text{T24} \to
\text{T25} \to \text{T26} \to \text{T28} \to \text{T29}.
\]
T04+T05 (gestión, redacción) son concurrentes durante todas las
22~semanas. T12 ejecuta en paralelo con T11. G5 y parte de G6 son
ejecutables en paralelo con G3--G4.

\begin{landscape}
\begin{figure}[H]
\centering
\resizebox{\linewidth}{!}{%
\begin{ganttchart}[
    hgrid, vgrid,
    x unit=0.9cm, y unit title=0.62cm, y unit chart=0.44cm,
    title height=1,
    title label font=\scriptsize\bfseries,
    bar label font=\scriptsize,
    group label font=\bfseries\scriptsize,
    milestone label font=\itshape\scriptsize,
    bar height=0.65, bar top shift=0.175,
    group right shift=0, group top shift=0.65, group height=0.35, group peaks height=0.2,
    milestone height=0.55, milestone top shift=0.225,
]{1}{22}
\gantttitle{Febrero}{4}\gantttitle{Marzo}{4}\gantttitle{Abril}{5}\gantttitle{Mayo}{4}\gantttitle{Junio}{5}\\
\gantttitlelist{1,...,22}{1}\\
%--- G0 ---
\ganttgroup[group/.append style={fill=ganttG0,draw=black}]{G0: Gestión}{1}{22}\\
\ganttbar[name=t01,bar/.append style={fill=ganttG0}]{T01--T03 Entregas GEP}{1}{7}\\
\ganttbar[name=t04,bar/.append style={fill=ganttG0!60,draw=gray!60}]{T04+T05 Reuniones+Doc.\ [conc.]}{1}{22}\\
%--- G1 ---
\ganttgroup[group/.append style={fill=ganttG1,draw=black}]{G1: Baseline}{1}{7}\\
\ganttbar[name=t06,bar/.append style={fill=ganttG1}]{T06 Config.\ entorno}{1}{2}\\
\ganttbar[name=t07,bar/.append style={fill=ganttG1}]{T07 Fine-tuning BERT}{2}{6}\\
\ganttbar[name=t08,bar/.append style={fill=ganttG1}]{T08 Eval.\ baseline}{6}{7}\\
%--- G2 ---
\ganttgroup[group/.append style={fill=ganttG2,draw=black}]{G2: Módulo SVD}{7}{10}\\
\ganttbar[name=t09,bar/.append style={fill=ganttG2}]{T09 Impl.\ SVDLinear}{7}{9}\\
\ganttbar[name=t10,bar/.append style={fill=ganttG2}]{T10 Tests validación}{9}{10}\\
%--- G3 ---
\ganttgroup[group/.append style={fill=ganttGE!70,draw=black}]{G3: Uniforme+Espectral}{10}{13}\\
\ganttbar[name=t11,bar/.append style={fill=ganttGE!70}]{T11 Comp.\ uniforme}{10}{12}\\
\ganttbar[name=t12,bar/.append style={fill=ganttGE!50}]{T12 Análisis espectral}{10}{11}\\
\ganttbar[name=t13,bar/.append style={fill=ganttGE!50}]{T13--T14 Adaptativa+Pareto}{12}{13}\\
%--- G4 ---
\ganttgroup[group/.append style={fill=ganttGE,draw=black}]{G4: Comp.+Prof.}{12}{16}\\
\ganttbar[name=t15,bar/.append style={fill=ganttGE}]{T15 Exp.\ componente (24)}{12}{14}\\
\ganttbar[name=t16,bar/.append style={fill=ganttGE}]{T16 Exp.\ profundidad (9)}{14}{15}\\
\ganttbar[name=t17,bar/.append style={fill=ganttGE}]{T17 Heatmaps+t-SNE}{15}{16}\\
%--- G5 ---
\ganttgroup[group/.append style={fill=ganttGE!85,draw=black}]{G5: Lesion study}{10}{13}\\
\ganttbar[name=t18,bar/.append style={fill=ganttGE!85}]{T18--T19 Lesiones+Mapa}{10}{13}\\
%--- G6 ---
\ganttgroup[group/.append style={fill=ganttGE!90,draw=black}]{G6: Interpretabilidad}{8}{17}\\
\ganttbar[name=t20,bar/.append style={fill=ganttGE!90}]{T20 Probing lineal}{8}{10}\\
\ganttbar[name=t20b,bar/.append style={fill=ganttGE!95}]{T20bis Logit lens [nuevo]}{14}{15}\\
\ganttbar[name=t21,bar/.append style={fill=ganttGE!90}]{T21 Act.\ patching}{12}{14}\\
\ganttbar[name=t22,bar/.append style={fill=ganttGE!90}]{T22--T23 Heads+Neurons}{10}{13}\\
%--- G7 ---
\ganttgroup[group/.append style={fill=ganttGE,draw=black}]{G7: Síntesis}{16}{20}\\
\ganttbar[name=t24,bar/.append style={fill=ganttGE}]{T24--T25 Informada+Greedy}{16}{18}\\
\ganttbar[name=t26,bar/.append style={fill=ganttGE}]{T26 Fine-tuning post-SVD}{18}{20}\\
%--- G9 ---
\ganttgroup[group/.append style={fill=ganttG9,draw=black}]{G9: Cierre}{20}{22}\\
\ganttbar[name=t28,bar/.append style={fill=ganttG9}]{T28 Memoria final}{20}{22}\\
\ganttbar[name=t29,bar/.append style={fill=ganttG9}]{T29 Presentación oral}{22}{22}\\
%--- Hito ---
\ganttmilestone[name=m1,milestone/.append style={fill=red!80,draw=red}]{Lectura TFG}{22}\\
%--- Dependencias ---
\ganttlink{t06}{t07}\ganttlink{t07}{t08}\ganttlink{t08}{t09}\ganttlink{t09}{t10}
\ganttlink{t10}{t11}\ganttlink{t11}{t15}\ganttlink{t15}{t16}\ganttlink{t16}{t17}
\ganttlink{t17}{t24}\ganttlink{t24}{t26}\ganttlink{t26}{t28}\ganttlink{t28}{t29}
\end{ganttchart}
}%
\caption{Diagrama de Gantt actualizado. T20bis (logit lens) y T26 (fine-tuning) reflejan los cambios respecto a la planificación inicial. Grises: opcionales no ejecutados. Flechas: cadena crítica.}
\label{fig:gep-gantt}
\end{figure}
\end{landscape}


% -------------------------------------------------------------
\subsection{Gestión económica (presupuesto)}
\label{sec:gep-presupuesto}
% -------------------------------------------------------------

El presupuesto se estructura en cuatro partidas siguiendo la
metodología del módulo: costes de personal por actividad (CPA),
costes genéricos (CG), contingencia e imprevistos.

\subsubsection{Costes por hora por perfil profesional}

\begin{table}[H]
\centering
\caption{Coste por hora estimado por perfil profesional.}
\label{tab:gep-perfiles}
\small
\begin{tabular}{@{}l r r r@{}}
\toprule
\textbf{Perfil} & \textbf{Bruto anual} & \textbf{Con SS ($\times$1,3)} & \textbf{Coste/hora} \\
\midrule
Jefe de proyecto (JP)    & 32.000\,\euro{} & 41.600\,\euro{} & 24,19\,\euro{}/h $\approx$ \textbf{24\,\euro{}/h} \\
Científico de datos (CD) & 28.000\,\euro{} & 36.400\,\euro{} & 21,16\,\euro{}/h $\approx$ \textbf{21\,\euro{}/h} \\
Desarrollador (Dev)      & 25.000\,\euro{} & 32.500\,\euro{} & 18,90\,\euro{}/h $\approx$ \textbf{19\,\euro{}/h} \\
\bottomrule
\end{tabular}
\\[0.15cm]{\small Fuente: elaboración propia a partir de Hays España 2025 y Glassdoor España. 1.720~horas laborales anuales.}
\end{table}

\subsubsection{Costes de personal por actividad (CPA)}

\begin{longtable}{@{}p{0.5cm} p{4.2cm} c c r r@{}}
    \caption{Costes de personal por actividad (CPA).}
    \label{tab:gep-cpa}\\
    \toprule
    \textbf{Cód.} & \textbf{Tarea} & \textbf{Perfil} & \textbf{h} & \textbf{\euro{}/h} & \textbf{Coste} \\
    \midrule
    \endfirsthead
    \multicolumn{6}{c}{\footnotesize(Tabla~\ref{tab:gep-cpa}, continuación)}\\
    \toprule
    \textbf{Cód.} & \textbf{Tarea} & \textbf{Perfil} & \textbf{h} & \textbf{\euro{}/h} & \textbf{Coste} \\
    \midrule
    \endhead
    \midrule\endfoot
    \bottomrule
    \endlastfoot
    \rowcolor{gray!10}\multicolumn{5}{@{}l}{\footnotesize\textbf{G0 --- Gestión (130~h)}} & \textbf{3.120\,\euro{}} \\
    T01 & Contextualización                  & JP  & 20 & 24 & 480\,\euro{} \\
    T02 & Planificación                       & JP  & 15 & 24 & 360\,\euro{} \\
    T03 & Presupuesto + sostenibilidad        & JP  & 15 & 24 & 360\,\euro{} \\
    T04 & Reuniones director                 & JP  & 20 & 24 & 480\,\euro{} \\
    T05 & Redacción memoria                  & JP  & 60 & 24 & 1.440\,\euro{} \\
    \midrule
    \rowcolor{gray!10}\multicolumn{5}{@{}l}{\footnotesize\textbf{G1 --- Baseline (39~h)}} & \textbf{803\,\euro{}} \\
    T06 & Config.\ entorno                    & Dev &  8 & 19 & 152\,\euro{} \\
    T07 & Fine-tuning BERT (23 emociones)     & CD  & 25 & 21 & 525\,\euro{} \\
    T08 & Eval.\ baseline                     & CD  &  6 & 21 & 126\,\euro{} \\
    \midrule
    \rowcolor{gray!10}\multicolumn{5}{@{}l}{\footnotesize\textbf{G2 --- Módulo SVD (30~h)}} & \textbf{570\,\euro{}} \\
    T09 & Impl.\ \texttt{SVDLinear}            & Dev & 20 & 19 & 380\,\euro{} \\
    T10 & Tests validación                    & Dev & 10 & 19 & 190\,\euro{} \\
    \midrule
    \rowcolor{gray!10}\multicolumn{5}{@{}l}{\footnotesize\textbf{G3 --- Uniforme + espectral (50~h)}} & \textbf{1.050\,\euro{}} \\
    T11--T14 & Uniforme, espectro, adapt., Pareto & CD & 50 & 21 & 1.050\,\euro{} \\
    \midrule
    \rowcolor{gray!10}\multicolumn{5}{@{}l}{\footnotesize\textbf{G4 --- Comp.\ + profundidad (45~h)}} & \textbf{945\,\euro{}} \\
    T15--T17 & Componente, profundidad, t-SNE  & CD & 45 & 21 & 945\,\euro{} \\
    \midrule
    \rowcolor{gray!10}\multicolumn{5}{@{}l}{\footnotesize\textbf{G5 --- Lesion study (23~h)}} & \textbf{483\,\euro{}} \\
    T18--T19 & Lesiones + mapa emocional        & CD & 23 & 21 & 483\,\euro{} \\
    \midrule
    \rowcolor{gray!10}\multicolumn{5}{@{}l}{\footnotesize\textbf{G6 --- Interpretabilidad (65~h)}} & \textbf{1.365\,\euro{}} \\
    T20 & Probing lineal (276 probes)          & CD & 15 & 21 & 315\,\euro{} \\
    T20bis & Logit lens [no previsto]          & CD & 10 & 21 & 210\,\euro{} \\
    T21 & Activation patching                  & CD & 20 & 21 & 420\,\euro{} \\
    T22 & Ablación cabezas                     & CD & 10 & 21 & 210\,\euro{} \\
    T23 & Análisis neuronal FFN                & CD & 10 & 21 & 210\,\euro{} \\
    \midrule
    \rowcolor{gray!10}\multicolumn{5}{@{}l}{\footnotesize\textbf{G7 --- Síntesis informada (55~h)}} & \textbf{1.155\,\euro{}} \\
    T24 & Heurística informada                 & CD & 20 & 21 & 420\,\euro{} \\
    T25 & Greedy + Pareto comp.                & CD & 10 & 21 & 210\,\euro{} \\
    T26 & Fine-tuning post-SVD                 & CD & 25 & 21 & 525\,\euro{} \\
    \midrule
    \rowcolor{gray!10}\multicolumn{5}{@{}l}{\footnotesize\textbf{G9 --- Cierre (55~h)}} & \textbf{1.320\,\euro{}} \\
    T28 & Memoria final                        & JP & 40 & 24 & 960\,\euro{} \\
    T29 & Presentación oral                    & JP & 15 & 24 & 360\,\euro{} \\
    \midrule
    \multicolumn{3}{@{}l}{\textbf{Total CPA núcleo (sin opcional T27)}} & \textbf{492} & & \textbf{10.811\,\euro{}} \\
\end{longtable}

\subsubsection{Costes genéricos (CG)}

\begin{table}[H]
\centering
\caption{Costes genéricos del proyecto.}
\label{tab:gep-cg}
\small
\begin{tabular}{@{}l l r r r@{}}
\toprule
\textbf{Concepto} & \textbf{Detalle} & \textbf{Coste total} & \textbf{Vida útil} & \textbf{Imputado} \\
\midrule
Amortización PC portátil & 1.200\,\euro{}, uso 5 meses & 1.200\,\euro{} & 4 años (48~m.) & 125\,\euro{} \\
Google Colab Pro         & 11,99\,\euro{}/mes $\times$ 5 meses & --- & --- & 60\,\euro{} \\
Electricidad             & 0,1\,kW $\times$ 492\,h $\times$ 0,20\,\euro{}/kWh & --- & --- & 10\,\euro{} \\
Software open source     & Python, PyTorch, HF, scikit-learn & --- & --- & 0\,\euro{} \\
Overleaf, GitHub         & Plan gratuito & --- & --- & 0\,\euro{} \\
Internet                 & Uso parcial conexión doméstica & --- & --- & 25\,\euro{} \\
Espacio de trabajo       & Uso parcial domicilio (10\,m$^2$, 5 meses) & --- & --- & 50\,\euro{} \\
\midrule
\multicolumn{4}{@{}l}{\textbf{Total costes genéricos}} & \textbf{270\,\euro{}} \\
\bottomrule
\end{tabular}
\\[0.15cm]{\small Amortización: $1.200 \times 5/48 = 125$\,\euro{}. Trabajo desde domicilio (sin desplazamientos). Fuente: elab.\ propia.}
\end{table}

\subsubsection{Contingencia e imprevistos}

Se aplica un ratio de contingencia del \textbf{15\%} sobre la suma de
CPA y CG, dentro del rango habitual del 10--20\% en proyectos
informáticos:
\[
\text{Contingencia} = (10.811 + 270) \times 0{,}15 = \textbf{1.662\,\euro{}}.
\]

Las partidas de imprevistos se vinculan a los riesgos identificados,
ponderadas por probabilidad:

\begin{table}[H]
\centering
\caption{Partidas de imprevistos.}
\label{tab:gep-imprevistos}
\small
\begin{tabular}{@{}l c r r r@{}}
\toprule
\textbf{Riesgo} & \textbf{Prob.} & \textbf{Coste si ocurre} & \textbf{Detalle} & \textbf{Imputado} \\
\midrule
R1 GPU insuf.       & 50\% & 150\,\euro{} & Upgrade Colab Pro+ (3~m.) & 75\,\euro{} \\
R2 Inest.\ num.     & 20\% & 105\,\euro{} & 5~h extra CD              & 21\,\euro{} \\
R4 Métricas insuf.  & 20\% & 168\,\euro{} & 8~h extra CD              & 34\,\euro{} \\
R5 Bug capas        & 50\% & 76\,\euro{}  & 4~h extra Dev             & 38\,\euro{} \\
\midrule
\multicolumn{4}{@{}l}{\textbf{Total imprevistos}} & \textbf{168\,\euro{}} \\
\bottomrule
\end{tabular}
\\[0.15cm]{\small R3, R6 y R7 se mitigan reduciendo alcance, sin coste adicional. Fuente: elab.\ propia.}
\end{table}

\subsubsection{Presupuesto total}

\begin{table}[H]
\centering
\caption{Resumen del presupuesto total del proyecto.}
\label{tab:gep-total}
\small
\begin{tabular}{@{}l r r@{}}
\toprule
\textbf{Partida} & \textbf{Importe} & \textbf{\%} \\
\midrule
Costes de personal (CPA)    & 10.811\,\euro{} & 84{,}1\% \\
Costes genéricos (CG)       & 270\,\euro{}    & 2{,}1\% \\
Contingencia (15\%)          & 1.662\,\euro{}  & 12{,}9\% \\
Imprevistos                  & 168\,\euro{}    & 1{,}3\% \\
\midrule
\textbf{Total presupuesto}  & \textbf{12.911\,\euro{}} & \textbf{100\%} \\
\bottomrule
\end{tabular}
\\[0.15cm]{\small El coste de personal representa $\sim$84\% del total, habitual en proyectos de investigación aplicada. Fuente: elab.\ propia.}
\end{table}

\subsubsection{Control de gestión}

Para detectar y corregir desviaciones presupuestarias durante la
ejecución se definieron los siguientes mecanismos.

\textbf{Procedimiento.} Al finalizar cada fase del Gantt (G1--G9) se
registran las horas reales dedicadas a cada tarea y se comparan con
las estimadas. Las reuniones quincenales con el director (T04) sirven
como puntos formales de revisión.

\textbf{Indicadores.}

\begin{itemize}[nosep, leftmargin=1.2em]
    \item \textbf{Desviación en horas por tarea}:
    $\Delta h_i = h_i^{\text{real}} - h_i^{\text{est}}$. Si
    $\Delta h_i > 0{,}2 \times h_i^{\text{est}}$, se activa una revisión
    de la planificación restante.
    \item \textbf{Desviación en coste por fase}: si
    $\Delta C_g > 0{,}15 \times C_g^{\text{est}}$, se evalúa el plan B
    correspondiente.
    \item \textbf{Índice de rendimiento del coste (CPI)}:
    $\text{CPI} = \text{EV}/\text{AC}$. Un CPI~$<0{,}85$ indica
    sobrecostes que requieren acción correctiva.
    \item \textbf{\% de contingencia consumida}: si supera el 50\% antes
    de completar el 75\% del proyecto, se priorizan O1--O11 y se
    eliminan tareas opcionales.
\end{itemize}


% -------------------------------------------------------------
\subsection{Sostenibilidad y compromiso social}
\label{sec:gep-sostenibilidad}
% -------------------------------------------------------------

\subsubsection{Autoevaluación}

Como \textbf{puntos fuertes}, dispongo de sensibilidad hacia el impacto
ambiental de la computación: el propio tema del TFG ---compresión de
modelos--- está motivado en parte por la reducción del coste energético
de la inferencia, y he incorporado referencias a Strubell et
al.~\cite{strubell2019energy} y Schwartz et al.~\cite{schwartz2020green}
sobre \emph{Green AI} desde la definición del proyecto. A nivel
económico, he estimado el presupuesto con detalle, diferenciando
perfiles y justificando cada partida.

Como \textbf{puntos débiles}, mi formación en sostenibilidad ha sido
más teórica que práctica: no había realizado anteriormente un análisis
de ciclo de vida energético de un proyecto informático y mi
conocimiento sobre normativas ambientales aplicables a centros de datos
es limitado. La dimensión social ---más allá de la accesibilidad---
requiere reflexión más profunda sobre sesgos en los datos de emociones
y sobre quién se beneficia realmente de estas tecnologías.

\subsubsection{Dimensión económica}

\textbf{PPP --- Coste estimado del proyecto.} El presupuesto total de
12.911\,\euro{} es razonable para un proyecto académico de
$\sim$497~h. El coste está dominado por el recurso humano (84\%), lo
que refleja un proyecto de investigación aplicada sin necesidad de
infraestructura costosa. El uso de herramientas open source y servicios
gratuitos o de bajo coste (Google Colab Pro, GitHub, Overleaf) mantiene
los costes genéricos en mínimos.

\textbf{VU --- ¿Cómo se resuelven actualmente los costes del problema?}
Desplegar modelos Transformer de gran tamaño requiere GPUs costosas
(A100, H100) con consumos energéticos elevados. Empresas y centros de
investigación invierten miles de euros mensuales en infraestructura
cloud para servir estos modelos. Las técnicas de compresión existentes
(destilación, cuantización) requieren re-entrenamiento costoso o
hardware especializado.

\textbf{¿En qué mejorará económicamente la solución?} La compresión SVD
post-entrenamiento que estudia este proyecto no requiere
re-entrenamiento ni hardware especial. El mapa de sensibilidad
producido permite a ingenieros comprimir selectivamente solo los
componentes menos críticos, maximizando la reducción de parámetros con
mínima pérdida de calidad, y por tanto reduciendo costes de inferencia
sin inversión adicional en entrenamiento.

\subsubsection{Dimensión ambiental}

\textbf{PPP --- Impacto ambiental de la realización del proyecto.} El
principal impacto proviene del consumo eléctrico de las GPUs durante
los experimentos. Estimamos un uso de GPU de $\sim$80--100~horas. Con
un consumo típico de una GPU T4 en Colab de $\sim$70\,W, esto supone
$\sim$7\,kWh. Sumando el PC personal (492~h~$\times$~0,1\,kW =
49\,kWh), el consumo total estimado es de $\sim$56\,kWh, equivalente a
$\sim$12\,kg de CO$_2$ (factor de emisión medio en España: 0,22\,kg
CO$_2$/kWh). Es un impacto modesto comparado con el entrenamiento de
modelos grandes.

\textbf{¿Te has planteado minimizar el impacto?} Sí. Se reutiliza un
modelo BERT-base ya pre-entrenado por Google, evitando el coste enorme
del pre-entrenamiento. La compresión SVD se aplica post-hoc sobre el
modelo fine-tuneado, sin requerir entrenamiento adicional para la
mayoría de experimentos. Se liberan copias del modelo de memoria tras
cada evaluación (RNF4) y se utiliza subsampleo para visualizaciones
(1.500 muestras en t-SNE).

\textbf{VU --- Mejora ambiental.} La contribución ambiental principal
del proyecto es indirecta: los mapas de sensibilidad permiten comprimir
modelos de forma informada, reduciendo el tamaño y el consumo
energético durante la inferencia. Dado que la inferencia representa la
mayor parte del gasto energético a largo plazo de un modelo desplegado,
guiar la compresión de forma óptima tiene un impacto ambiental positivo
acumulativo significativo, alineado con la filosofía
\emph{Green AI}~\cite{schwartz2020green}.

\subsubsection{Dimensión social}

\textbf{PPP --- Aportación personal.} Este proyecto me permite
profundizar en tres áreas clave para mi desarrollo profesional:
(1)~comprensión profunda de la arquitectura Transformer y sus
mecanismos internos, conocimiento fundamental en la IA actual;
(2)~habilidades en álgebra lineal aplicada y análisis numérico,
conectando teoría matemática con aplicaciones prácticas; y
(3)~experiencia en diseño y ejecución de estudios experimentales
sistemáticos.

\textbf{VU --- Mejora social.} Actualmente, los modelos de lenguaje
grandes son accesibles principalmente para grandes empresas con
recursos computacionales significativos, lo que genera una brecha
tecnológica. La compresión eficiente democratiza el acceso a estos
modelos, permitiendo su despliegue en dispositivos con recursos
limitados (móviles, sistemas embebidos) y en organizaciones con
presupuestos reducidos (universidades, ONGs, PYMES). Un modelo de
clasificación de emociones comprimido podría usarse en aplicaciones de
salud mental, moderación de contenido o análisis de opinión en
contextos donde antes era inviable.

\textbf{VU --- Necesidad real.} Sí. La tendencia hacia modelos cada
vez más grandes hace que la compresión sea no solo deseable sino
necesaria para su adopción generalizada. Además, la clasificación de
emociones tiene aplicaciones reales en salud mental, educación y
moderación de plataformas. Este proyecto no solo demuestra
\emph{cómo} comprimir, sino \emph{dónde} hacerlo de forma segura, lo
cual es una contribución con utilidad práctica directa.


% -------------------------------------------------------------
\subsection{Ejecución real frente a planificación}
\label{sec:gep-ejecucion-real}
% -------------------------------------------------------------

Esta sección, redactada al cierre del proyecto, contrasta lo
ejecutado con la planificación inicial.

\subsubsection{Ajustes de alcance respecto a la Entrega 4}

Durante la ejecución se introdujeron cuatro modificaciones de alcance
respecto a lo recogido en la Entrega 4 GEP, todas en la dirección de
\emph{ampliar la profundidad del análisis} y todas justificadas por los
hallazgos intermedios:

\begin{enumerate}[nosep]
    \item \textbf{Filtrado a 23 emociones (afecta O1, O3, T07, T08).}
    Tras el fine-tuning sobre las 28 emociones de GoEmotions, cuatro
    emociones (\emph{grief}, \emph{nervousness}, \emph{pride},
    \emph{relief}) producían $F1=0$ y la categoría \emph{neutral}
    distorsionaba las métricas macro. Se decidió filtrar las cinco y
    trabajar con 23 emociones funcionales (Cap.~\ref{sec:metodologia-tfg}).
    El filtrado no añadió horas significativas de cómputo.

    \item \textbf{Probes efectivos: 276 en lugar de 364.} La
    planificación inicial estimaba 28~emociones $\times$ 13~capas. Tras
    el filtrado a 23~emociones y el descarte del embedding (donde el
    F1 lineal es nulo uniformemente), el número efectivo de
    clasificadores entrenados es 276 (Cap.~\ref{sec:resultados-interpretabilidad}).
    El coste computacional se redujo en consecuencia.

    \item \textbf{Logit lens (T20bis, +10\,h, no prevista).} La brecha
    entre el probing (información presente) y el activation patching
    (suficiencia funcional) motivó incorporar el logit lens como cuarta
    técnica de interpretabilidad para distinguir información
    \emph{presente} de información \emph{leíble} por la cabeza de
    clasificación entrenada (\S\ref{sec:logit-lens}). Esto extendió G6
    de 55\,h a 65\,h.

    \item \textbf{Fine-tuning post-SVD (T26) promovido a tarea
    principal.} Inicialmente clasificada como opcional dentro de G8,
    el \emph{fine-tuning} de tres épocas sobre el modelo
    \texttt{greedy\_90} produjo una observación que merecía análisis:
    el modelo comprimido al 86{,}4\% de los parámetros alcanza un F1
    macro comparable al baseline original (0{,}591 vs 0{,}577), con
    una mejora sistemática y notable en las emociones
    infrarrepresentadas (\S\ref{sec:finetuning-recovery}). El alcance
    causal del hallazgo está pendiente de un experimento de control
    (baseline + 3 epochs adicionales) propuesto en
    \S\ref{sec:trabajo-futuro}. Dada la relevancia descriptiva, T26
    se reclasificó como tarea principal de G7; como contrapartida,
    T27 (benchmarks de latencia) se pospuso como trabajo futuro.
\end{enumerate}

\subsubsection{Resumen del balance de horas}

\begin{table}[H]
\centering
\caption{Comparativa entre planificación inicial (Entrega 4) y ejecución
real.}
\label{tab:gep-balance-horas}
\small
\begin{tabular}{@{}l r r r@{}}
\toprule
\textbf{Fase} & \textbf{Plan inicial (h)} & \textbf{Ejecutado (h)} & \textbf{$\Delta$} \\
\midrule
G0 Gestión                    & 130 & 130 & 0 \\
G1 Baseline                   & 39  & 39  & 0 \\
G2 Módulo SVD                 & 30  & 30  & 0 \\
G3 Uniforme + espectral       & 50  & 50  & 0 \\
G4 Componente + profundidad   & 45  & 45  & 0 \\
G5 Lesion study               & 23  & 23  & 0 \\
G6 Interpretabilidad          & 55  & 65  & $+$10 (logit lens) \\
G7 Síntesis informada         & 30  & 55  & $+$25 (T26 promovido) \\
G8 Opcionales                 & 40  & 0   & $-$40 (T27 no ejecutado) \\
G9 Cierre                     & 55  & 55  & 0 \\
\midrule
\textbf{Total}                & \textbf{497} & \textbf{492} & $-$5 \\
\bottomrule
\end{tabular}
\\[0.15cm]{\small La incorporación del logit lens y la promoción del fine-tuning post-SVD se compensaron con la exclusión de los benchmarks de latencia, manteniendo el total dentro del rango planificado.}
\end{table}

\subsubsection{Planes B activados}

De los 7 riesgos identificados, sólo \textbf{R3 (tiempo de ejecución
excesivo)} se materializó parcialmente en G4: el experimento por
componente (T15) se ejecutó secuencialmente con liberación explícita de
memoria entre configuraciones para evitar saturar la GPU. No fue
necesario reducir el alcance de T15.

\subsubsection{Indicadores de control finales}

\begin{itemize}[nosep, leftmargin=1.2em]
    \item \textbf{Desviación total en horas}: $-$5\,h sobre 497\,h
    planificadas ($-$1\%). Dentro de tolerancia.
    \item \textbf{CPI estimado}: $\approx$1{,}01 (proyecto en
    presupuesto).
    \item \textbf{Contingencia consumida}: 0\% (no se necesitaron
    horas extra ponderadas por riesgo). El margen del 15\% queda como
    holgura.
    \item \textbf{Tareas opcionales canceladas}: T27 (benchmarks de
    latencia), redirigida a trabajo futuro.
    \item \textbf{Tareas no previstas incorporadas}: T20bis (logit
    lens), por necesidad metodológica detectada en G6.
\end{itemize}


\end{document}